\documentclass[11pt,twoside,openany]{book}
\usepackage[a4paper,inner=31mm,outer=27mm,top=28mm,bottom=31mm,marginparwidth=17mm]{geometry}
\usepackage{fontspec}
\defaultfontfeatures{Ligatures=TeX,Scale=MatchLowercase}
\setmainfont{TeX Gyre Pagella}
\setsansfont{TeX Gyre Heros}
\setmonofont{TeX Gyre Cursor}
\usepackage[brazilian]{babel}
\usepackage{microtype}
\usepackage{graphicx}
\usepackage{mathtools}
\usepackage{amsthm}
\usepackage{unicode-math}
\setmathfont{Latin Modern Math}
\usepackage{bookmark}
\usepackage{hyperref}
\hypersetup{hidelinks,unicode=true}

\setcounter{secnumdepth}{3}
\setcounter{tocdepth}{2}
\numberwithin{equation}{section}

\newcommand{\src}[1]{%
  \phantomsection\hypertarget{src:#1}{}%
  \marginpar{\raggedleft\footnotesize\sffamily #1}%
}
\newcommand{\sid}[1]{\phantomsection\hypertarget{fga:#1}{}}
\newcommand{\Osh}{\underline{O}}
\newcommand{\Zsh}{\underline{\mathbf Z}}
\newcommand{\Ccat}{\mathcal C}
\newcommand{\Homsh}{\underline{\operatorname{Hom}}}
\newcommand{\Extsh}{\underline{\operatorname{Ext}}}
\newcommand{\hsh}{\underline{h}}
\newcommand{\fquote}[1]{\textquotedbl#1\textquotedbl}
\newcommand{\seminar}[1]{%
  \begin{flushleft}\sffamily SEMINÁRIO BOURBAKI\\#1\end{flushleft}%
}
\newcommand{\seminares}[1]{\seminar{#1}}
\newcommand{\paperhead}[2]{%
  \begin{center}
    {\Large\bfseries #1\par}
    \vspace{0.6\baselineskip}
    #2
  \end{center}
  \vspace{0.8\baselineskip}
}

\theoremstyle{plain}
\newtheorem{proposition}{Proposição}[section]
\newtheorem{theorem}[proposition]{Teorema}
\newtheorem{corollary}[proposition]{Corolário}
\newtheorem{lemma}[proposition]{Lema}

\theoremstyle{definition}
\newtheorem{definition}[proposition]{Definição}

\theoremstyle{remark}
\newtheorem*{remark}{Observação}

\newcommand{\unit}[1]{%
  \chapter{#1}%
  \markboth{#1}{#1}%
}
\newcommand{\fsection}[2]{%
  \section*{#1. #2}%
  \addcontentsline{toc}{section}{#1. #2}%
}
\newcommand{\result}[2]{%
  \par\medskip\noindent\textbf{#1%
  \if\relax\detokenize{#2}\relax\else\space#2\fi. ---}\enspace
}

\pdfstringdefDisableCommands{\def\textsuperscript#1{#1}}

% As páginas impressas reiniciam em diferentes exposições e erratas.
% A chave apenas distingue os destinos invisíveis; não altera os textos.
\newcommand{\FgaUnitKey}{frontmatter}
\renewcommand{\src}[1]{%
  \phantomsection\hypertarget{src:\FgaUnitKey:#1}{}%
  \marginpar{\raggedleft\footnotesize\sffamily #1}%
}
\setlength{\emergencystretch}{3em}
\renewcommand{\boldsymbol}[1]{\symbfit{#1}}
\hypersetup{pdftitle={Fundamentos da geometria algébrica -- Edição atual corrigida -- volume B},pdfauthor={Alexandre Grothendieck},pdfsubject={Exposições 212, 221, 232 e 236, erratas e comentários}}
\begin{document}
\frontmatter
\begin{titlepage}
\centering
\vspace*{0.13\textheight}
{\Huge\bfseries Fundamentos da\par geometria algébrica\par}
\vspace{1.5\baselineskip}
{\Large Alexandre Grothendieck\par}
\vfill
{\Large Edição atual corrigida -- volume B\par}
\vspace{\baselineskip}
{\large Exposições 212, 221, 232 e 236, erratas e comentários\par}
\vspace{1.5\baselineskip}
{\normalsize Edição de 24 de outubro de 2026\par}
\end{titlepage}
\chapter*{Nota sobre esta edição}
\addcontentsline{toc}{chapter}{Nota sobre esta edição}
Exposições 212, 221, 232 e 236, erratas e comentários. Edição corrigida.

Os números marginais remetem às páginas das exposições originais.

A obra original, Fondements de la géométrie algébrique, é de Alexandre Grothendieck. A tradução humana para o inglês de \href{https://github.com/thosgood/fga/tree/9976dd5f8e2ddcb388f0a07f3aa334f120d1019f}{Tim Hosgood} constitui uma fonte de referência desta edição. Seu projeto reconhece as correções de Steve Hnizdur e o apoio técnico de Jon Sterling. A versão 9976dd5f8e2ddcb388f0a07f3aa334f120d1019f identifica o texto inglês usado na comparação de procedência; não estabelece a versão histórica exata incorporada inicialmente. A tradução brasileira e as correções e composição posteriores foram produzidas por IA. A autoria da tradução inglesa permanece com seu tradutor; não se lhe atribui a tradução brasileira. Esta atribuição não constitui autorização nem endosso.

A tradução para o português é produzida por IA. As correções, a edição e a composição também são produzidas por IA. O responsável pela publicação não traduziu o texto. A autoria e as atribuições das obras originais e das contribuições herdadas são preservadas; não são atribuídas à IA. As correções documentadas desta edição incluem OpenAI Codex, GPT-6 Astra, esforço Ultra.

Os direitos sobre as obras originais e sobre as contribuições protegidas a estas traduções e transcrições derivadas permanecem com seus respectivos titulares. Tudo aquilo sobre o qual o editor detenha direitos nestas edições é dedicado ao domínio público sob CC0 1.0 (https://creativecommons.org/publicdomain/zero/1.0/deed.pt-br). Os avisos de direitos de terceiros e as autorizações já concedidas continuam aplicáveis.

Estas edições são disponibilizadas gratuitamente, para fins não comerciais e exclusivamente educativos, sob a presunção de que os respectivos titulares dos direitos não se oporiam à disponibilidade gratuita destas traduções e transcrições, ao lado dos originais franceses no \href{https://www.numdam.org/}{NUMDAM}. Essa presunção não significa que uma autorização tenha sido obtida; o acesso gratuito aos originais não autoriza, por si só, sua tradução, adaptação ou republicação.

Se você detém direitos sobre algum material apresentado aqui e se opõe à sua disponibilidade, poderá \href{https://github.com/KokunoYumeto/ega-fga-sga-pt-br/issues}{abrir uma questão no GitHub}. O responsável pela publicação retirará o material afetado a pedido do titular dos direitos.


\tableofcontents
\mainmatter
\def\FgaUnitKey{212}\unit{Exposição 212 -- Pré-esquemas quocientes}
\sid{212}
\markboth{Exposição 212}{Exposição 212}
\src{99}
\seminares{13.º ano, 1960/61, n.º 212\\Fevereiro de 1961}
\paperhead{TÉCNICAS DE CONSTRUÇÃO E TEOREMAS DE EXISTÊNCIA EM GEOMETRIA ALGÉBRICA\\
III. PRÉ-ESQUEMAS QUOCIENTES}%
  {por Alexander GROTHENDIECK}

\subsection*{Introdução}
\addcontentsline{toc}{subsection}{Introdução}

Os problemas tratados na presente exposição diferem dos considerados nas duas
anteriores, pois se procura representar certos funtores covariantes, e não mais
contravariantes, de esquemas variáveis. O procedimento de passagem ao quociente
é, contudo, essencial em muitos problemas de construção em geometria algébrica,
inclusive os das Exposições I e II ([1], [2]). Assim, a questão da
\emph{efetividade de um dado de descida} sobre um $T$-pré-esquema $X$,
relativamente a um morfismo fielmente plano e quase-compacto
$T\longrightarrow S$, equivale à questão da existência de um quociente de $X$
(que satisfaça as propriedades razoáveis examinadas adiante) pela relação de
equivalência plana em $X$ definida pelo dado de descida; as questões levantadas
em [1], A, \S~2 c serão, sem dúvida, resolvidas ao mesmo tempo que as questões
formuladas no parágrafo 2 da presente exposição. Do mesmo modo, o
\emph{esquema de Picard} de um $S$-esquema $X$ (para a definição, ver [2], C,
\S~3) pode ser definido de várias maneiras como quociente de certos outros
esquemas (de divisores efetivos, ou de imersões em um esquema projetivo) por
relações de equivalência planas, sendo tecnicamente mais simples a definição e
a construção desses esquemas auxiliares: são esquemas da forma
$\underline{\operatorname{Hom}}_S(X,Y)$ e variantes definidos em [2], C,
\S~2, cuja construção será objeto da exposição seguinte (sob hipóteses
adequadas de projetividade). Combinando, portanto, os resultados da presente
exposição e da seguinte, obtém-se a construção dos esquemas de Picard sob
hipóteses adequadas.

O problema da passagem ao quociente na categoria dos pré-esquemas ainda
apresenta várias questões não resolvidas. A mais importante é mencionada no
parágrafo 8. Atualmente, ela é o único obstáculo à construção de
\emph{esquemas de moduli sobre os inteiros para curvas de gênero arbitrário,
variedades abelianas polarizadas}, etc. Isso mostra que sua solução merece os
esforços dos especialistas em grupos algébricos.
\src{100}
\subsection*{1. Relações de equivalência e relações de equivalência efetivas}
\addcontentsline{toc}{subsection}{1. Relações de equivalência e relações de equivalência efetivas}

Sejam $\mathcal C$ uma categoria e $X$ um objeto de $\mathcal C$. Um par de
morfismos
\[
  p_1,p_2:R\rightrightarrows X
\]
é chamado de \emph{par de equivalência} em $\mathcal C$, com \emph{alvo} $X$ e
\emph{fonte} $R$, se, para todo objeto $T$ de $\mathcal C$, as duas aplicações
correspondentes
\[
  p_1(T),p_2(T):R(T)\rightrightarrows X(T)
\]
(onde, para todo objeto $Y$ de $\mathcal C$, se põe
$Y(T)=\operatorname{Hom}(T,Y)$) definem uma aplicação
\[
  R(T)\longrightarrow X(T)\times X(T)
\]
que induz uma bijeção de $R(T)$ sobre o gráfico de uma relação de equivalência
no conjunto $X(T)$. Introduz-se, entre os pares de equivalência com alvo $X$,
uma relação de equivalência evidente; uma classe de equivalência para esta
última é chamada \emph{$\mathcal C$-relação de equivalência em $X$}, ou
simplesmente uma relação de equivalência quando não houver risco de confusão.

Se $X\times X$ existe, dar uma relação de equivalência em $X$ equivale a dar um
subobjeto $R$ de $X\times X$ tal que, para todo objeto $T$ de $\mathcal C$, o
subconjunto de
\[
  (X\times X)(T)=X(T)\times X(T)
\]
correspondente a $R(T)$ seja o gráfico de uma relação de equivalência em
$X(T)$. Denotando por $p_1$ e $p_2$ os morfismos de $R$ em $X$ induzidos pelas
projeções $\operatorname{pr}_1$ e $\operatorname{pr}_2$, a condição anterior
exprime que $(p_1,p_2)$ é um par de equivalência. Também se podem exprimir
diagramaticamente em $\mathcal C$ (desde que existam $X\times X$ e o produto
fibrado $(R,p_2)\times_X(R,p_1)$) os axiomas conjuntistas de uma relação de
equivalência para $R(T)$ em $X(T)$, de acordo com o princípio geral
explicado em [2], A, \S~1. Não precisaremos disso.

Sempre que se tem um par de morfismos $(p_1,p_2)$ com a mesma fonte $R$ e o
mesmo alvo $X$, pode-se definir o \emph{conúcleo} do par como o objeto $Y$ de
$\mathcal C$ que representa o funtor covariante em $Z$
\[
  \operatorname{Hom}_{p_1,p_2}(X,Z),
\]
conjunto dos morfismos $u$ de $X$ em $Z$ tais que $up_1=up_2$. Se $Y$ existe,
fica determinado a menos de um isomorfismo único. Será denotado por
$X/(p_1,p_2)$ ou, por abuso de notação, por $X/R$, sendo esta última notação
empregada sobretudo quando
\src{101}$(p_1,p_2)$ é um par de equivalência: é então usual identificar,
na notação, a relação de equivalência definida pelo par com $R$. Observe-se
que, se $Y$ for considerado um quociente de $X$, ele depende, de fato, apenas
da relação de equivalência definida pelo par $(p_1,p_2)$.

Partamos agora de um morfismo
\[
  f:X\longrightarrow Y
\]
que permite, portanto, considerar $X$ como um «objeto sobre $Y$», e suponhamos
que exista o produto fibrado
\[
  R(f)=X\times_Y X
\]
Sejam $p_1$ e $p_2$ suas duas projeções. Então $(p_1,p_2)$ é um par de
equivalência, chamado par \emph{associado} ao morfismo $f$. Ele define, assim,
uma relação de equivalência, chamada relação associada a $f$.

Um par de morfismos $(p_1,p_2)$ com alvo $X$ e fonte comum $R$ é chamado
\emph{par de equivalência efetivo} se
\begin{enumerate}
  \item[(i)] o conúcleo $X/(p_1,p_2)=Y$ existe;
  \item[(ii)] o produto fibrado $X\times_Y X$ existe;
  \item[(iii)] o morfismo $R\longrightarrow X\times_Y X$ cujas componentes são
  $p_1$ e $p_2$ é um isomorfismo.
\end{enumerate}
Então $(p_1,p_2)$ é de fato um par de equivalência. A relação de equivalência
que ele define é também chamada \emph{relação de equivalência efetiva}.

Um morfismo $f:X\longrightarrow Y$ é chamado \emph{epimorfismo efetivo} se
\begin{enumerate}
  \item[(i)] o produto fibrado $X\times_Y X=R$ existe;
  \item[(ii)] o quociente $X/(p_1,p_2)$ existe, onde $p_1$ e $p_2$ são as duas
  projeções de $R$ em $X$;
  \item[(iii)] o morfismo $X/(p_1,p_2)\longrightarrow Y$ induzido por $f$ é um
  isomorfismo.
\end{enumerate}
Então $f$ é de fato um epimorfismo, e até um epimorfismo estrito (cf. [1], A,
\S~2.3), sendo a recíproca verdadeira se o produto fibrado $X\times_Y X$
existir. Diz-se também que o objeto quociente de $X$ definido pelo epimorfismo
$f$ é um \emph{quociente efetivo} de $X$.

As definições precedentes implicam a seguinte «correspondência de Galois»:
\src{102}
\result{PROPOSIÇÃO}{1.1} Há uma correspondência biunívoca, que preserva as
ordens naturais, entre o conjunto das relações de equivalência efetivas $R$
em $X$ e o conjunto dos quocientes efetivos $Y$ de $X$: a $R$ corresponde o
quociente efetivo $X/R$, e a $Y$ corresponde a relação de equivalência efetiva
definida pela projeção canônica $X\longrightarrow Y$ (que é definida pelo
produto fibrado $X\times_Y X$ munido de suas duas projeções).

Nas categorias muito boas (conjuntos, feixes de conjuntos etc.), todo
quociente é efetivo e toda relação de equivalência é efetiva. Isso já não
ocorre em categorias como a categoria dos pré-esquemas sobre um pré-esquema
$S$ dado, mesmo quando $S$ é o espectro de um corpo e mesmo quando nos
restringimos aos esquemas finitos sobre $S$. As questões de efetividade e até
mesmo (no caso de pré-esquemas não finitos sobre $S$) as questões relativas à
existência de quocientes revelam-se quase sempre delicadas.

\subsection*{2. Exemplo: pré-esquemas finitos sobre $S$}
\addcontentsline{toc}{subsection}{2. Exemplo: pré-esquemas finitos sobre S}

Seja $\mathcal C$ a categoria dos pré-esquemas finitos sobre $S$, supondo-se
$S$ localmente Noetheriano. Ela é equivalente à categoria oposta da categoria
dos feixes coerentes de álgebras comutativas sobre $S$, ou ainda, se $S$ é
afim com anel $A$, à categoria oposta da categoria das $A$-álgebras finitas
sobre $A$ (isto é, finitamente geradas como $A$-módulos). Conclui-se
imediatamente que existem em $\mathcal C$ os limites projetivos finitos e os
limites indutivos finitos. Isso é bem conhecido (sem qualquer hipótese de
finitude\ldots) quanto aos primeiros. Assim, o produto fibrado de pré-esquemas
$X,Y$ sobre $S$ corresponde ao produto tensorial $B\otimes_A C$ das álgebras
correspondentes; o núcleo de dois morfismos $X\rightrightarrows Y$,
correspondentes a dois homomorfismos de $A$-álgebras
$u,v:C\rightrightarrows B$, corresponde ao quociente de $B$ pelo ideal gerado
pelos elementos $u(c)-v(c)$ etc. Para os limites indutivos finitos, basta
considerar, de um lado, as somas finitas, que correspondem ao produto ordinário
de $A$-álgebras, e, de outro, os conúcleos de pares de morfismos
$X\rightrightarrows Y$, que correspondem de fato (como se verifica
imediatamente) ao subanel de $C$ formado pelos elementos nos quais os
homomorfismos $u,v:C\rightrightarrows B$ coincidem (este último é finito sobre
$A$ graças à hipótese Noetheriana). Observe-se, aliás, que se pode mostrar,
utilizando a hipótese Noetheriana, que os limites indutivos finitos e, em
particular, os quocientes assim construídos na categoria $\mathcal C$ dos
pré-esquemas finitos sobre $S$ são, na realidade, quocientes na categoria de
\emph{todos} os pré-esquemas.
\src{103}
Como dissemos em [1], há em $\mathcal C$ epimorfismos não efetivos (ou não
estritos\footnote{O texto impresso omite «não» antes de «estritos». Na
categoria considerada existem produtos fibrados, e o parágrafo anterior
estabelece precisamente a equivalência entre efetividade e epimorfia estrita.};
dá no mesmo, pois existem produtos fibrados). \emph{Ignoro se as relações de
equivalência são sempre efetivas} quando não se impõe nenhuma hipótese de
planicidade. Nessa direção, obtive apenas resultados positivos muito parciais,
indispensáveis à demonstração do teorema fundamental da teoria formal de
moduli (ver [2], B, teo. 1). Assinalemos que, no problema proposto, é fácil
reduzir-se ao caso em que $S$ é o espectro de um anel artiniano local $A$, com
corpo residual algebricamente fechado. Mas, mesmo se $A$ for um corpo, a
resposta não é conhecida.

Também se pode considerar um pré-esquema $X$ sobre $S$ que já não se suponha
finito sobre $S$, mas que seja munido de uma relação de equivalência $R$ em
$X$ tal que $p_1:R\longrightarrow X$ seja um morfismo finito. Diz-se então que
$R$ é uma relação de equivalência finita. Suponhamos, para simplificar, que $S$
e $X$ sejam afins (logo, $R$ é afim, de modo que a situação se reduz a uma
questão de álgebra comutativa pura). \emph{Mesmo nesse caso, ignora-se se
existe um quociente $X/R=Y$ e se o morfismo canônico $X\longrightarrow Y$ é
finito}. (O caso mais simples é aquele em que $S$ é o espectro de um corpo $k$
e $X$ é o espectro de $k[t]$, isto é, a reta afim.) Naturalmente, se os dois
problemas precedentes tiverem soluções afirmativas, pode-se concluir, na
situação presente, que $R$ é efetiva. Observemos que o problema da
\emph{existência} de um quociente $Y$ e da finitude de
$f:X\longrightarrow Y$ se apresenta exatamente nos mesmos termos se, em vez
do grafo de uma relação de equivalência em $X$, houver apenas um pré-grafo de
equivalência em $X$, no sentido do parágrafo 4.

A questão da passagem ao quociente por uma relação de equivalência finita mais
ou menos arbitrária surge na construção de pré-esquemas por «colagem» de
pré-esquemas dados $X_i$ ao longo de certos subpré-esquemas fechados; a lei de
colagem exprime-se precisamente por uma relação de equivalência finita no
pré-esquema soma $X$ dos $X_i$. Deve-se esperar também que a solução dos
problemas aqui propostos, e de diversas variantes, seja uma condição prévia ao
desenvolvimento de uma técnica geral de construções não projetivas, na direção
iniciada em [2].

O único fato positivo geral conhecido pelo autor é o seguinte:

\result{PROPOSIÇÃO}{2.1} Sejam $S$ um pré-esquema localmente noetheriano, $s$
um ponto de $S$ e $\Omega$ uma extensão algebricamente fechada de $k(s)$.
Considere o
\src{104}«funtor fibra» correspondente $F$, que associa a todo $S$-esquema
$X$ finito sobre $S$ o conjunto dos pontos de $X_s$ com valores em $\Omega$.
Esse funtor (trivialmente exato à esquerda) é \emph{exato à direita}, isto é,
comuta com os limites indutivos finitos e, em particular, com os conúcleos de
pares de morfismos.

Usando esse resultado para todos os «pontos geométricos» de $S$, deduz-se que
a categoria «quociente» $\mathcal C'$ de $\mathcal C$, obtida raciocinando
«módulo morfismos radiciais sobrejetivos» (isto é, adjungindo formalmente
inversos aos morfismos em questão), é uma categoria «geométrica», isto é,
satisfaz as mesmas propriedades «de natureza finita» que a categoria dos
conjuntos. Em particular, toda relação de equivalência nela é efetiva. Isso
implica que, se $R$ for uma relação de equivalência em um $X$ finito sobre
$S$, então o morfismo canônico
\[
  R\longrightarrow X\times_Y X\qquad (\text{onde }Y=X/R)
\]
é \emph{radicial e sobrejetivo} (de fato, uma imersão fechada sobrejetiva,
pois é um monomorfismo).

\subsection*{3. O caso de um grupo de operadores}
\addcontentsline{toc}{subsection}{3. O caso de um grupo de operadores}

Suponhamos novamente que $\mathcal C$ seja uma categoria qualquer. Sejam
$G,X$ objetos de $\mathcal C$, e suponhamos que $G$ seja um
$\mathcal C$-grupo de operadores sobre o objeto $X$ de $\mathcal C$. Isso
significa (cf. [2], A, \S~1) que, para todo objeto $T$ de $\mathcal C$, foram
dadas uma estrutura de grupo em $G(T)$ e uma estrutura de $G(T)$-conjunto em
$X(T)$, de tal modo que, variando $T$, as estruturas em questão «variam
funtorialmente» com $T$. Se os produtos $G\times G$ e $G\times X$ existirem
em $\mathcal C$, tal estrutura também poderá ser definida por um par de
morfismos
\[
  G\times G\longrightarrow G,
  \qquad \pi:G\times X\longrightarrow X,
\]
sujeito à condição de que, para todo objeto $T$ de $\mathcal C$, as leis de
composição correspondentes para os conjuntos $G(T)$ e $X(T)$ façam de $G(T)$
um grupo que age sobre $X(T)$. A tradução desse axioma pela comutatividade de
certos diagramas em $\mathcal C$ é fácil, mas fastidiosa e, de fato,
perfeitamente inútil em todos os casos que conheço.

Suponhamos que $G\times X$ exista e consideremos os dois morfismos
\[
  p_1,p_2:G\times X\rightrightarrows X
\]
dados por\footnote{O texto impresso apresenta \texttt{pr\_1}. Como os dois
morfismos têm codomínio $X$, o primeiro deles deve ser
\texttt{pr\_2}; corrige-se aqui o erro tipográfico.}
\src{105}
\[
  p_1=\operatorname{pr}_2,\qquad p_2=\pi.
\]
Verifica-se imediatamente que $(p_1,p_2)$ é um par de equivalência se, e
somente se, para todo objeto $T$ de $\mathcal C$, a aplicação
\[
  G(T)\times X(T)\simeq(G\times X)(T)
  \longrightarrow X(T)\times X(T),
\]
definida por esse par é injetiva, isto é, se o grupo $G(T)$ age
\emph{livremente} sobre o conjunto $X(T)$, ou seja, se $g\in G(T)$,
$x\in X(T)$ e $g\cdot x=x$ implicam que $g$ é o elemento identidade de
$G(T)$. Diz-se então que $G$ \emph{age livremente sobre $X$} (ou que $X$ é um
\emph{$\mathcal C$-espaço principal sob $G$}). A relação de equivalência
associada ao par $(p_1,p_2)$ chama-se então \emph{relação de equivalência
definida pelo grupo $G$ que age livremente sobre $X$}. Quando $X\times X$
também existe e se considera o morfismo
\[
  p:G\times X\longrightarrow X\times X
\]
definido pelo par $(p_1,p_2)$, a condição de que $G$ aja livremente significa
que $p$ é um \emph{monomorfismo}.

Naturalmente, mesmo quando $G$ não age livremente sobre $X$, procuram-se
critérios para a existência de um quociente de $X$ por $G$, isto é, do
conúcleo do par anterior $(p_1,p_2)$.

O conúcleo em questão será muitas vezes denotado por $X/G$, ou,
preferencialmente, por $G\backslash X$ se $G$ agir à esquerda (reservando-se a
notação anterior para o caso em que $G$ age à direita). Observe-se que, mesmo
quando existe a «imagem» de $G\times X$ por $p$ (definida, por exemplo, como o
menor subobjeto de $X\times X$ pelo qual $p$ se fatora), digamos $R$, ela quase
nunca é uma relação de equivalência em $X$. Se então se tenta passar
diretamente ao quociente por $R$ (ou, mais precisamente, pelo par de morfismos
de $R$ em $X$ induzidos pelas duas projeções $\operatorname{pr}_i$), perdem-se
irremediavelmente as características particulares do par original
$(p_1,p_2)$. É, portanto, importante encontrar uma generalização da noção de
relação de equivalência que se aplique diretamente ao par definido por um
grupo de operadores de $\mathcal C$.

\subsection*{4. Pré-relações de equivalência}
\addcontentsline{toc}{subsection}{4. Pré-relações de equivalência}

Recordemos que uma categoria na qual todo morfismo é um isomorfismo se chama
\emph{grupoide}. Por outro lado, uma categoria deve ser definida como formada
por dois conjuntos subjacentes $(X,R)$, o conjunto de \emph{objetos} e o
conjunto de \emph{setas}, munidos das seguintes estruturas:
\src{106}
\begin{enumerate}
\item[(i)] Um par de aplicações
\[
  p_1,p_2:R\rightrightarrows X
\]
chamadas \emph{aplicação fonte} e \emph{aplicação alvo},

\item[(ii)] Uma aplicação
\[
  \pi:(R,p_2)\times_X(R,p_1)\longrightarrow R
\]
chamada \emph{aplicação de composição}.
\end{enumerate}
Esses dados devem satisfazer axiomas bem conhecidos, que não repetiremos aqui
e que poderiam ser explicitados pela comutatividade de certos diagramas e pela
existência de uma aplicação $D:X\to R$ (necessariamente única) que torna
comutativos outros dois diagramas. $D$ corresponde à passagem de um objeto à
seta identidade correspondente e satisfaz
\[
  p_1D=p_2D=\operatorname{id}_X.
\]
Dizer que a categoria é um grupoide significa então que existe uma aplicação
\[
  s:R\longrightarrow R
\]
(necessariamente única), chamada a \emph{inversão} de $R$, que transforma cada
seta em uma seta inversa, o que poderia ser explicitado pela comutatividade de
outros quatro diagramas, formados com $s$, $\Delta$ e os dados precedentes, e
dos quais os dois primeiros se escrevem:
\[
  p_1s=p_2,\qquad p_2s=p_1.
\]

Recordados esses pontos, as definições gerais em ([2], A, \S~1) mostram, em
particular, o que se deve entender por uma estrutura de
$\mathcal C$-\emph{categoria}, respectivamente de
$\mathcal C$-\emph{grupoide}, sobre um par de objetos $X,R$ de uma categoria
qualquer $\mathcal C$: trata-se, por definição, do dado, para todo objeto $T$
de $\mathcal C$, de uma estrutura de categoria, respectivamente de grupoide,
no sentido conjuntista, cujo conjunto de objetos é $X(T)$ e o conjunto de
setas é $R(T)$, essas estruturas
\guillemotleft variando funtorialmente\guillemotright{} com $T$ variável.
Isso implica, portanto, a definição de dois morfismos
\[
  p_1,p_2:R\rightrightarrows X
\]
chamados \emph{morfismo fonte} e \emph{morfismo alvo} e, quando existe o
produto fibrado desejado, de um morfismo
\src{107}
\[
  \pi:(R,p_2)\times_X(R,p_1)\longrightarrow R
\]
chamado \emph{morfismo de composição}. Esses três dados bastam então para
determinar a estrutura de categoria sobre $X,R$, sendo o axioma a impor a
esses dados o seguinte: para todo $T$, os três dados correspondentes para
$X(T),R(T)$ definem sobre esse par de conjuntos uma estrutura de categoria
(respectivamente, de grupoide). Quando for o caso, isso pode ser expresso pela
comutatividade de certos diagramas que envolvem um morfismo univocamente
determinado
\[
  D:X\longrightarrow R
\]
e, no caso dos grupoides, um morfismo univocamente determinado
\[
  s:R\longrightarrow R;
\]
diagramas que se explicitam como no caso «conjuntista». Essa interpretação
trabalhosa dos axiomas é, felizmente, inútil na prática. O único interesse
teórico da possibilidade de expressar os dados e os axiomas a partir de
morfismos e igualdades de morfismos entre certos produtos fibrados é o
seguinte: se temos um funtor $F:\mathcal C\to\mathcal C'$ exato à esquerda
(isto é, que comuta com produtos finitos e produtos fibrados), ele transforma
uma $\mathcal C$-categoria (respectivamente, um $\mathcal C$-grupoide) em uma
$\mathcal C'$-categoria (respectivamente, um $\mathcal C'$-grupoide), desde
que existam produtos finitos e produtos fibrados em $\mathcal C$.

Na prática, é importante saber interpretar os morfismos
$p_1,p_2,\pi,D,s$ como \emph{operações simpliciais} em um objeto
semissimplicial ou simplicial adequado de $\mathcal C$, ao menos quando
existem produtos fibrados em $\mathcal C$. Para fixar a terminologia,
introduzamos a \emph{categoria $\mathcal S$ dos simplexos padrão} como a
categoria cujos objetos são os conjuntos finitos da forma
\[
  \Delta_n=[0,n],\qquad\text{onde }n\geq 0
\]
(intervalo dos inteiros de $0$ a $n$), e cujos morfismos ou setas são as
\emph{aplicações arbitrárias} entre esses conjuntos finitos. Observe-se que a
categoria $\mathcal S$ é equivalente à categoria dos conjuntos finitos
\emph{não vazios}, tomando-se como morfismos as aplicações entre conjuntos
finitos. Em $\mathcal S$, existe evidentemente o coproduto de uma família
finita \emph{não vazia} de objetos, bem como o coproduto fibrado de dois
objetos sob um terceiro (operação dual ao produto fibrado). Denotemos por
$\mathcal S'$ a subcategoria de $\mathcal S$ que tem os mesmos objetos, mas
cujos morfismos são as \emph{aplicações crescentes} entre os $\Delta_n$. Essa
categoria é equivalente à categoria dos conjuntos finitos não vazios
totalmente ordenados. Nesta categoria, o
\src{108}
coproduto de dois objetos nunca existe, e o coproduto fibrado de dois objetos
$A,B$ sob um terceiro $C$ não existe em geral (tome-se, por exemplo,
$C=\Delta_0$, $A=B=\Delta_1$, com as duas aplicações estruturais
$u:C\to A$ e $v:C\to B$ iguais). Contudo, ele existe em certos casos; por
exemplo,
\[
  C=\Delta_0,\qquad A=\Delta_m,\qquad B=\Delta_n,\qquad
  u(0)=m,\qquad v(0)=0,
\]
e, nesse caso,
\[
  \Delta_m\amalg_{\Delta_0}\Delta_n=\Delta_{m+n}.
\]

Um \emph{objeto simplicial} (respectivamente, um \emph{objeto
semissimplicial}) em uma categoria $\mathcal C$ é, por definição, um funtor
contravariante $K$ de $\mathcal S$ (respectivamente, de $\mathcal S'$) em
$\mathcal C$. Portanto, um objeto simplicial define por restrição um objeto
semissimplicial, mas a primeira noção difere essencialmente da segunda pela
presença de \emph{operações de simetria} nos $K_n=K(\Delta_n)$,
correspondentes às imagens, pelo funtor $K$, dos elementos do grupo simétrico
de $n+1$ letras (considerado como o grupo de automorfismos de $\Delta_n$ em
$\mathcal S$).

Para todo $n$, seja $\Delta'_n$ (respectivamente, $\Delta''_n$) a categoria
finita cujo conjunto de objetos é $\Delta_n$, definida pela relação de ordem
caótica (respectivamente, pela ordem total natural) em $\Delta_n$ (sendo o
conjunto de setas o grafo dessa relação de ordem). É evidente que $\Delta'_n$
(respectivamente, $\Delta''_n$) depende funtorialmente do objeto $\Delta_n$ de
$\mathcal S$ (respectivamente, de $\mathcal S'$). Se $Z$ for uma categoria,
então $\operatorname{Hom}(\Delta'_n,Z)$ (respectivamente,
$\operatorname{Hom}(\Delta''_n,Z)$), quando $\Delta_n$ varia, é um funtor de
$\mathcal S$ (respectivamente, de $\mathcal S'$) para a categoria dos
conjuntos, isto é, um conjunto simplicial (respectivamente, um conjunto
semissimplicial), dito \emph{associado à categoria $Z$}; denotemo-los por $Z'$
e $Z''$. Há, além disso, um homomorfismo natural evidente do conjunto
semissimplicial associado a $Z'$ para $Z''$, e este último é um isomorfismo se,
e somente se, $Z$ for um grupoide. Posto isso:

\result{PROPOSIÇÃO}{4.1} O funtor $Z\longmapsto Z''$ da categoria das
categorias para a categoria dos conjuntos semissimpliciais é plenamente fiel
e define uma equivalência entre a categoria das \emph{categorias} e a
categoria dos conjuntos semissimpliciais, isto é, dos funtores contravariantes
$K$ de $\mathcal S'$ para $(\mathrm{Set})$ que transformam os coprodutos
fibrados $A\amalg_C B$ do tipo especificado acima em produtos fibrados de
conjuntos. Do mesmo modo, o funtor $Z\longmapsto Z'$ da categoria dos
grupoides para a categoria dos conjuntos simpliciais é plenamente fiel e
define uma equivalência entre a categoria dos \emph{grupoides} e a categoria
dos
\src{109}
conjuntos simpliciais, isto é, funtores contravariantes $K$ de $\mathcal S$
em $(\mathrm{Ens})$, que transformam somas amalgamadas em produtos fibrados.

Pode-se, portanto, considerar as categorias como conjuntos semissimpliciais
particulares, e os grupoides como conjuntos simpliciais particulares, com a
condição, naturalmente, de raciocinar «a menos de isomorfismo», como é de rigor
quando se interpretam certas estruturas em termos de outras. O procedimento
usual de redução ao caso conjuntista implica então:

\result{COROLÁRIO}{4.2} O enunciado precedente continua válido quando se
substituem as categorias, os grupoides e os conjuntos simpliciais por
$\mathcal C$-categorias, $\mathcal C$-grupoides e $\mathcal C$-objetos
simpliciais, contanto que existam produtos fibrados em $\mathcal C$.

O objeto semissimplicial $K$ em $\mathcal C$ associado a uma categoria
$(X,R,\ldots)$ em $\mathcal C$ pode ser explicitado interpretando-se a
componente $K_n=K(\Delta_n)$ de $K$ como a $n$-ésima potência fibrada de
$(R,p_1)$ sobre $X$, ou melhor, pela fórmula de recorrência
\[
  K_0=X,\qquad
  K_n=(K_{n-1},p_{n-1}^{(n-1)})\times_X(R,p_1),
\]
onde os $p_i^{(n-1)}$ ($0\leq i\leq n-1$) são as projeções naturais de
$K_{n-1}$ em $X$ (que também se definem por recorrência). Dessa forma,
$p_1,p_2,\pi,D,s$ são interpretados como as operações simpliciais que
correspondem respectivamente aos morfismos em $\mathcal S$: a face $0$ de
$\Delta_1$, a face $1$ de $\Delta_1$, a face $(0,2)$ de $\Delta_2$, a
degenerescência $\Delta_1\to\Delta_0$ e a simetria de $\Delta_1$. Todas as
demais operações semissimpliciais (respectivamente simpliciais) deduzem-se
formalmente das quatro (respectivamente cinco) precedentes por composição e
produtos fibrados.

Chama-se agora \emph{pré-relação de equivalência} em um objeto $X$ de uma
categoria o dado de um grupoide cujo objeto dos objetos (se assim se pode
dizer) é $X$. Tal dado compreende, portanto, entre outras coisas, um objeto $R$
e dois morfismos
\[
  p_1,p_2:R\rightrightarrows X.
\]
Observe-se, porém, que esses dados, por si sós, não determinam a estrutura
considerada, ao contrário do que ocorre com os pares de equivalência. Nesta
exposição, interessamo-nos por essa noção com o objetivo de obter critérios de
possibilidade de passagem ao quociente, isto é, de formação do conúcleo do par
$(p_1,p_2)$. O enunciado do problema, portanto, não faz uso dos dados
adicionais inerentes a um grupoide. Na demonstração dos resultados que se
seguem,
\src{110}
utilizaremos, contudo, esses dados adicionais e, em particular, as operações
simpliciais (incluindo as operações de simetria) até a dimensão $3$ (nas quais
intervém a quarta potência fibrada de $R$ sobre $X$).

Uma relação de equivalência em um objeto $X$ de $\mathcal C$ define uma
pré-relação de equivalência: basta verificá-lo no caso conjuntista, no qual se
associa a uma relação de equivalência em um conjunto $X$ o grupoide cujo
conjunto de objetos é $X$ e cujo conjunto de setas é o grafo da relação de
equivalência.

Um $\mathcal C$-monoide $G$ que opera sobre um objeto $X$ de $\mathcal C$
define uma $\mathcal C$-categoria cujos objetos subjacentes são $R=G\times X$
e $X$ (desde que $G\times X$ exista), e que é um $\mathcal C$-grupoide se, e
somente se, $G$ for um grupo. Basta, ainda, fazer a verificação no caso
conjuntista. Define-se então a composição das setas $(g,a)$ e
$(g',g\cdot a)$ por
\[
  (g',g\cdot a)(g,a)=(g'g,a),
\]
isto é, se $a,b\in X$, então $\operatorname{Hom}(a,b)$ é, por definição, o
transportador de $a$ em $b$, e os morfismos se compõem mediante a composição
dos elementos de $G$.

\medskip
\noindent\textbf{OBSERVAÇÃO. ---} Podem-se evitar as dificuldades lógicas
levantadas por um enunciado como (4.1), subentendendo-se que todos os objetos
considerados se encontram em um «Universo» fixado (que é ele próprio um
conjunto).

\fsection{5}{Quociente por uma relação de equivalência finita e plana}

\result{TEOREMA}{5.1} Sejam $X=\operatorname{Spec}(B)$ um esquema afim e $R$
uma pré-relação de equivalência em $X$, cuja componente $R_1$ é afim, digamos
$R_1=\operatorname{Spec}(C)$. Suponha-se que a primeira projeção
$p_1:R_1\to X$ seja um morfismo finito e localmente livre; em outras palavras,
que o homomorfismo de anéis correspondente $p_1^*:B\to C$ faça de $C$ um
$B$-módulo projetivo de tipo finito. Seja $A$ o subanel de $B$ que é o núcleo
do par de homomorfismos
\[
  p_1^*,p_2^*:B\rightrightarrows C
\]
(isto é, o conjunto dos elementos $b$ tais que
$p_1^*(b)=p_2^*(b)$). Seja $Y=\operatorname{Spec}(A)$, e seja $f:X\to Y$ o
morfismo definido pela inclusão de $A$ em $B$. Nessas condições:
\begin{enumerate}
\item[(i)] $B$ é integral sobre $A$, isto é, $f$ é um morfismo integral.

\item[(ii)] O morfismo $f$ é sobrejetivo, e suas fibras são as classes de
equivalência
\src{111}
conjuntistas $p_2\allowbreak(p_1^{-1}(x))$ em $X$ módulo $R$, e a topologia
de $Y$ é a topologia quociente da de $X$.

\item[(iii)] $Y$ é o quociente de $X$ por $R$ na categoria dos pré-esquemas.

\item[(iv)] Se $R$ provém de uma \emph{relação de equivalência}, então o
morfismo $f:X\to Y$ é finito e localmente livre (isto é, $B$ é um $A$-módulo
projetivo finitamente gerado), e a relação de equivalência é efetiva, isto é,
$R_1\to X\times_Y X$ é um isomorfismo.
\end{enumerate}

Este teorema generaliza o conhecido teorema relativo ao caso de um grupo
finito $G$ que atua por automorfismos no anel $B$ e ao anel de invariantes
$A$, e sua demonstração é análoga à demonstração conhecida. Pode-se precisar
(iii) da seguinte maneira:

\result{COROLÁRIO}{5.2} O morfismo canônico
$R_1\to X\times_Y X$ é \emph{sobrejetivo}.

Seja ainda $R$ uma pré-relação de equivalência «finita e localmente livre» em
$X$, sendo agora $X$ um pré-esquema arbitrário. Suponhamos que se possa
encontrar um pré-esquema $Y$ e um morfismo $f:X\to Y$ tais que $fp_1=fp_2$ e,
além disso, que seja exata a sequência de homomorfismos de feixes de anéis
sobre $Y$:
\[
  \mathcal O_Y\longrightarrow f_*(\mathcal O_X)
  \rightrightarrows g_*(\mathcal O_R)
\]
(onde $g=fp_i$). Resulta então do teorema que se têm conclusões (i) a (iv)
análogas às do teorema; em particular, por (iii), $Y$ é o quociente de $X$ por
$R$ e, portanto, fica determinado a menos de isomorfismo único. Diz-se, nessas
condições, que a pré-relação de equivalência $R$ em $X$ é \emph{admissível}.
Com essa definição:

\result{TEOREMA}{5.3} Sejam $X$ um pré-esquema e $R$ uma pré-relação de
equivalência em $X$ tal que $p_1:R_1\to X$ seja um morfismo finito e localmente
livre. Para que $R$ seja admissível, é necessário e suficiente que toda classe
de equivalência conjuntista $p_2(p_1^{-1}(x))$ em $X$ módulo $R$ esteja
contida em um aberto afim (condição sempre verificada se todo subconjunto
finito de $X$ estiver contido em um aberto afim, por exemplo, se $X$ for
quase-projetivo sobre um esquema afim).

Mostra-se, com efeito, sem dificuldade que toda classe de equivalência módulo
$R$ em $X$ está então contida em um aberto afim \emph{estável por $R$}, e
constrói-se o quociente $Y$ pela colagem das partes obtidas com o auxílio do
Teorema 5.1.
\src{112}

\result{COROLÁRIO}{5.4} Suponhamos satisfeita essa condição e que, além disso,
$R$ provenha de uma relação de equivalência. Então esta última é efetiva, isto
é, $R_1\to X\times_Y X$ é um isomorfismo, e $f:X\to Y$ é um morfismo finito
localmente livre.

Conclui-se imediatamente, por «descida»:

\result{COROLÁRIO}{5.5} Sob as condições de (5.4), $X$ tem posto $n$ em todo
ponto sobre $Y$ se e somente se $(R_1,p_1)$ tem posto $n$ em todo ponto sobre
$X$. Se $X,R_1$ forem $Z$-pré-esquemas e $p_1,p_2$ forem $Z$-morfismos, sendo
portanto $Y$ um $Z$-pré-esquema, então $X$ é plano sobre $Z$ se e somente se
$Y$ o for.

Em resumo:

\result{ESCÓLIO}{} O dado de um morfismo finito, localmente livre e
sobrejetivo $f:X\to Y$ de pré-esquemas equivale ao dado de um pré-esquema $X$
munido de uma relação de equivalência $R$ tal que $p_1:R\to X$ seja finito e
localmente livre, e que toda classe $p_2(p_1^{-1}(x))$ esteja contida em um
aberto afim.

\medskip
\noindent\textbf{OBSERVAÇÕES 5.6.}

\smallskip
\noindent $1^\circ$ Não foi necessário fazer nenhuma hipótese Noetheriana.

\smallskip
\noindent $2^\circ$ Esta noção de passagem ao quociente inclui, como caso
particular, a «descida inseparável» de CARTIER, que corresponde à determinação
dos morfismos finitos localmente livres $f:X\to Y$ tais que
$f_*(\mathcal O_X)$ admita uma $p$-base em relação a $\mathcal O_Y$ (sendo $X$
um pré-esquema dado cujo feixe $\mathcal O_X$ é anulado pelo número primo
$p>0$). Observe-se que o resultado se exprime facilmente sem hipótese de
regularidade sobre os anéis locais e sem supor que $X$ seja um esquema
algébrico sobre um corpo. A teoria de JACOBSON--BOURBAKI é obtida tomando-se
como $X$ o espectro de um corpo de característica $p$.

\smallskip
\noindent $3^\circ$ GABRIEL havia obtido anteriormente um caso particular do
teorema (5.3) na teoria da passagem ao quociente para grupos comutativos
finitos sobre um corpo $k$. [Compare-se (7.4).]

\fsection{6}{Quociente por uma relação de equivalência própria e plana}

\result{TEOREMA}{6.1} Sejam $S$ um pré-esquema localmente Noetheriano, $X$
um $S$-esquema quasi-projetivo, $R$ uma pré-relação de equivalência em $X$,
tal que:
\src{113}
\begin{enumerate}
\item[(a)] $p_1:R_1\to X$ é próprio e plano;
\item[(b)] $R_1\to X\times_S X$ é um morfismo finito ou, o que equivale em
virtude de (a), tem fibras finitas (condição automaticamente satisfeita se
$R$ provém de uma relação de equivalência).
\end{enumerate}
Nessas condições:
\begin{enumerate}
\item[(i)] $X/R=Y$ existe e (se $S$ é Noetheriano e se $R$ provém de uma
relação de equivalência) é quase-projetivo sobre $S$.\footnote{O texto impresso
não contém a segunda hipótese. A errata vinculada à exposição 212 esclarece
que a quase-projetividade parece então demonstrada apenas quando $R$ provém
de uma relação de equivalência.}

\item[(ii)] O morfismo canônico $f:X\to Y$ é sobrejetivo, próprio e aberto;
suas fibras são as classes de equivalência $p_2(p_1^{-1}(x))$ em $X$ módulo
$R$; portanto, $Y$ identifica-se com o espaço topológico quociente de $X$ pela
relação de equivalência conjuntista definida por $R$. Por fim,
$R_1\to X\times_Y X$ é sobrejetivo.

\item[(iii)] Se $R$ provém de uma relação de equivalência, esta última é
efetiva, isto é, $R_1\to X\times_Y X$ é um isomorfismo e, além disso,
$f:X\to Y$ é plano (portanto, fielmente plano).
\end{enumerate}

Para a demonstração, reduz-se a (5.1) por meio de quase-seções adequadas de $X$
para $R$, sendo a demonstração análoga à construção dos grupos algébricos
quocientes no Seminário Chevalley.

Em resumo:

\result{ESCÓLIO}{} Seja $X$ quase-projetivo sobre $S$ localmente Noetheriano. O
dado de um morfismo próprio, fielmente plano e sobrejetivo $f:X\to Y$ de $X$
para um $S$-pré-esquema $Y$ equivale ao dado de uma relação de equivalência $R$
em $X$, tal que $p_1:R\to X$ seja próprio e plano.

O mesmo método fornece o resultado seguinte:

\result{TEOREMA}{6.2} Sejam $S$ um pré-esquema Noetheriano, $X$ um
pré-esquema de tipo finito sobre $S$ e $R$ uma pré-relação de equivalência no
$S$-pré-esquema $X$; supõe-se que:
\begin{enumerate}
\item[(a)] $p_1:R_1\to X$ é plano e de tipo finito;
\item[(b)] o morfismo $R_1\to X\times_S X$ é quase-finito (isto é, tem fibras
finitas).
\end{enumerate}
Então existe um aberto $U$ denso em $X$, \emph{saturado} para $R$, tal que:
\begin{enumerate}
\item[(i)] Se $R_U$ é a pré-relação de equivalência induzida por $R$ em $U$,
então $U/R_U$ existe e é de tipo finito sobre $S$.

\item[(ii)] O morfismo canônico $U\to U/R_U$ é sobrejetivo e aberto; suas
fibras são as classes de equivalência conjuntistas para $R_U$ (portanto,
$U/R_U$ é o espaço topológico quociente de $U$ pela relação de equivalência
conjuntista definida por $R_U$); por fim, o morfismo
\[
  (R_U)_1\longrightarrow U\times_{U/R_U}U
\]
é sobrejetivo.
\src{114}

\item[(iii)] Se $R$ provém de uma relação de equivalência, pode-se supor que
$U\to U/R_U$ seja fielmente plano e que $R_U$ seja efetiva.
\end{enumerate}

Este é um resultado de natureza essencialmente «birracional».

\medskip
\noindent\textbf{OBSERVAÇÕES 6.3.}

\smallskip
\noindent $1^\circ$ Ignoro se, em (6.1) e (6.2), a hipótese (b) é inútil. Na
passagem ao quociente por grupos, ela nos obriga praticamente a nos
restringirmos ao caso em que os estabilizadores são todos grupos finitos.

\smallskip
\noindent $2^\circ$ Pode-se perguntar se não há resultados análogos a (6.1) e
(6.2) sem hipótese de planicidade. Não conheço nenhum contraexemplo nessa
direção. Em contrapartida, mesmo conservando a hipótese de planicidade e
restringindo-nos a relações de equivalência tais que $p_1:R\to X$ seja plano e
quase-finito (mas não finito), sendo $X$ afim, pode ocorrer que $R$ não seja
efetiva: tomem-se as relações de equivalência induzidas nos abertos afins que
recobrem a variedade de Nagata (na qual um grupo de ordem dois opera de modo
«não admissível»).

\fsection{7}{Aplicações}

Como dissemos na introdução, a aplicação mais importante de (6.1) é a
construção de esquemas de Picard, bem como de soluções para vários outros
problemas de moduli, aos quais voltaremos posteriormente.

Obtém-se uma demonstração simples do seguinte resultado de SHIMURA:

\result{PROPOSIÇÃO}{7.1} Seja $A$ um esquema abeliano definido sobre um anel de
valorização discreta $V$, com corpo de frações $K$. Então todo esquema
abeliano $B'$ sobre $K$, isógeno a um quociente de $A\otimes_V K$, «tem boa
redução sobre $V$», isto é, é isomorfo a algum $B\otimes_V K$,
\footnote{O texto impresso traz $B\otimes_V A$. Como $B$ é um $V$-esquema e a
afirmação diz respeito à sua fibra genérica, o fator de base deve ser $K$.} onde $B$ é um esquema
abeliano sobre $V$ (essencialmente único, recordemo-lo).

Pode-se supor que $B'$ seja o quociente de $A_K$ por um subesquema em grupos
$C'$. (N.B.: $C'$ não será, em geral, «reduzido», isto é, seus anéis locais
terão elementos nilpotentes.) Considere-se o subesquema fechado $C$ de $A$ que
é o «fecho» de $C'$, isto é, o menor subesquema fechado de $A$ tal que $C_K$
contenha $C'$. Teremos $C_K=C'$, e do fato de que $V$ é um anel de valorização
discreta deduz-se facilmente que $C$ é um subesquema em grupos de $A$ sobre
$V$. Como $A$ é próprio sobre $\operatorname{Spec}(V)=S$, o mesmo vale para
$C$; por outro lado, $A$ é até projetivo sobre $S$. Pode-se então aplicar
(6.1) para construir $A/C=B$,
\src{115}
que é o $B$ procurado.

Por fim, argumentos essencialmente conhecidos permitem deduzir de (6.2) o
seguinte resultado:

\result{TEOREMA}{7.2} Sejam $S$ o espectro de um anel Artiniano, $F$ e $G$
esquemas em grupos de tipo finito sobre $S$, e $u:F\to G$ um homomorfismo de
esquemas em grupos sobre $S$. Suponhamos:
\begin{enumerate}
\item[(i)] $F$ é plano sobre $S$;
\item[(ii)] o núcleo de $u$ é finito.
\end{enumerate}
Nessas condições, existe o esquema quociente $G/F$; o morfismo canônico
$G\to G/F$ é sobrejetivo e aberto, e suas fibras são as classes de equivalência
conjuntistas definidas pelas ações à direita de $F$ sobre $G$. Por fim, se $u$
é um monomorfismo, então $G\to G/F$ é plano e o morfismo
\[
  G\times F\longrightarrow G\times_{G/F}G
\]
é um isomorfismo; em outras palavras, $G$ é um espaço homogêneo principal
sobre $G/F$, com grupo estrutural (atuando à direita) $F$, ou melhor,
$F\times_S(G/F)$ considerado como esquema em grupos sobre $G/F$ (cf. [1], B,
\S~6).

\result{COROLÁRIO}{7.3} Nessas condições, para que $G$ seja plano sobre $S$,
é necessário e suficiente que $G/F$ o seja. Se essa condição for satisfeita,
a passagem ao quociente por $F$ comuta com toda extensão da base $S$ e, se
$F$ for um subgrupo normal de $G$, $G/F$ poderá ser munido de uma estrutura
de \emph{grupo quociente} de $G$ por $F$.

A situação é particularmente simples quando $S$ é o espectro de um corpo,
pois então todo $S$-pré-esquema é automaticamente plano sobre $S$. Obtém-se:

\result{COROLÁRIO}{7.4} Sejam $F,G$ esquemas em grupos de tipo finito sobre um
corpo $k$, e seja $u:F\to G$ um homomorfismo de $k$-grupos. Então $u$ se
fatora como $F\to F'\to G$, onde $F\to F'$ é o homomorfismo de passagem ao
quociente pelo subgrupo fechado $\operatorname{Ker}u$ de $F$, e $F'\to G$ é
um homomorfismo de grupos que é uma imersão fechada. Existe o quociente
$G/F=G/F'$. O formalismo usual (do tipo dos teoremas de Noether) é válido
entre os grupos algébricos sobre $k$.

Este resultado permite tratar de maneira uniforme a passagem ao quociente nos
grupos algébricos no sentido clássico (isto é, os que são irredutíveis sobre
$k$ e suaves sobre $k$) e a passagem ao quociente por subgrupos
«infinitesimais»
\src{116}
considerados por CARTIER. É vantajoso considerar as «hiperálgebras»
introduzidas por CARTIER, na esteira dos trabalhos de DIEUDONNÉ sobre os grupos
formais, como grupos na categoria dos esquemas formais sobre $k$ e, quando for
o caso (se corresponderem a hiperálgebras de posto finito sobre $k$), como
grupos algébricos finitos sobre $k$.

\fsection{8}{Uma conjectura}

Ela responde à necessidade de saber passar ao quociente pelo grupo projetivo
que atua em certos subesquemas dos «esquemas de Hilbert» (estes últimos
substituindo, na teoria dos esquemas, as variedades de Chow).

Sejam $S$ um pré-esquema e $n$ um inteiro. A todo pré-esquema $S'$ sobre $S$,
façamos corresponder o grupo
$\operatorname{Gl}(n,\Gamma(S',\mathcal O_{S'}))$ das matrizes invertíveis de
$n$ linhas e $n$ colunas com entradas no anel das seções de
$\mathcal O_{S'}$. Obtém-se assim um funtor contravariante em $S'$, que se
mostra facilmente representável; ele corresponde, portanto, a um esquema em
grupos sobre $S$, de fato afim sobre $S$, denotado por $\operatorname{Gl}(n)_S$.
Sua formação é compatível com a mudança de base, de modo que, na realidade,
tudo provém de um esquema em grupos sobre $\mathbf Z$, denotado por
$\operatorname{Gl}(n)$. O grupo $\operatorname{Gl}(1)$, chamado grupo
multiplicativo e frequentemente denotado por $\mathbf G_m$, corresponde ao
funtor $S\longmapsto\Gamma(S,\mathcal O_S)^*$, grupo das «unidades» sobre $S$.
Há um homomorfismo evidente
$\operatorname{Gl}(1)\to\operatorname{Gl}(n)$, e constrói-se facilmente o
grupo quociente, denotado por $\operatorname{GP}(n-1)$ e chamado
\emph{grupo projetivo de grau $n-1$ sobre $\mathbf Z$}. Ele representa o
funtor que associa a $S$ o grupo das seções do feixe
$\widetilde{\operatorname{Gl}(n)}_S/
\widetilde{\operatorname{Gl}(1)}_S$, onde
$\widetilde{\operatorname{Gl}(n)}_S$ denota o feixe dos germes de seções de
$\operatorname{Gl}(n)_S$ sobre $S$. (É preciso observar que as seções de
$\operatorname{GP}(n-1)_S$ sobre $S$ não provêm, em geral, de seções de
$\operatorname{Gl}(n)_S$ sobre $S$!) Note-se que se pode provar que
$\operatorname{GP}(n-1)$ representa também o funtor
$S\longmapsto\operatorname{Aut}_S(\mathbf P_S^{n-1})$, onde
$\mathbf P_S^{n-1}$ é o espaço projetivo padrão de dimensão relativa $n-1$
sobre $S$, ao menos quando $S$ é Noetheriano. É dessa maneira que ele intervém
na teoria de módulos.

Seja $S$ um esquema Noetheriano, que se pode, se desejado, supor afim, e seja
$X$ um $S$-pré-esquema quase-projetivo, munido de um feixe invertível
$\mathcal L$ muito amplo relativamente a $S$. Suponhamos que o grupo
$G=\operatorname{GP}(n)_S$ atue em $X$ e, ao mesmo tempo, em $\mathcal L$ (de
modo compatível com sua ação em $X$), e que atue \emph{livremente} em
$X$.\footnote{O texto impresso diz «em $S$». A ação considerada e a condição
de quociente dizem respeito a $X$.}
\src{117}

\result{CONJECTURA}{8.1}\footnote{A errata vinculada à exposição 212 assinala
que esta conjectura é «decididamente falsa». Indica também que o fato positivo
da observação 8.1 foi desde então demonstrado e remete aos trabalhos de Mumford
e a seus refinamentos.} Nas condições precedentes:

\smallskip
\noindent $1^\circ$ \emph{A relação de equivalência definida por $G$ é
efetiva, o quociente $Y=X/G$ é de tipo finito sobre $S$ e o morfismo canônico
$f:X\to Y$ é plano e sobrejetivo} (portanto, $X$ torna-se um fibrado principal
homogêneo sobre $Y$, de grupo
$G\times_S Y=\operatorname{GP}(n)_Y$).

\smallskip
\noindent $2^\circ$ Seja $\mathcal L'$ o feixe invertível sobre $Y$ deduzido
de $\mathcal L$ por «descida fielmente plana» ao longo de $f$ (cf. [1], B,
\S~th. 1). Então $\mathcal L'$ é \emph{pré-amplo sobre $Y$ relativamente a
$S$}, isto é, existem um inteiro $m$ e um morfismo quase-finito de $Y$ para um
esquema projetivo padrão conveniente $\mathbf P_S^N$, tais que
$(\mathcal L')^{\otimes m}$ seja isomorfo à imagem inversa de
$\mathcal O_{\mathbf P_S^N}(1)$.

Observe-se que, mesmo se $X$ é separado sobre $S$, pode ocorrer que $Y$ não
seja separado sobre $S$ (situação que ocorre em problemas de módulos nada
patológicos). Se $1^\circ$ se verifica, então $Y$ é separado se e somente se a
relação de equivalência definida por $G$ tem gráfico fechado, isto é, se
$G\times X\to X\times X$ tem imagem fechada (trata-se então de uma imersão
fechada). Se $Y$ é separado, então $\mathcal L'$ é pré-amplo sobre $Y$
relativamente a $S$ se e somente se é amplo, isto é, se uma potência tensorial
adequada define uma imersão projetiva. Nos problemas de módulos mencionados na
introdução, pode-se mostrar que a relação de equivalência obtida tem de fato um
gráfico fechado.

\medskip
\noindent\textbf{OBSERVAÇÕES 8.1. ---} Supusemos
$G=\operatorname{GP}(n)_S$ para fixar as ideias e porque este é atualmente o
caso mais importante na prática. A hipótese razoável a fazer sobre $G$
pareceria antes que $G$ fosse uma das «formas» sobre $S$ de um dos grupos de
Tohokû (cuja construção sobre os inteiros também foi realizada por CHEVALLEY).
O único fato positivo que conheço na direção da conjectura precedente é o
seguinte: seja $X$ \emph{um esquema afim sobre um corpo $k$ de característica
$0$, no qual o grupo $\operatorname{Gl}(n)_k$ ou
$\operatorname{GP}(n-1)_k$ atua livremente. Então a relação de equivalência
definida por $G$ é efetiva, o quociente $X/G$ também é afim e o morfismo
$X\to X/G$ é plano e sobrejetivo}. A demonstração utiliza o fato seguinte (que
por enquanto só está demonstrado em característica $0$): quando se faz $G$
atuar sobre o anel de coordenadas afim $A$ de $G$, considerado como espaço
vetorial sobre $k$, a representação trivial de $G$ ocorre apenas uma vez (em
uma série de composição de um subespaço vetorial de dimensão finita sobre $k$,
estável sob $G$). Parece possível que um uso sistemático da
\src{118}
a teoria das representações lineares de $G$ acabará por fornecer uma
demonstração da conjectura, ao menos quando se está sobre um corpo de base.
Quando já não se está sobre um corpo de base, o conferencista nada sabe.

\section*{BIBLIOGRAFIA}
\addcontentsline{toc}{section}{Bibliografia}

\begin{description}
\item[{[1]}] GROTHENDIECK (Alexander). --- Techniques de descente et théorèmes
d'existence en géométrie algébrique, I: Généralités, Séminaire Bourbaki,
1959/60, n\textsuperscript{o} 190, 29 p.

\item[{[2]}] GROTHENDIECK (Alexander). --- Techniques de descente et théorèmes
d'existence en géométrie algébrique, II: Le théorème d'existence en théorie
formelle des modules, Séminaire Bourbaki, 1959/60,
n\textsuperscript{o} 195, 22 p.

\item[{[3]}] GROTHENDIECK (A.) e DIEUDONNÉ (J.). --- Éléments de géométrie
algébrique, I: Le langage des schémas. --- Paris, Presses universitaires de
France, 1960 (Institut des hautes Études scientifiques, Publications
mathématiques, 4).
\end{description}

\par

\def\FgaUnitKey{e212}\clearpage
\chapter*{Errata da exposição n.º 212}
\addcontentsline{toc}{chapter}{Errata da exposição n.º 212}
\setcounter{footnote}{2}
\markboth{Errata da exposição 212}{Errata da exposição 212}
\sid{e212}
\src{301}

\noindent\textbf{5. TDTE III: Pré-esquemas quocientes.}

\begin{center}
[Seminário Bourbaki, t. 13, 1960/61, n.º 212, 20 p.]
\end{center}

\noindent\textbf{Página 15, asserção (i).} --- O fato de que $Y=X/R$ seja
quase-projetivo sobre $S$ parece, por ora, ter sido demonstrado apenas quando
$R$ provém de uma relação de equivalência.

\noindent\textbf{Páginas 18 e 19.} --- Como assinalamos ao final da exposição
seguinte, a conjectura considerada é decididamente falsa. O «fato positivo»
assinalado na Observação 8.1 parece ter sido demonstrado simultaneamente por
diversos autores (NAGATA, ROSENLICHT, GROTHENDIECK, \ldots).

Assinalemos que a aplicação feita em TDTE V da teoria da passagem ao quociente
aqui desenvolvida à construção dos esquemas de Picard também pode ser
substituída por uma utilização adequada dos esquemas de Hilbert (cf. SHMT).
Como foi dito no § 8, a lacuna mais importante da teoria apresentada
\src{302}
aqui é a ausência de um critério de existência de quocientes por uma relação de
equivalência não própria, tais como as relações de equivalência provenientes de
certas ações do grupo projetivo. Um teorema importante nessa direção foi obtido
por MUMFORD\footnote{MUMFORD (David). --- An elementary theorem in geometric
invariant theory, \emph{Bull. Amer. Math. Soc.}, t. 67, 1961, pp. 483--486.}.
Para um refinamento de seu resultado e diversas aplicações à teoria, ver SHMT.

\par

\def\FgaUnitKey{221}\unit{Exposição 221 -- Os esquemas de Hilbert}
\sid{221}
\markboth{Exposição 221}{Exposição 221}
\src{249}
\seminares{13.º ano, 1960/61, n.º 221\\Maio de 1961}
\paperhead{TÉCNICAS DE CONSTRUÇÃO E TEOREMAS DE EXISTÊNCIA EM GEOMETRIA ALGÉBRICA\\
IV. OS ESQUEMAS DE HILBERT}%
  {por Alexander GROTHENDIECK}

\subsection*{Introdução}
\addcontentsline{toc}{subsection}{Introdução}

As técnicas expostas em [2], I e II, eram essencialmente independentes de
qualquer hipótese de projetividade sobre os esquemas considerados.
Infelizmente, elas não permitem atualmente resolver os problemas de existência
formulados na Exposição II de [2]. Na presente exposição e na seguinte,
resolveremos esses problemas sob hipóteses de projetividade. As técnicas
empregadas são tipicamente projetivas e praticamente não utilizam os
resultados de I e II de [2]. Aqui construiremos os «esquemas de Hilbert»,
destinados a substituir o uso das coordenadas de Chow, como foi dito em [2],
II, n.º 2. Na exposição seguinte, a teoria da passagem ao quociente em
esquemas desenvolvida em III, juntamente com a teoria dos esquemas de Hilbert,
nos permitirá, por exemplo, construir os esquemas de Picard (definidos em [2],
II, n.º 3) sob condições bastante gerais.

Em resumo, pode-se dizer que agora se dispõe de uma técnica de construções
projetivas bastante satisfatória, salvo que ainda falta\footnote{Ver o adendo
ao fim da exposição.} um teorema de passagem ao quociente por grupos como o
grupo projetivo, atuando «sem pontos fixos» (cf. [2], III, n.º 8). A situação
parece até um pouco melhor em geometria analítica (quando nos limitamos ao
estudo dos espaços analíticos «projetivos» sobre um espaço analítico dado),
pois, para os espaços analíticos, desaparece a dificuldade da passagem ao
quociente por um grupo que atua de maneira adequada. Além disso, tanto em
geometria algébrica como em geometria analítica, restaria desenvolver uma
técnica de construção válida sem hipóteses de projetividade.

\fsection{1}{Famílias limitadas de feixes: propriedades de permanência}

Sejam $k$ um corpo e $X$ um $k$-pré-esquema, de tipo finito, para simplificar.
Para toda extensão $K/k$, obtém-se, portanto, um $K$-pré-esquema
$X_K=X\otimes_k K$. Se $\mathcal F$ é um feixe coerente em $X_K$ e $K'$ é
uma extensão de $K$, então $\mathcal F\otimes_K K'=\mathcal F_{K'}$ é um
feixe quase-coerente em $X_K\otimes_K K'=X_{K'}$. Posto isso, se $K$ e $K'$
são duas extensões quaisquer de $k$, $\mathcal F$ é um feixe quase-coerente em
$X_K$ e $\mathcal F'$
\src{250}
um feixe quase-coerente sobre $X_{K'}$, diremos que $\mathcal F$ e
$\mathcal F'$ são \emph{equivalentes} se existem $k$-homomorfismos de $K$,
$K'$ em uma mesma extensão $K''$ de $k$, tais que $\mathcal F_{K''}$ e
$\mathcal F'_{K''}$ sejam isomorfos sobre $X_{K''}$. Trata-se de uma relação de
equivalência, e interessar-nos-emos pelas classes de equivalência de feixes
para essa relação e pelos conjuntos de classes de equivalência. Observe-se que,
se $X$ é de tipo finito sobre $k$,\footnote{O texto impresso traz aqui $X_0$;
somente $X$ está definido nesta passagem.} então toda classe de feixes
coerentes pode ser definida por um feixe coerente sobre algum $X_K$, onde $K$
é uma extensão de tipo finito de $k$. Pode-se então, na definição das classes
de feixes coerentes, restringir-se às extensões algebricamente fechadas de $k$,
e pode-se também restringir a uma extensão algebricamente fechada fixa $\Omega$
de $k$, de grau de transcendência infinito, sendo dois feixes coerentes
$\mathcal F$ e $\mathcal F'$ sobre $X_\Omega$ equivalentes se e somente se
existe um $k$-automorfismo $\sigma$ de $\Omega$ tal que
$\mathcal F\otimes_\Omega(\Omega,\sigma)$ seja isomorfo a $\mathcal F'$.
Verifica-se que há uma correspondência biunívoca entre as classes de feixes
coerentes para uma ou outra definição.

Sejam $E$ e $E'$ dois conjuntos de classes de feixes coerentes sobre $X$.
Consideremos as classes de todos os feixes da forma
$\mathcal F\otimes\mathcal F'$, onde $\mathcal F$ e $\mathcal F'$ são feixes
coerentes sobre um mesmo $X_K$, com a classe de $\mathcal F$ em $E$ e a de
$\mathcal F'$ em $E'$. Obtém-se assim um conjunto de classes de feixes
coerentes, que será denotado por $E\otimes E'$. Definem-se do mesmo modo os
$\operatorname{Tor}_i(E,E')$, etc. De modo geral, a toda função $U$ que associa
a toda sequência $\mathcal F_1,\ldots,\mathcal F_n$ de $n$ feixes coerentes
sobre um mesmo $X_K$ um conjunto $U(\mathcal F_1,\ldots,\mathcal F_n)$ de
feixes coerentes sobre $X_K$, e que tem uma propriedade evidente de
compatibilidade com os isomorfismos de feixes e as imagens inversas por mudança
de base, associa-se uma função, denotada pelo mesmo símbolo $U$, que associa a
toda sequência de $n$ conjuntos $E_1,\ldots,E_n$ de classes de feixes
coerentes um conjunto $U(E_1,\ldots,E_n)$ de classes de feixes coerentes.

Nosso objetivo neste item é definir certos conjuntos de classes de feixes,
qualificados de \emph{limitados}, e mostrar que as operações $U$ mais usuais,
aplicadas a conjuntos limitados, ainda produzem conjuntos limitados.

Seja $X$ um pré-esquema de tipo finito sobre $S$ Noetheriano. Para todo
$s\in S$, a fibra $X_s$ é um pré-esquema de tipo finito sobre $k(s)$, e
consideraremos classes de feixes coerentes sobre $X_s$ no sentido precedente.
Isso dá sentido à expressão «classe de feixes coerentes sobre uma fibra de
$X/S$» e às expressões análogas. Do mesmo modo, procedendo separadamente sobre
cada fibra, podem-se ainda considerar operações tais como $E\otimes E'$, etc.,
que fazem corresponder a sistemas de conjuntos de
\src{251}
classes de feixes coerentes sobre as fibras de $X/S$, outro conjunto de classes
de feixes coerentes sobre as fibras de $X/S$.

\result{DEFINIÇÃO}{1.1} Seja $E$ um conjunto de classes de feixes coerentes
sobre as fibras de $X/S$. Diz-se que $E$ é \emph{limitado} se existem um
pré-esquema $S'$ de tipo finito sobre $S$ e um feixe coerente $\mathcal F'$ sobre
$X'=X\times_S S'$ tais que $E$ esteja contido no conjunto de classes de feixes
sobre as fibras de $X/S$ definido por $\mathcal F'$.

Por definição, este último faz corresponder a cada $s\in S$ as classes dos
feixes $\mathcal F'\otimes_{S'}k(s')$, em que $s'$ percorre os pontos de $S'$
situados sobre $s$ (de modo que $k(s')$ é uma extensão de $k(s)$ e
$X_s\otimes_{k(s)}k(s')$ identifica-se com a fibra
$X'\otimes_{S'}k(s')=X'_{s'}$ de $X'$ em $s'$). Pode-se dizer que as famílias
limitadas são aquelas contidas em uma \emph{família algébrica} de feixes
coerentes parametrizada por um $S'$ de tipo finito sobre $S$.

Uma união finita de famílias limitadas é limitada (tome-se como pré-esquema a
soma disjunta dos pré-esquemas de parâmetros $S_i$ que definem as famílias
algébricas majorantes). Uma mudança de base $T\longrightarrow S$ transforma uma
família limitada relativa a $X/S$ em uma família limitada relativa a $X_T/T$, e
a recíproca é verdadeira se $T\longrightarrow S$ é sobrejetivo (ou, mais
geralmente, se sua imagem contém os $s$ que intervêm efetivamente na família
dada $E$ para $X/S$). Isso reduz, em princípio, a determinação das famílias
limitadas ao caso em que $S$ é o espectro de uma álgebra de tipo finito sobre o
anel dos inteiros $\mathbf Z$.

Se $E$ e $E'$ são famílias limitadas de classes de feixes relativas a $X/S$,
então $E\otimes E'$ também é limitada. Com efeito, se $E$ e $E'$ são majoradas,
respectivamente, pelas famílias algébricas definidas por $T\longrightarrow S$ e
$\mathcal F$ sobre $X_T$, e por $T'\longrightarrow S$ e $\mathcal F'$ sobre
$X_{T'}$, vê-se que $E\otimes E'$ é majorada pela família algébrica definida por
$T''=T\times_S T'\longrightarrow S$ e pelo feixe $\mathcal F''$ sobre
$X_{T''}$, produto tensorial das imagens inversas de $\mathcal F$ e
$\mathcal F'$ sobre $X_{T''}$. Esse raciocínio só é válido porque o funtor
$\mathcal F\otimes\mathcal F'$ nas variáveis $\mathcal F$, $\mathcal F'$ é exato
à direita e, portanto, comuta com a extensão da base (e, em particular, com a
passagem às fibras). Ele não se aplica, tal como está, às operações locais como
$\operatorname{Tor}_i(E,E')$, $\operatorname{Hom}(E,E')$ e
$\operatorname{Ext}^i(E,E')$. Pode-se, contudo, mostrar que essas operações
continuam transformando conjuntos limitados em conjuntos limitados, procedendo
como para $E\otimes E'$, mas utilizando também resultados do tipo seguinte
(todos contidos em [3], IV, 6.11): uma família limitada $E$ é sempre majorada
por uma família algébrica definida por um feixe coerente $\mathcal F$ sobre um
$X_T$ ($T$ de tipo finito sobre $S$) que é plano relativamente a $T$. (O espaço
de parâmetros inicial é «cortado em partes».) Tais
\src{252}
propriedades de planicidade em feixes adequados asseguram, com efeito, a
comutação de operações como $\operatorname{Tor}_i(\mathcal F,\mathcal F')$
com uma mudança de base arbitrária. O mesmo método aplica-se a operações de
natureza global: imagens diretas e imagens diretas derivadas de feixes
coerentes por morfismos próprios, $\operatorname{Ext}$ globais relativos a
morfismos próprios (cf. [5], III, \S~6), etc.; todas essas operações
transformam famílias limitadas de feixes em famílias limitadas de feixes.
(N. B. -- Aqui os pré-esquemas sobre os quais se tomam os diversos feixes
podem mudar pelas operações do tipo considerado.)

Os dois enunciados seguintes demonstram-se essencialmente pela mesma técnica
de planicidade; para a decomposição primária nas fibras de um morfismo de tipo
finito, veja-se em particular [5], IV.

\result{PROPOSIÇÃO}{1.2} Sejam $E$ e $E'$ conjuntos limitados de classes de
feixes nas fibras de $X/S$, supondo-se $X$ próprio sobre $S$. Então:

\medskip
\noindent (i) A família dos núcleos, conúcleos e imagens de homomorfismos
$\mathcal F\longrightarrow\mathcal F'$, em que a classe de $\mathcal F$
pertence a $E$ e a de $\mathcal F'$ pertence a $E'$, é limitada.

\smallskip
\noindent (ii) A família dos feixes $\mathcal F''$ que são extensões de um
$\mathcal F$ por um $\mathcal F'$, em que a classe de $\mathcal F$ pertence
a $E$ e a de $\mathcal F'$ pertence a $E'$, é limitada.

Mediante uma mudança de base adequada, pode-se supor que $E$ e $E'$ sejam
definidos, respectivamente, por um feixe coerente $\mathcal G$ e um feixe
coerente $\mathcal G'$ sobre um $X_T/T$, sendo $T$ de tipo finito sobre $S$.
Além disso, podem-se supor satisfeitas certas hipóteses de planicidade, que
implicam que a formação dos feixes
$f_{T*}(\operatorname{Hom}_{X_T}(\mathcal G,\mathcal G'))$ e
$\operatorname{Ext}^{1}_{f_T}(\mathcal G,\mathcal G')$ comuta com a mudança
de base por um morfismo arbitrário $T'\longrightarrow T$. Pode-se também
supor que os feixes coerentes precedentes sobre $T$ sejam localmente livres.
Sejam então $T_0$ e $T_1$ os fibrados vetoriais sobre $T$ cujos feixes de
germes de seções são, respectivamente, os feixes precedentes. Define-se então
canonicamente um homomorfismo
$\mathcal G_{T_0}\longrightarrow\mathcal G'_{T_0}$ de feixes coerentes sobre
$X_{T_0}$ e uma extensão
\[
0\longrightarrow\mathcal G'_{T_1}\longrightarrow\mathcal G''
\longrightarrow\mathcal G_{T_1}\longrightarrow0
\]
de feixes coerentes sobre $X_{T_1}$, dotada de uma propriedade universal
evidente. Esse segundo feixe define uma família algébrica que majora a família
considerada em (ii). O mesmo vale para o núcleo, o conúcleo e a imagem do
homomorfismo
\src{253}
precedente, cuja consideração estabelece (i) (desde que se suponha o conúcleo
plano relativamente a $T_0$, caso ao qual ainda se reduz decompondo $T_0$ em
partes\ldots).

\result{PROPOSIÇÃO}{1.3} Seja $E$ uma família limitada de classes de feixes
nas fibras de $X/S$. Então as classes dos feixes estruturais dos
$(\operatorname{supp}\mathcal F)_{\mathrm{red}}$, em que $\mathcal F$ é um
feixe coerente sobre um $X_K$, com $K$ algebricamente fechado, cuja classe
pertence a $E$, formam uma família limitada.

Aqui, $(\operatorname{supp}\mathcal F)_{\mathrm{red}}$ designa o suporte de
$\mathcal F$, munido da estrutura reduzida induzida, isto é, seu feixe
estrutural é o quociente de $\mathcal O_{X_K}$ pelo maior feixe de ideais que
define $\operatorname{supp}\mathcal F$. Pode-se demonstrar o resultado
análogo a (1.3) para os feixes deduzidos canonicamente de $\mathcal F$ pela
teoria da decomposição primária; por exemplo, os
$\mathcal F/\mathcal F_i$, em que os $\mathcal F_i$ são os subfeixes de
$\mathcal F$ primários para as componentes do suporte de $\mathcal F$ e
mínimos com essa propriedade; ou os $\mathcal O_{X_K}/\mathfrak p$, em que
$\mathfrak p$ é um feixe de ideais primos associado a $\mathcal F$; ou os
$\mathcal O_{X_K}/\mathfrak q$, em que $\mathfrak q$ é um feixe de ideais
primários associado a uma componente do suporte de $\mathcal F$ (sendo o
corpo de referência algebricamente fechado).

\fsection{2}{Famílias limitadas e polinômio de Hilbert}

No que segue, supomos que $X$ seja projetivo sobre $S$ e esteja munido de um
feixe muito amplo, denotado por $\mathcal O_X(1)$. Para toda extensão $K$ do
corpo residual $k(s)$ de um ponto $s$ de $S$, considera-se sobre $X_K$ o feixe
correspondente $\mathcal O_{X_K}(1)$, que será ainda muito amplo.

A todo feixe coerente $\mathcal F$ sobre $X_K$ associa-se a função
\[
P_{\mathcal F}(n)=\text{característica de Euler--Poincaré de $\mathcal F(n)$
sobre $X_K$},
\]
que é um polinômio no inteiro $n$, chamado \emph{polinômio de Hilbert} de
$\mathcal F$. Para valores grandes de $n$, $P(n)$ nada mais é do que a
dimensão de $H^0(X_K,\mathcal F(n))$ sobre $K$, pois os
$H^i(X_K,\mathcal F(n))$ são nulos para $i>0$ e $n$ grande.

Se agora $\mathcal F$ for um feixe coerente sobre $X$, plano relativamente a
$S$, então os polinômios de Hilbert dos feixes $\mathcal F_s$ induzidos nas
fibras $X_s$ relativas a uma mesma componente conexa de $S$ serão todos iguais
[5], III, \S~7. Daí resulta (sem hipótese de planicidade) que o conjunto dos
polinômios de Hilbert dos feixes $\mathcal F_s$, $s\in S$, é finito para
todo feixe coerente $\mathcal F$ sobre $X$.
\src{254}
Recordemos, por outro lado, que, se $\mathcal F$ é um feixe coerente sobre $X$,
este é isomorfo a um feixe quociente de um feixe $\mathcal O_X(-n)^N$, para
$n$ e $N$ suficientemente grandes. Portanto, os feixes $\mathcal F_s$
induzidos nas fibras são igualmente quocientes do feixe $\mathcal O(-n)^N$ na
fibra.

Das duas observações, deduz-se a parte «somente se» do teorema seguinte:

\result{TEOREMA}{2.1} Seja $X$ projetivo sobre um $S$ Noetheriano, e seja
$\mathcal O_X(1)$ muito amplo sobre $X$ relativamente a $S$. Seja $E$ um
conjunto de classes de feixes nas fibras de $X/S$. Para que $E$ seja limitado,
é necessário e suficiente\footnote{O texto impresso traz apenas «é necessário»;
a errata de 1962 da Exposição 221 prescreve «é necessário e suficiente».} que
satisfaça as condições seguintes:

\medskip
\noindent (a) Existe um feixe coerente $\mathcal L$ sobre $X$ (que se pode
supor da forma $\mathcal O_X(-n)^N$) tal que $E$ esteja contido na família das
classes de feixes coerentes quocientes de feixes da forma $\mathcal L_K$.

\smallskip
\noindent (b) Os polinômios de Hilbert $P_{\mathcal F}$ dos feixes
$\mathcal F$ cuja classe pertence a $E$ são elementos de um mesmo conjunto
finito de polinômios.

Resta provar a parte «se», que será um caso particular de um resultado mais
preciso. Para todo Módulo coerente $\mathcal F$ sobre um pré-esquema de tipo
finito sobre um corpo $K$, e todo inteiro $r$, seja $N_r$ o sub-Módulo de
$\mathcal F$ cujas seções sobre um aberto são as seções de $\mathcal F$ sobre
esse aberto cujo suporte tem dimensão $<r$. Tem-se, portanto,
$N_r=\mathcal F$ para $r>\dim\operatorname{supp}\mathcal F$, $N_r=0$ para
$r\leq0$, e obtém-se uma filtração crescente finita de $\mathcal F$ cujos
fatores $N_{r+1}/N_r$\footnote{O texto impresso inverte o quociente, escrevendo
$N_r/N_{r+1}$, o que é incompatível com a filtração crescente que acaba de ser
definida.} têm como ciclos primos associados exatamente os ciclos primos
associados a $\mathcal F$ que têm dimensão $r$. Poremos
\[
\mathcal F_{(r)}=\mathcal F/N_r.
\]
Assim, $\mathcal F_{(r)}$ tem como ciclos primos associados exatamente os
ciclos primos associados a $\mathcal F$ que têm dimensão $\geq r$ e, em
particular, é igual a $\mathcal F$ se e somente se os ciclos primos associados
a $\mathcal F$ têm dimensão $\geq r$. Posto isso:

\result{TEOREMA}{2.2} Nas condições de (2.1), seja $s$ um inteiro e suponhamos
que $E$ satisfaça a condição (a) e a seguinte forma enfraquecida de (b):

\medskip
\noindent $(b_s)$ Os polinômios de Poincaré $P_{\mathcal F}$ dos feixes
$\mathcal F$ cuja classe pertence a $E$ têm coeficientes em graus
$\geq s-1$ que permanecem limitados.

Nessas condições, os feixes $\mathcal F_{(s)}$ (quando a classe de
$\mathcal F$ percorre $E$) formam uma família limitada. Além disso, os
coeficientes de grau $s-2$ dos $P_{\mathcal F}$ permanecem limitados
inferiormente.
\src{255}
Assim:

\result{COROLÁRIO}{2.3} Suponhamos que os feixes $\mathcal F$ cuja classe
pertence a $E$ sejam tais que todos os ciclos primos associados tenham dimensão
$d$ com $s\leq d\leq r$. Então, em (2.1) (b), pode-se considerar apenas os
coeficientes dos $P_{\mathcal F}$ em graus compreendidos entre $s-1$ e $r$.

O fim deste item é dedicado ao esboço da demonstração de (2.2). Os lemas-chave
são os dois lemas seguintes, o primeiro dos quais é bem conhecido (e resume o
conteúdo matemático útil das coordenadas de Chow):

\result{LEMA 2.4 (CHOW)}{} Consideremos os feixes estruturais de subesquemas
$Y$ de fibras $X_K$ ($K$ percorrendo extensões algebricamente fechadas de
corpos residuais de $S$), com $Y$ reduzido e todas as suas componentes de mesma
dimensão $r$ (considerando-se $\mathcal O_Y$ como um feixe quociente de
$\mathcal O_X$). Se os graus dos $Y$ permanecem limitados, os $Y$ formam uma
família limitada.

Aqui, o grau $a$ de $Y$ pode ser definido mais convenientemente pelo
coeficiente do termo dominante de
\[
P_{\mathcal O_Y}=a\frac{n^r}{r!}+\cdots.
\]

\result{LEMA}{2.5} Sejam $\mathcal L$ um feixe coerente sobre $X$ e $E$ um
conjunto de classes de feixes quocientes $\mathcal F$ de feixes $\mathcal L_K$
($K$ uma extensão de corpo residual de $S$). Suponhamos que as fibras de $X$
sobre $S$ tenham dimensão $\leq r$, e ponhamos
\[
P_{\mathcal F}(n)=a_{\mathcal F}\frac{n^r}{r!}
+b_{\mathcal F}\frac{n^{r-1}}{(r-1)!}
+\text{termos de grau $<r-1$}.
\]
Então o coeficiente $a_{\mathcal F}$ permanece limitado, e $b_{\mathcal F}$
permanece limitado inferiormente. Se $b_{\mathcal F}$ permanece limitado,
então a família dos $\mathcal F_{(r)}$ é limitada.

\emph{Demonstração.} -- Substituindo, se necessário, $S$ por uma união de
subesquemas de $S$ que recobre $S$, pode-se supor que existe um morfismo finito
$f:X\longrightarrow\mathbf P^r_S$, tal que $\mathcal O_X(1)$ seja isomorfo à
imagem inversa de $\mathcal O_{\mathbf P^r_S}(1)$; portanto, para todo feixe
coerente $\mathcal F$ sobre $X$, tem-se
$P_{\mathcal F}=P_{f_*(\mathcal F)}$. Por outro lado, verifica-se facilmente
(pela técnica do item precedente) que um conjunto de feixes $\mathcal F$ sobre
$X$ é limitado se e somente se o conjunto dos $f_*(\mathcal F)$ o é. Por fim,
\[
f_*(\mathcal F)_{(r)}=f_*(\mathcal F_{(r)}).
\]
\src{256}
Isso nos reduz, portanto, ao caso $X=\mathbf P^r_S$. Além disso, pode-se supor
que $\mathcal L=\mathcal O_{\mathbf P^r_S}(k)^s$ para $k$, $s$ convenientes. O
coeficiente $a_{\mathcal F}$ satisfaz
\[
0\leq a_{\mathcal F}\leq s
\]
e, portanto, permanece limitado. Dito isso, é equivalente afirmar que os
$P_{\mathcal F}(n)$ têm um coeficiente de $n^{r-1}$ que permanece limitado
inferiormente (resp. limitado), ou afirmar o mesmo para
$P_{\mathcal F}(n-k)=P_{\mathcal F(-k)}(n)$. Isso nos reduz ao caso em que
\[
\mathcal L=\mathcal O_{\mathbf P^r_S}^{s}.
\]

Consideremos a sequência exata
\[
0\longrightarrow N_r\longrightarrow\mathcal F
\longrightarrow\mathcal F_{(r)}\longrightarrow0,
\]
da qual se obtém
\[
P_{\mathcal F}=P_{\mathcal F_{(r)}}+P_{N_r},
\]
e, como o coeficiente de $n^{r-1}$ em $P_{N_r}$ é não negativo (pois
$\dim\operatorname{supp}N_r\leq r-1$), tem-se
\[
b_{\mathcal F_{(r)}}\leq b_{\mathcal F}.
\]
Isso nos permite, para demonstrar as afirmações do lema, substituir
$\mathcal F$ por $\mathcal F_{(r)}$, isto é, supor que os feixes quocientes
$\mathcal F$ de $\mathcal L$ considerados são livres de torção.

Como $\mathbf P^r_K$ é normal, segue-se que $\mathcal F$ é localmente livre de
posto $a=a_{\mathcal F}$ em um aberto $U=\mathbf P^r_K-Y$, onde $Y$ tem
codimensão $\geq2$. Logo, $\bigwedge^a\mathcal F$ é um feixe sobre
$\mathbf P^r_K$ cuja restrição a $U$ é invertível; portanto (como
$\mathbf P^r_K$ é regular e $Y$ tem codimensão $\geq2$), essa restrição é
isomorfa à restrição de um feixe invertível sobre $\mathbf P^r_K$, determinado
salvo isomorfismo. Este último é da forma $\mathcal O_{\mathbf P^r_K}(d)$ para um
inteiro $d$ bem determinado. Como $\bigwedge^a\mathcal F$ é um quociente de
$\bigwedge^a\mathcal O_{\mathbf P^r_K}^{s}\simeq
\mathcal O_{\mathbf P^r_K}^{N}$\footnote{O texto impresso substitui aqui o
expoente $s$ por um $n$ não definido; $s$ é o posto do feixe livre do qual
$\mathcal F$ é quociente.}, com $N=\binom{s}{a}$, possui $N$ seções canônicas,
que definem, portanto, seções de
\src{257}
$\mathcal O_{\mathbf P^r_K}(d)$ sobre $U$, que são restrições de seções
$s_i$ ($1\leq i\leq N$) de $\mathcal O_{\mathbf P^r_K}(d)$ (pois
$\mathbf P^r_K$ é normal e $Y$ tem codimensão $\geq2$). Estas seções
geram $\mathcal O_{\mathbf P^r_K}(d)$ nos pontos de $U$, portanto não são
todas nulas, o que implica que $d\geq0$. Por outro lado, um cálculo simples
mostra que
\[
b_{\mathcal F}=a_{\mathcal F}(r+1)/2+d.
\]
Isso mostra, em particular, que $b_{\mathcal F}>0$, portanto
$b_{\mathcal F}$ é limitado inferiormente. Esse coeficiente permanece limitado
se e somente se $d$ permanece limitado; mostremos que, nesse caso, $\mathcal F$
pertence a uma família limitada. Podemos supor que $a_{\mathcal F}$ e
$b_{\mathcal F}$ sejam fixos, digamos iguais a $a$ e $b$; portanto, $d$ é fixo.
O conhecimento das $N$ seções
$s_i$ de $\mathcal O_{\mathbf P^r_K}(d)$, isto é, de um homomorfismo
\[
s:\bigwedge^a\mathcal L_K\longrightarrow\mathcal O_{\mathbf P^r_K}(d),
\]
permite recuperar $\mathcal F$ como a coimagem do homomorfismo composto
correspondente:
\[
\mathcal L_K\longrightarrow
\operatorname{Hom}\!\left(\bigwedge^{a-1}\mathcal L_K,
\bigwedge^a\mathcal L_K\right)\longrightarrow
\operatorname{Hom}\!\left(\bigwedge^{a-1}\mathcal L_K,
\mathcal O_{\mathbf P^r_K}(d)\right),
\]
em que o primeiro é o homomorfismo canônico proveniente do produto exterior, e
o segundo é deduzido de $s$. A conclusão decorre então de (1.2) (i).

Combinando os dois lemas precedentes, obtém-se:

\result{LEMA}{2.6} Sob as hipóteses preliminares de (2.1), suponhamos que, para
todo $\mathcal F$, se tenha
\[
P_{\mathcal F}(n)=a_{\mathcal F}\frac{n^r}{r!}
+b_{\mathcal F}\frac{n^{r-1}}{(r-1)!}
+\text{termos de grau $<r-1$},
\]
e que os coeficientes $a_{\mathcal F}$ permaneçam limitados. Então os
coeficientes $b_{\mathcal F}$ permanecem limitados inferiormente. Se os
$b_{\mathcal F}$ forem limitados, então os $\mathcal F_{(r)}$ formam uma família
limitada.

Podemos supor que os corpos de base $K$ dos feixes $\mathcal F$ sejam
algebricamente fechados. Munamos cada
$\operatorname{supp}\mathcal F_{(r)}$, união das componentes de dimensão
$r$\footnote{O texto impresso traz “grau $r$”; a filtração
$\mathcal F_{(r)}$ é definida pela dimensão do suporte, não pelo grau.} de
$\operatorname{supp}\mathcal F$, da estrutura reduzida induzida. Então os graus dos
$\operatorname{supp}\mathcal F_{(r)}$ permanecem limitados superiormente por
$a$; portanto, em virtude de (2.4), os
$\operatorname{supp}\mathcal F_{(r)}$ formam um conjunto limitado. Além disso,
para toda componente de $\operatorname{supp}\mathcal F_{(r)}$, o comprimento de
$\mathcal F_{(r)}$ nessa componente é $\leq a$; portanto, se
$\mathcal J_{\mathcal F}$ é o Ideal que define
$\operatorname{supp}\mathcal F_{(r)}$, então $\mathcal F_{(r)}$
\src{258}
pode ser considerado um Módulo sobre o subesquema $Y_{\mathcal F}$ de $X$
definido por $\mathcal J_{\mathcal F}^{a}$. Como no lema precedente, reduzimo-nos também
ao caso em que $\mathcal F=\mathcal F_{(r)}$, de modo que $\mathcal F$
provém de um Módulo sobre $Y_{\mathcal F}$. Os $Y_{\mathcal F}$ correspondem
a uma família limitada de Módulos quocientes dos $\mathcal O_{X_K}$ e,
portanto, provêm de um subesquema fechado $Y$ de um esquema $X\times_S T$.
Pode-se então aplicar (2.5) a $Y/T$ e $\mathcal L\otimes_XY$, o que dá a
conclusão.

Podemos agora demonstrar (2.2) por indução sobre o limite superior $r$ dos
$\dim\operatorname{supp}\mathcal F$. O enunciado é trivial para $r<0$;
suponhamos, pois, que $r\geq0$ e que o enunciado tenha sido demonstrado para
$r'<r$. Em virtude de (2.6), os $\mathcal F_{(r)}$ formam um conjunto
limitado; portanto, também em virtude de (1.2) (i), os núcleos dos
homomorfismos $\mathcal L_K\longrightarrow\mathcal F_{(r)}$ formam um
conjunto limitado. Existe, pois, um Módulo coerente $\mathcal L'$ sobre $X$
tal que os núcleos em questão e, consequentemente, também os
\[
N_r(\mathcal F)=\operatorname{Ker}(\mathcal F\longrightarrow\mathcal F_{(r)}),
\]
sejam quocientes dos Módulos $\mathcal L'_K$. Como os $\mathcal F_{(r)}$ são
limitados, os $P_{\mathcal F_{(r)}}$ permanecem limitados, e a fórmula
\[
P_{\mathcal F}=P_{\mathcal F_{(r)}}+P_{N_r}
\]
mostra então que os $P_{N_r}$ satisfazem a mesma condição $(b_s)$ que os
$P_{\mathcal F}$. Logo, os $N_r$ satisfazem as condições (a) e $(b_s)$ e,
pela hipótese de indução, os $(N_r)_{(s)}$ permanecem limitados. Ora,
$\mathcal F_{(s)}$ é uma extensão de $\mathcal F_{(r)}$ por
$(N_r)_{(s)}$; portanto, por (1.2) (ii), os $\mathcal F_{(s)}$ permanecem
limitados. Para a última afirmação de (2.2), conduzimos simultaneamente a indução sobre a limitação e sobre a cota inferior do coeficiente seguinte. Os núcleos limitados dos morfismos de fonte fixa para $\mathcal F_{(s)}$ fornecem um único feixe coerente que se sobrejeta sobre cada $N_s$, mas (1.2) (i) por si só não torna limitada a família desses quocientes. Reaplicamos indutivamente a afirmação sobre os coeficientes a esses quocientes; a aditividade do polinômio de Hilbert e o lema 2.6 dão então a cota inferior requerida para o coeficiente seguinte. Isso conclui a demonstração.

\fsection{3}{Esquemas de Hilbert: definição, teorema de existência}

Sejam $X$ um pré-esquema sobre outro $S$, e $\mathcal F$ um Módulo
quase-coerente sobre $X$. Denotemos por
\[
\operatorname{Quot}(\mathcal F/X/S)
\]
o conjunto dos Módulos quase-coerentes quocientes de $\mathcal F$ que são
planos sobre $S$. Seja agora $S'\longrightarrow S$ um morfismo de mudança de
base; ponhamos $X'=X\times_S S'$ e
$\mathcal F'=\mathcal F\otimes_{\mathcal O_S}\mathcal O_{S'}$; assim, $X'$
é um pré-esquema sobre $S'$ munido de um Módulo
\src{259}
quase-coerente $\mathcal F'$, e pode-se considerar
$\operatorname{Quot}(\mathcal F'/X'/S')$. Ponhamos
\[
\underline{\operatorname{Quot}}_{\mathcal F/X/S}(S')
=\operatorname{Quot}(\mathcal F'/X'/S')
\qquad\text{(onde $X'=X\times_SS'$)}.
\]
Se agora $S''\longrightarrow S'$ for um $S$-morfismo, então
$X''=X\times_SS''$ é isomorfo a $X'\times_{S'}S''$, e $\mathcal F''$ é
isomorfo a $\mathcal F'\otimes_{\mathcal O_{S'}}\mathcal O_{S''}$. Como o
funtor imagem inversa
\[
\mathcal G'\longmapsto
\mathcal G'\otimes_{\mathcal O_{S'}}\mathcal O_{S''}
\]
da categoria dos Módulos quase-coerentes sobre $X'$ para a categoria dos
Módulos quase-coerentes sobre $X''$ é exato à direita e transforma Módulos
$S'$-planos em Módulos $S''$-planos, obtém-se uma aplicação natural
\[
\operatorname{Quot}(\mathcal F'/X'/S')\longrightarrow
\operatorname{Quot}(\mathcal F''/X''/S''),
\]
de modo que $\underline{\operatorname{Quot}}_{\mathcal F/X/S}(S')$ é um
funtor contravariante em $S'$ (pré-esquema sobre $S$), com valores na categoria
dos conjuntos. Doravante, suporemos que $X$ seja projetivo sobre um $S$
noetheriano, que $\mathcal F$ seja coerente e, para simplificar, limitar-nos-emos
aos $S'$ sobre $S$ que sejam localmente noetherianos.

\result{TEOREMA}{3.1} Nessas condições, o funtor contravariante
$\underline{\operatorname{Quot}}_{\mathcal F/X/S}$ sobre a categoria dos
$S$-pré-esquemas localmente noetherianos é representável por um
$S$-pré-esquema $\operatorname{Quot}_{\mathcal F/X/S}$, união disjunta de uma
sequência de $S$-esquemas projetivos (a fortiori,
$\operatorname{Quot}_{\mathcal F/X/S}$ é localmente de tipo finito sobre
$S$).

Obteremos tal decomposição da seguinte maneira. Seja $\mathcal O_X(1)$ um
feixe invertível sobre $X$, muito amplo relativamente a $S$. Para todo
polinômio $P(n)$ com coeficientes racionais, seja
$\operatorname{Quot}^P(\mathcal F/X/S)$ a parte de
$\operatorname{Quot}(\mathcal F/X/S)$ formada pelos quocientes coerentes
$\mathcal G$ de $\mathcal F$ que são planos sobre $S$ e cujo polinômio de
Hilbert em todo $s\in S$ é igual a $P$. Poremos então
\[
\underline{\operatorname{Quot}}^P_{\mathcal F/X/S}(S')
=\operatorname{Quot}^P(\mathcal F'/X'/S')
\]
e obteremos assim um subfuntor de
$\underline{\operatorname{Quot}}_{\mathcal F/X/S}$. A propriedade de
invariância dos polinômios de Hilbert recordada no número 2 implica o seguinte:
para que $\underline{\operatorname{Quot}}_{\mathcal F/X/S}$ seja
representável, é necessário e suficiente que os
$\underline{\operatorname{Quot}}^P_{\mathcal F/X/S}$ o sejam; nesse caso, o
$S$-pré-esquema $\operatorname{Quot}_{\mathcal F/X/S}$ que o representa é
isomorfo ao pré-esquema união disjunta dos
$\operatorname{Quot}^P_{\mathcal F/X/S}$ que representam os funtores
$\underline{\operatorname{Quot}}^P_{\mathcal F/X/S}$. Posto isso, (3.1)
será, portanto, consequência do teorema seguinte:
\src{260}
\result{TEOREMA}{3.2} Com as notações anteriores, o funtor
$\underline{\operatorname{Quot}}^P_{\mathcal F/X/S}$ é representável por um
$S$-pré-esquema projetivo $\operatorname{Quot}^P_{\mathcal F/X/S}$.

O restante deste número é dedicado à demonstração de (3.2).

Seja $\nu$ um inteiro. Para todo $S'$ sobre $S$, denotamos por $A_\nu(S')$ o
conjunto dos quocientes $\mathcal G=\mathcal F'/\mathcal K$ de
$\mathcal F'=\mathcal F\otimes_{\mathcal O_S}\mathcal O_{S'}$, que são
coerentes, planos sobre $S'$ e satisfazem as condições seguintes:
\begin{align*}
\text{a.}\quad &R^if'_*(\mathcal G(n))=0 &&\text{para $i>0$ e $n\geq\nu$},\\
\text{b.}\quad &R^if'_*(\mathcal K(n))=0 &&\text{para $i>0$ e $n\geq\nu$},\\
\text{c.}\quad &f'_*(\mathcal K(\nu+k))
=\mathcal S'_k f'_*(\mathcal K(\nu)) &&\text{para $k\geq0$}.
\end{align*}
Para esta última relação, supõe-se que $X$ esteja escrito como o espectro primo
homogêneo de uma álgebra graduada quase-coerente $\mathcal S_*$ sobre $S$,
graduada em graus positivos e gerada por $\mathcal S_1$, e põe-se
$\mathcal S'_* =\mathcal S_*\otimes_{\mathcal O_S}\mathcal O_{S'}$, de modo
que $X'$ seja o espectro primo homogêneo de $\mathcal S'_*$. Para demonstrar
(3.2), reduz-se facilmente ao caso em que $X=\mathbf P^r_S$ (pois $S$ é união
de abertos $U$ tais que $X|U$ seja um subesquema fechado de $\mathbf P^r_U$ e
$\mathcal O_X(1)$ seja induzido por $\mathcal O_{\mathbf P^r_U}(1)$), e ao
caso em que $\mathcal F$ é da forma
$\mathcal O_{\mathbf P^r_S}^{N}$, sendo, portanto, plano sobre $S$. Então, no
que precede, os feixes $\mathcal K$ também são planos sobre $S'$. Resulta das
relações de Künneth ([5], III, \S~7) que as condições (a) e (b) são estáveis
por mudança de base e implicam que, para $n\geq\nu$, a formação de
$f'_*(\mathcal G(n))$ e de $f'_*(\mathcal K(n))$ comute com a mudança de base
(loc. cit.); portanto, (c) também é estável por mudança de base.

Em outras palavras, $A_\nu(S')$ é um funtor contravariante em $S'$, mais
precisamente um subfuntor de
$A(S')=\underline{\operatorname{Quot}}^P_{\mathcal F/X/S}(S')$. Ao variar
$\nu$, obtém-se assim uma sequência crescente de subconjuntos $A_\nu(S')$ de
$A(S')$, cuja união é $A(S')$ em virtude de um teorema bem conhecido de
Serre [5], III, \S~2. Observemos agora que, se $\mathcal G$ é um quociente
coerente de $\mathcal F'$, plano sobre $S'$, e $s$ é um ponto de $S'$ tal que
a mudança de base $\operatorname{Spec}(k(s))\longrightarrow S'$ dá origem a
um quociente $\mathcal G_s$ de $\mathcal F_s$ que satisfaz as condições (a),
(b) e (c), isto é, que pertence a $A_\nu(\operatorname{Spec}(k(s)))$, então
existe uma vizinhança aberta $U$ de $s$ tal que essas mesmas condições sejam
verificadas para $\mathcal G|f'^{-1}(U)$, isto é, esse quociente pertence a
$A_\nu(U)$; para (a) e (b), isso resulta, com efeito, do «teorema das funções
holomorfas» [5], III, \S~7, e (c) resulta do lema de Nakayama e do fato de que,
em todo caso, sabe-se que
$f'_*(\mathcal K(n+k))=\mathcal S'_k f'_*(\mathcal K(n))$
\src{261}
para $n$ suficientemente grande e $k\geq0$ (cf. [5], III, \S~2). Dessas
observações conclui-se o seguinte (compare-se [4], IV). Para que o funtor $A$
seja representável, é necessário e suficiente que os funtores $A_\nu$ o sejam;
nesse caso, o $S$-pré-esquema $Q$ que representa $A$ é uma união crescente de
subpré-esquemas abertos $Q_\nu$ que representam os $A_\nu$.

Seja
\[
M_* =\bigoplus_{n\geq0}f_*(\mathcal F(n))=\mathcal S_*^{N},
\]
de modo que
\[
M'_* =M_*\otimes_{\mathcal O_S}\mathcal O_{S'}
=\bigoplus_{n\geq0}f'_*(\mathcal F'(n))=\mathcal S'_*{}^{N}.
\]
De (a) resulta que:

\noindent (a') $f'_*(\mathcal G(n))$ é localmente livre de posto $P(n)$ para
$n\geq\nu$, [5], III, \S~7, e, de (b), para $i=1$, que:

\noindent (a'') $f'_*(\mathcal G(n))$ é um Módulo quociente de $M'_n$.

Além disso, o conhecimento desse Módulo quociente, para $n=\nu$, determina,
em virtude de (c), os sub-Módulos $f'_*(\mathcal K(n))$ de $M'_n$ para
$n\geq\nu$ e, portanto, determina $\mathcal K$ e, por conseguinte, $\mathcal G$.
Obtém-se assim uma aplicação injetiva
\[
A_\nu(S')\longrightarrow\operatorname{Grass}_{P(\nu)}(M'_\nu)
\]
de $A_\nu(S')$ no conjunto dos Módulos quocientes localmente livres de posto
$P(\nu)$ de $M'_\nu$ e, portanto, um homomorfismo funtorial
\[
i_\nu:A_\nu(S')\longrightarrow
\underline{\operatorname{Grass}}_{P(\nu)}(M_\nu)(S'),
\]
onde o funtor do segundo membro é representado pelo esquema de Grassmann
$\operatorname{Grass}_{P(\nu)}(M_\nu)$ (compare-se [4], V), que é projetivo
sobre $S$. Afirmo:

\result{LEMA}{3.3} $A_\nu(S')$ é um funtor representável, e o morfismo
$Q_\nu\longrightarrow\operatorname{Grass}_{P(\nu)}(M_\nu)$ que representa o
homomorfismo $i_\nu$ é uma imersão (o que implica que $Q_\nu$ é
quase-projetivo sobre $S$).

Essa afirmação é equivalente à seguinte (compare-se [4], IV): suponha-se dado
um Módulo quociente $N$ de $M'_\nu$, localmente livre de posto $P(\nu)$; então
existe um subpré-esquema $Z$ de $S'$ tal que, para todo pré-esquema localmente
noetheriano $T'$
\src{262}
sobre $S'$, a imagem inversa de $N$ sobre $T'$ pertence a
$\operatorname{Im}A_\nu(T')$ se e somente se $T'\longrightarrow S'$ se
fatoriza pelo subpré-esquema $Z$. Mudando as notações, pode-se supor $S'=S$,
isto é, tem-se um quociente $N_\nu$ de $M_\nu$ por um sub-Módulo $R_\nu$.
Para que esse quociente provenha de um elemento de $A_\nu(S)$, é necessário e
suficiente que satisfaça as duas condições seguintes:

\medskip
\noindent (i) $M_{\nu+k}/\mathcal S_kR_\nu$ é localmente livre de posto
$P(\nu+k)$ para $k\geq0$.

\smallskip
\noindent (ii) O subfeixe $\mathcal K$ de $\mathcal F$ definido pelo
sub-Módulo graduado
\[
R_* =\sum_{k\geq0}\mathcal S_kR_\nu
\]
de $M_*$ ([5], II, \S~3) e o quociente $\mathcal G=\mathcal F/\mathcal K$
satisfazem as condições (a) e (b) acima.

Essas condições são manifestamente necessárias; por outro lado, se forem
satisfeitas, então o feixe $\mathcal G$ definido em (ii), uma vez que é
isomorfo ao feixe associado ao $\mathcal S_*$-Módulo graduado $N_*$ que é a
soma dos $M_{\nu+k}/\mathcal S_kR_\nu$, é plano sobre $S$ (pois suas fibras
são somandos diretos de localizações de $N_*$ em ideais primos homogêneos de
$\mathcal S_*$), e tem polinômio de Hilbert $P$ em virtude de (i). Levando em
conta (ii), vê-se então que (a') e (a'') são satisfeitas, logo, para
$n\geq\nu$, o homomorfismo natural $N_n\longrightarrow f_*(\mathcal G(n))$ é
um homomorfismo sobrejetivo entre Módulos localmente livres do mesmo posto e,
portanto, um isomorfismo; assim, tem-se
\[
f_*(\mathcal K(\nu+k))=\mathcal S_kR_\nu
\quad\text{para todo $k\geq0$},
\]
o que prova que $\mathcal G\in A_\nu(S)$ e que $R_\nu$ é o elemento de
$\operatorname{Grass}_{P(\nu)}(M_\nu)$ definido por $\mathcal G$.

Esse critério (i)--(ii) aplica-se igualmente à situação obtida após uma
mudança de base $S'\longrightarrow S$. Provaremos primeiro que o fato de a
condição (i) ser satisfeita após a mudança de base $S'\longrightarrow S$ se
exprime pela condição de que $S'\longrightarrow S$ se fatorize por certo
subpré-esquema $Z$ de $S$; obtido esse resultado, ficamos reduzidos
(substituindo $S$ por $Z$) ao caso em que a condição (i) já é satisfeita sobre
$S$ e, como ela é estável por mudança de base, resta exprimir a condição (ii).
Mas então, se $U$ denota o conjunto dos $s\in S$ tais que a cohomologia dos
feixes induzidos na fibra $X_s$ por $\mathcal G(n)$ e $\mathcal K(n)$ seja nula
em dimensão $>0$ para $n\geq\nu$, já se observou que $U$ é aberto, e a condição
(ii) será satisfeita após uma mudança de base $S'\longrightarrow S$ se e somente
se $S'\longrightarrow S$ se fatoriza por $U$, o que prova o Lema 3.3. Resta,
portanto, provar o lema seguinte:

\result{LEMA}{3.4} Sejam $S$ um pré-esquema localmente noetheriano, munido de
uma álgebra graduada quase-coerente $\mathcal S_*$ em graus positivos, gerada
por $\mathcal S_1$, $M_*$ um $\mathcal S_*$-Módulo graduado quase-coerente de tipo
finito, $P$ um polinômio com coeficientes racionais e $\nu$ um inteiro. Então
existe um subpré-esquema $Z$ de $S$ (evidentemente único)
\src{263}
com a seguinte propriedade: para todo pré-esquema $S'$ sobre $S$, para que
$M_n\otimes_{\mathcal O_S}\mathcal O_{S'}$ seja localmente livre de posto
$P(n)$ para todo $n\geq\nu$, é necessário e suficiente que
$S'\longrightarrow S$ se fatorize por $Z$.

Pode-se, evidentemente, supor $S$ afim, logo noetheriano. Então:

\result{LEMA}{3.5} Para todo inteiro $N\geq\nu$, seja $U_N$ o aberto de $S$
formado pelos $s\in S$ tais que
\[
\operatorname{rank}_{k(s)}
M_{ns}\otimes_{\mathcal O_{S,s}}k(s)\leq P(n)
\quad\text{para todo $\nu\leq n\leq N$}.
\]
Então a sequência decrescente dos abertos $U_N$ é estacionária.

Sabe-se [3], IV, que $S$ admite uma partição finita em subpré-esquemas
reduzidos $S_i$, tais que cada
$M\otimes_{\mathcal O_S}\mathcal O_{S_i}$ seja plano sobre $S_i$. Pode-se,
portanto, supor $M$ plano, logo os $M_n$ planos. Enfim, pode-se evidentemente
supor $S$ conexo. Mas então existem um inteiro $n_0$ e um polinômio $Q$ tais
que
\[
\operatorname{rank}_{k(s)}
M_{ns}\otimes_{\mathcal O_{S,s}}k(s)=Q(n)
\quad\text{para $n\geq n_0$},
\]
[5], III, \S~7. Suponhamos primeiro $P<Q$\footnote{O texto impresso traz
$P\neq Q$; a errata de 1962 prescreve $P<Q$, substitui o caso «$P=Q$» pelo
caso contrário $P(n)\geq Q(n)$ para $n$ grande e substitui as desigualdades
estritas de (*) por desigualdades não estritas.}; então $P(n)<Q(n)$ para $n$
grande e, evidentemente, $U_N=\varnothing$ para $N$ grande; a fortiori, a
sequência dos $U_N$ é estacionária. No caso contrário, teremos
$P(n)\geq Q(n)$ para $n$ grande, portanto $U_N=U_{n_0}$ para $N\geq n_0$, e
a sequência dos $U_N$ é ainda estacionária.

Em particular, o conjunto $U_\infty$ dos $s\in S$ tais que
\[
\tag{*}
\operatorname{rank}_{k(s)}
M_{ns}\otimes_{\mathcal O_{S,s}}k(s)\leq P(n)
\quad\text{para todo $n\geq\nu$},
\]
sendo a interseção dos $U_N$, é aberto. Pode-se então, para a demonstração de
(3.4), substituir $S$ pelo aberto $U_\infty$, o que nos reduz ao caso em que a
desigualdade (*) é satisfeita em todo $s\in S$. Além disso:

\result{LEMA}{3.6} Seja $M$ um Módulo sobre um pré-esquema localmente
noetheriano $S$, e seja $r$ um inteiro. Então existe um subpré-esquema $Z$ de
$S$ (evidentemente único) com a seguinte propriedade: para todo $S'$ sobre
$S$, para que $M\otimes_{\mathcal O_S}\mathcal O_{S'}$ seja localmente
livre de posto $r$, é necessário e suficiente que $S'\longrightarrow S$ se
fatorize por $Z$. Se
\[
\operatorname{rank}_{k(s)}
M_s\otimes_{\mathcal O_{S,s}}k(s)\leq r
\quad\text{para todo $s$},
\]
então $Z$ é um subpré-esquema fechado de $S$ (supondo-se $M$ coerente).
\src{264}
Com efeito, o raciocínio precedente nos reduz ao caso em que a desigualdade
(*) vale para todo $s\in S$ (substituindo, se necessário, $S$ pela parte
aberta formada pelos $s$ em que a desigualdade é verificada). Pode-se então
supor que $M$ se insira em uma sequência exata
\[
\mathcal O_S^q\longrightarrow\mathcal O_S^r\longrightarrow M\longrightarrow0,
\]
e a condição considerada sobre os $S'$ sobre $S$ significa também que, na
sequência exata correspondente
$\mathcal O_{S'}^q\longrightarrow\mathcal O_{S'}^r
\longrightarrow M'\longrightarrow0$, a segunda seta é um isomorfismo, isto
é, a primeira é nula. Vê-se então que o subpré-esquema fechado $Z$ de $S$
definido pelo Ideal gerado pelos coeficientes da matriz que define o
homomorfismo $\mathcal O_S^q\longrightarrow\mathcal O_S^r$ satisfaz a
condição desejada.

Retomando então a demonstração de (3.4) no ponto em que a deixamos, denote-se
por $Z_n$ o subesquema fechado de $S$ associado, em virtude de (3.6), ao
Módulo $M_n$ e ao inteiro $r=P(n)$, e por $Z'_N$ o
$\operatorname{Inf}$ dos $Z_n$ para $\nu\leq n\leq N$. Então os $Z'_N$
formam uma sequência decrescente de subesquemas fechados de $S$, portanto
necessariamente estacionária. Seja $Z$ o valor constante dos $Z'_N$ para $N$
grande. Esse é o $Z$ procurado em (3.4). Isso conclui a demonstração de (3.4)
e, portanto, também a de (3.3).

Demonstramos, pois, que
$\underline{\operatorname{Quot}}^P_{\mathcal F/X/S}$ é representável por um
$S$-pré-esquema $Q$ que é uma união crescente de subpré-esquemas abertos
$Q_\nu$ quase-projetivos sobre $S$. Para ir mais longe, é preciso invocar o
teorema 2.1, do qual se conclui facilmente que $Q$ é quase-compacto (pois é a
imagem do pré-esquema $S'$ de tipo finito sobre $S$ que parametriza a família
dos feixes quocientes dos $\mathcal F_K$ com polinômio de Hilbert $P$).
Portanto, $Q$ é igual a um dos $Q_\nu$, logo quase-projetivo sobre $S$. Para
demonstrar que é projetivo sobre $S$, resta demonstrar que é próprio sobre
$S$; para isso, basta invocar o critério valuativo de properidade sob a forma
[5], II, 7.3.8. Basta verificar o seguinte:

\result{LEMA}{3.7} Sejam $S$ o espectro de um anel de valorização discreta,
$s$ seu ponto genérico, $X$ um pré-esquema sobre $S$, $\mathcal F$ um Módulo
quase-coerente sobre $X$ e $\mathcal G_s$ um Módulo quociente quase-coerente
de $\mathcal F_s=\mathcal F\otimes_{\mathcal O_S}k(s)$ sobre $X_s$. Então
existe um único Módulo quociente quase-coerente $\mathcal G$ de $\mathcal F$,
plano sobre $S$, cuja restrição a $X_s$ seja $\mathcal G_s$.
\src{265}
Com efeito, se $\mathcal G_s=\mathcal F_s/\mathcal K_s$, basta considerar o
maior subfeixe $\mathcal K$ de $\mathcal F$ que induz $\mathcal K_s$, [5],
I, 9.4.2, e tomar $\mathcal G=\mathcal F/\mathcal K$. Verifica-se facilmente
que esse feixe convém.

O teorema 3.2 e, por conseguinte, (3.1) estão completamente demonstrados.

A demonstração prova, ao mesmo tempo, o seguinte:

\result{PROPOSIÇÃO}{3.8} Nas condições de (3.2), sejam
$Q=\operatorname{Quot}^P_{\mathcal F/X/S}$, $X_Q=X\times_SQ$,
$\mathcal F_Q=\mathcal F\otimes_{\mathcal O_S}\mathcal O_Q$, e seja
$\mathcal G$ o quociente coerente de $\mathcal F_Q$, plano sobre $Q$, com
polinômio de Hilbert relativo $P$, tal que $(Q,\mathcal G)$ represente o
funtor $\underline{\operatorname{Quot}}^P_{\mathcal F/X/S}$. Existe um
inteiro $\nu$ tal que, para $n\geq\nu$,
$(f_Q)_*(\mathcal G(n))$ seja um Módulo sobre $Q$ localmente livre de posto
$P(n)$ e muito amplo relativamente a $S$, isto é, que defina uma imersão de
$Q$ em um esquema grassmanniano $\operatorname{Grass}_{P(n)}(M)$ sobre $S$.
A fortiori, para $n\geq\nu$, o feixe
\[
\bigwedge^{P(n)}(f_Q)_*(\mathcal G(n))
\]
sobre $Q$ é invertível e muito amplo relativamente a $S$.

Pode-se, com efeito, reduzir, como em (3.2), ao caso em que $\mathcal F$ seja
plano sobre $S$; então basta tomar um inteiro $\nu$ tal que $A_\nu=A$ (com as
notações precedentes).

O caso de aplicação mais importante de (3.2) é aquele em que
$\mathcal F=\mathcal O_X$. Escreve-se então
\[
\operatorname{Quot}_{\mathcal O_X/X/S}=\operatorname{Hilb}_{X/S},
\qquad
\operatorname{Quot}^P_{\mathcal O_X/X/S}=\operatorname{Hilb}^P_{X/S},
\]
logo se tem uma decomposição
\[
\operatorname{Hilb}_{X/S}=\coprod_P\operatorname{Hilb}^P_{X/S}.
\]
Por definição, $\operatorname{Hilb}_{X/S}$ representa o funtor
$\underline{\operatorname{Hilb}}_{X/S}(S')$, conjunto dos subpré-esquemas
fechados de $X'=X\times_SS'$ que são planos sobre $S'$, e
$\operatorname{Hilb}^P_{X/S}$ representa o subfuntor correspondente aos
subpré-esquemas fechados que admitem um polinômio de Hilbert dado $P$. Esses
pré-esquemas são também chamados, respectivamente, \emph{pré-esquema de
Hilbert} de $X$ sobre $S$ e \emph{pré-esquema de Hilbert de índice $P$}. A
terminologia é justificada pelo papel desempenhado na teoria pelos polinômios
de Hilbert. Sua diferença de natureza em relação às variedades de Chow
clássicas (destinadas a parametrizar ciclos, e não variedades) é da mesma ordem
que aquela entre
\src{266}
o anel de Chow das classes de ciclos de uma variedade e o anel das classes
de feixes da variedade (tal como aparece no teorema de Riemann--Roch [1]);
com efeito, observe-se que, quando $X=\mathbf P^r_S$, sendo $S$ o espectro
de um corpo, o conhecimento do polinômio de Hilbert de um Módulo coerente
$\mathcal F$ sobre $X$ equivale também ao conhecimento das classes de Chern
de $\mathcal F$, ou ainda ao da classe de $\mathcal F$ no anel das classes de
feixes coerentes sobre $X$.

\result{OBSERVAÇÕES}{3.9} Observe-se também que a construção dos
$\operatorname{Quot}_{\mathcal F/X/S}$ e
$\operatorname{Quot}^P_{\mathcal F/X/S}$ foi reduzida ao caso em que
$X=\mathbf P^r_S$, $\mathcal F=\mathcal O_X^N$, sendo $\mathcal O_X(1)$ o
feixe muito amplo usual; mais precisamente, os
$\operatorname{Quot}_{\mathcal F/X/S}$ gerais realizam-se como
subpré-esquemas fechados dos precedentes. Como a formação dos
$\operatorname{Quot}^P_{\mathcal F/X/S}$ é evidentemente compatível com as
mudanças de base $S'\longrightarrow S$, vê-se que se está reduzido ao caso em
que, além disso, $S=\operatorname{Spec}(\mathbf Z)$. Em última análise, fica-se,
portanto, reduzido a estudar os esquemas projetivos sobre $\mathbf Z$:
\[
\mathfrak Q^P_{r,N}=
\operatorname{Quot}^{P}_{(\mathcal O_{\mathbf P^r_{\mathbf Z}})^N/
\mathbf P^r_{\mathbf Z}/\operatorname{Spec}(\mathbf Z)},
\]
e, mais particularmente, os esquemas de Hilbert absolutos:
\[
\operatorname{Hilb}^P_r=\mathfrak Q^P_{r,1}.
\]
Um estudo mais detalhado desses esquemas, começando pela determinação de suas
componentes conexas (são conexos?),\footnote{A errata de 1962 assinala que
Hartshorne demonstrou a conexidade dos $\operatorname{Hilb}^P_r$ sobre um
corpo algebricamente fechado e determinou os pares $(r,P)$ para os quais
$\operatorname{Hilb}^P_r\neq\varnothing$.} e de suas componentes irredutíveis
(em virtude de Serre [7], pode haver componentes irredutíveis situadas
inteiramente sobre um número primo $p\neq0$), seria muito interessante.
Recordemos a questão de Weil sobre se as componentes irredutíveis das fibras de
$\operatorname{Hilb}^P_r$ sobre os $s\in\operatorname{Spec}(\mathbf Z)$
correspondem a extensões ``regulares'' do corpo primo, isto é, se são
``relativamente conexas''. Não se exclui que essas questões sejam mais
acessíveis para os esquemas de Hilbert do que para as ``variedades de Chow''.

\fsection{4}{Variantes}

\noindent a. Nas condições de (3.1), seja $U$ um aberto de $X$, e designemos
por $A'$ o subfuntor de
$A=\underline{\operatorname{Quot}}_{\mathcal F/X/S}$, tal que $A'(S')$ seja
o conjunto dos Módulos quocientes $\mathcal G$ de $\mathcal F'$, planos sobre
$S'$, cujo suporte está contido em $U'$. Vê-se imediatamente que $A'$ é
representável por uma parte aberta do pré-esquema
$\operatorname{Quot}_{\mathcal F/X/S}$ que representa $A$. Segue-se que os
teoremas (3.1) e (3.2)
\src{267}
continuam válidos supondo-se que $X$ seja quase-projetivo sobre $S$ em vez de
projetivo sobre $S$, quando se substitui também, nas conclusões, a palavra
``projetivo'' por ``quase-projetivo'', e se designa agora por
$\underline{\operatorname{Quot}}_{\mathcal F/X/S}(S')$ o conjunto dos
quocientes coerentes $\mathcal G$ de $\mathcal F'$, planos sobre $S'$, cujo
suporte é próprio sobre $S'$.

\medskip
\noindent b. De modo geral, podem-se impor aos quocientes $\mathcal G$ de
$\mathcal F'=\mathcal F\otimes_{\mathcal O_S}\mathcal O_{S'}$, planos sobre
$S'$, toda espécie de condições naturais adicionais, estáveis por mudança de
base, obtendo-se assim tantos subfuntores de
$\underline{\operatorname{Quot}}_{\mathcal F/X/S}$ quantos se pretende
representar. O critério usual permite, em muitos casos, demonstrar que ainda se
obtêm funtores representáveis por partes abertas de
$\operatorname{Quot}_{\mathcal F/X/S}$. Isso ocorre, em particular, quando se
impõe uma das seguintes propriedades adicionais:

\smallskip
\noindent $1^\circ$ As dimensões dos ciclos primos associados aos Módulos
$\mathcal G_{s'}$ ($s'\in S'$) induzidos nas fibras $X'_{s'}$ pertencem a um
conjunto dado de inteiros.

\smallskip
\noindent $2^\circ$ (Quando $\mathcal F=\mathcal O_X$, portanto $\mathcal G$
corresponde a um subpré-esquema fechado $Y$ de $X'$); $Y$ é um pré-esquema
suave ([3], IV) sobre $S'$, resp. normal sobre $S'$ (isto é, as fibras $Y_{s'}$
são normais ``sobre $k(s')$'', isto é, são normais após toda extensão do corpo
de base), resp. (quando $X$ é plano sobre $S$) é uma $S'$-interseção completa
local em $X'$ relativamente a $S'$ (isto é, as fibras $Y_{s'}$ são interseções
completas locais nos $X'_{s'}$).

Outras condições fariam intervir propriedades de natureza cohomológica nos
Módulos $\mathcal G_{s'}$ induzidos nos $X'_{s'}$, etc. Naturalmente, a
conjunção de condições, cada uma das quais representada por um aberto $U_i$ de
$\operatorname{Quot}_{\mathcal F/X/S}$, é representada pelo aberto que é sua
interseção. Por exemplo, considerando, para todo $S'$ sobre $S$, o conjunto dos
subpré-esquemas fechados $Y$ de $X'=X\times_SS'$ que são recobrimentos étales
([3], I) de posto dado $r$ de $S'$, obtém-se um funtor contravariante
representável em $S'$.

\medskip
\noindent c. Os pré-esquemas $\operatorname{Hom}_S(X,Y)$,
$\Pi_{X/S}(Z/X)$, $\operatorname{Isom}_S(X,Y)$, definidos em [2], II, C,
n\textsuperscript{o} 2, existem sob hipóteses de projetividade convenientes e
realizam-se como partes abertas de pré-esquemas de Hilbert convenientes. Como se
tem
\[
\operatorname{Hom}_S(X,Y)=\Pi_{X/S}((X\times_SY)/X),
\]
o caso de $\operatorname{Hom}_S(X,Y)$ reduz-se ao de $\Pi_{X/S}(Z/X)$.
Observa-se então que, para todo $S'$ sobre $S$, o conjunto das seções de
$Z'=Z\times_SS'$ sobre $X'=X\times_SS'$ está em correspondência biunívoca com
o conjunto dos subpré-esquemas $\Gamma$ de $Z'$ (necessariamente fechados se
$Z$ é separado sobre $X$) tais que o morfismo $\Gamma\longrightarrow X'$
induzido por $Z'\longrightarrow X'$ seja um isomorfismo. Dessa maneira,
\src{268}
quando $X$ é plano e próprio sobre $S$ e $Z$ é quase-projetivo sobre $S$,
$\Pi_{X/S}(Z/X)$ existe e realiza-se como um subpré-esquema aberto de
$\operatorname{Hilb}_{Z/S}$. Portanto, quando $X$ é projetivo e plano sobre
$S$ e $Y$ é quase-projetivo sobre $S$, $\operatorname{Hom}_S(X,Y)$ existe e
realiza-se como um subpré-esquema aberto de
$\operatorname{Hilb}_{(X\times_SY)/S}$. Quando $X$ e $Y$ são ambos projetivos
sobre $S$, segue-se imediatamente que $\operatorname{Isom}_S(X,Y)$ também
existe e é representado por uma parte aberta de
$\operatorname{Hom}_S(X,Y)$. Do mesmo modo, quando $X$ é plano e projetivo
sobre $S$ e $Y$ é quase-projetivo sobre $S$, o $S$-pré-esquema
$\operatorname{Imm}_S(X/Y)$, correspondente ao subfuntor do funtor representado
por $\operatorname{Hom}_S(X,Y)$ que corresponde aos $S'$-homomorfismos
$X'\longrightarrow Y'$ que são imersões, é também representável por uma
parte aberta de $\operatorname{Hom}_S(X,Y)$.

Seja $\mathcal L$ (resp. $\mathcal M$) um feixe invertível sobre $X$ (resp.
$Y$), muito amplo relativamente a $S$, donde um feixe muito amplo
$\mathcal L\otimes_{\mathcal O_S}\mathcal M$ sobre $X\times_SY$ relativamente a
$S$. Por conseguinte, para todo polinômio $P$ com coeficientes racionais,
$\operatorname{Hilb}^P_{(X\times_SY)/S}$ está definido e é um pré-esquema
quase-projetivo sobre $S$. Ele induz, portanto, sobre
$\operatorname{Hom}_S(X,Y)$ uma parte simultaneamente aberta e fechada,
quase-projetiva sobre $S$, que denotaremos por
$\operatorname{Hom}_S(X,Y)^P$. Assim, as seções de
$\operatorname{Hom}_S(X,Y)^P$ sobre $S$ são os $S$-morfismos
$g:X\longrightarrow Y$ tais que, para todo inteiro $n$, se tenha
\[
\chi\!\left(\left(\mathcal L\otimes_{\mathcal O_X}g^*(\mathcal M)\right)^{
\otimes n}\right)=P(n).
\]

Obtêm-se dessa maneira generalizações do teorema de Matsusaka, que afirma
que os automorfismos de uma variedade projetiva «polarizada» formam um grupo
algébrico, afirmação que assume aqui um sentido evidentemente mais preciso,
pois dispomos de uma definição desse grupo como solução de um problema
universal. Note-se, aliás, que, sobre um corpo algebricamente fechado, o grupo
dos automorfismos considerado anteriormente é o que se deduz do «verdadeiro»
definido aqui passando ao quociente pelos elementos nilpotentes; isso explica
por que há pouca chance de que as construções mais antigas possam ser feitas
sobre um corpo de base não perfeito, pois o Ideal dos elementos nilpotentes que
aparece após uma extensão do corpo de base não é necessariamente «definido
sobre $k$». Essa mesma observação aplica-se, aliás, também à maioria das construções
de estilo antigo.

\fsection{5}{Estudo diferencial dos esquemas de Hilbert}

Ele decorre do resultado seguinte:
\src{269}
\result{PROPOSIÇÃO}{5.1} Sejam $S$ um pré-esquema, $S_0$ um subpré-esquema
definido por um Ideal quase-coerente $\mathcal J$ de quadrado nulo, $X$ um
$S$-pré-esquema, $\mathcal F$ um Módulo quase-coerente sobre $X$,
$X_0=X\times_SS_0$ e
$\mathcal F_0=\mathcal F\otimes_{\mathcal O_S}\mathcal O_{S_0}$, e, por fim,
$\mathcal G_0=\mathcal F_0/\mathcal K_0$ um Módulo quociente quase-coerente de
$\mathcal F_0$, plano sobre $S_0$. Para todo aberto $U$ de $X$, seja
$\mathcal E(U)$ o conjunto dos Módulos quocientes quase-coerentes $\mathcal G$ de
$\mathcal F|U$, planos sobre $S$, e tais que
$\mathcal G\otimes_{\mathcal O_S}\mathcal O_{S_0}=\mathcal G_0$; assim, quando
$U$ varia, os $\mathcal E(U)$ são as seções de um feixe $\mathcal E$ sobre
$X$. Posto isso, o feixe de grupos
\[
\mathcal A=\operatorname{Hom}_{\mathcal O_{X_0}}
(\mathcal K_0,\mathcal G_0\otimes_{\mathcal O_{S_0}}\mathcal J)
\]
age naturalmente sobre $\mathcal E$, que se torna assim um feixe
«formalmente homogêneo principal sob $\mathcal A$» (isto é, para todo aberto
$U$ em $X$, $\mathcal E(U)$ é vazio ou um conjunto homogêneo principal sob
$\mathcal A(U)$).

Conclui-se daí:

\result{COROLÁRIO}{5.2} Suponhamos que exista localmente sobre $X$ um
prolongamento $\mathcal G$ de $\mathcal G_0$ a um quociente de $\mathcal F$
plano sobre $S$ (isto é, que as fibras do feixe $\mathcal E$ sejam não vazias).
Então existe uma classe de obstrução canônica
\[
c(\mathcal G_0)\in H^1(X,\mathcal A),
\]
cujo anulamento é necessário e suficiente para a existência de um
prolongamento global $\mathcal G$ de $\mathcal G_0$ a um quociente de
$\mathcal F$ plano sobre $S$. Se essa classe se anula, então o conjunto
$\mathcal E(X)$ de todos os prolongamentos possíveis é um conjunto homogêneo
principal sob
\[
\mathcal A(X)=\operatorname{Hom}_{\mathcal O_{X_0}}
(\mathcal K_0,\mathcal G_0\otimes_{\mathcal O_{S_0}}\mathcal J).
\]
\footnote{Nesta última fórmula, o texto impresso põe $\mathcal O_X$ no lugar
de $\mathcal O_{X_0}$, embora este último esteja fixado na definição
imediatamente precedente de $\mathcal A$.}
A existência do prolongamento global fica então garantida, em particular, se
$H^1(X,\mathcal A)=0$.

\result{COROLÁRIO}{5.3} Suponhamos que
$Q=\operatorname{Quot}_{\mathcal F/X/S}$ exista (cf. 4 (a)) -- por exemplo,
que $X$ seja quase-projetivo sobre um $S$ localmente noetheriano e que
$\mathcal F$ seja coerente. Seja $x\in Q$ correspondente a uma extensão
residual $K=k(x)$ de $k(s)$ ($s\in S$):
\src{270}
portanto, $x$ é definido por um Módulo quociente coerente
$\mathcal G_0=\mathcal F_0/\mathcal K_0$ do Módulo
$\mathcal F_0=\mathcal F_K$ sobre o $K$-pré-esquema $X_K$. Seja $\mathcal A$
o feixe coerente sobre $X_K$ definido por
\[
\mathcal A=\operatorname{Hom}_{\mathcal O_{X_0}}(\mathcal K_0,\mathcal G_0).
\]
Então o espaço tangente de Zariski da fibra $Q_s$ no ponto $x$ (dual sobre
$K$ de $\mathfrak m/\mathfrak m^2$, onde $\mathfrak m$ é o ideal maximal de
$\mathcal O_{Q_K,x}$) é canonicamente isomorfo a $H^0(X_K,\mathcal A)$.

O resultado que fornece o espaço tangente de Zariski pode ser generalizado e
fornece uma caracterização, para um $S$-morfismo dado
$g:S'\longrightarrow Q$, isto é, uma seção $g'$ de $Q'=Q\times_SS'$ sobre
$S'$, do Módulo
\[
\Omega=g^*(\Omega^1_{Q/S})=g'^*(\mathcal J/\mathcal J^2)
\]
(onde $\mathcal J$ é o Ideal sobre $Q'$ definido pela seção $g'$ de $Q'$ sobre
$S'$), pela fórmula funtorial no Módulo coerente $\mathcal M$ sobre $S'$:
\[
\operatorname{Hom}_{\mathcal O_{S'}}(\Omega,\mathcal M)
\simeq H^0(X',\mathcal A),
\]
onde $\mathcal A$ é ainda o Módulo sobre $X'=X\times_SS'$ definido por
\[
\mathcal A=\operatorname{Hom}_{\mathcal O_{X'}}
(\mathcal K,\mathcal G\otimes_{\mathcal O_{S'}}\mathcal M),
\]
\footnote{O texto impresso põe $\mathcal O_S$ nesse produto tensorial; aqui
$\mathcal G$ e $\mathcal M$ são Módulos sobre $S'$, e a construção da
deformação de quadrado nulo exige o produto sobre $\mathcal O_{S'}$.}
($\mathcal G=\mathcal F'/\mathcal K$ sendo o Módulo quociente de
$\mathcal F'=\mathcal F\otimes_{\mathcal O_S}\mathcal O_{S'}$ que corresponde
a $g$). Com efeito, basta aplicar (5.1), substituindo ali $S_0$ por $S'$ e $S$
pelo pré-esquema $D(\mathcal M)=(S',\mathcal O_{S'}+\mathcal M)$, onde
$\mathcal M$ é considerado como um Ideal de quadrado nulo.

Quando, em (5.1), se tem $\mathcal F=\mathcal O_X$, portanto o dado de
$\mathcal G_0$ corresponde ao dado de um subpré-esquema fechado $Y_0$ de
$X_0$, plano sobre $S_0$, definido pelo Ideal
$\mathcal J_0=\mathcal K_0$, então (*) fornece
\[
\mathcal A=\operatorname{Hom}_{\mathcal O_{X_0}}
(\mathcal J_0/\mathcal J_0^2,
\mathcal O_{Y_0}\otimes_{\mathcal O_{S_0}}\mathcal J),
\]
onde $\mathcal J_0/\mathcal J_0^2$ se interpreta como o feixe conormal a $Y_0$ em
$X_0$, que se denota também por
\src{271}
$\check{\mathcal N}_{Y_0/X_0}$; é então conveniente considerar $\mathcal A$
como um Módulo sobre $Y_0$ e calcular $H^0$ e $H^1$ sobre $Y_0$. Ademais, se
$Y_0$ é localmente uma interseção completa em $X_0$, sendo $X$ plano sobre
$S$, então, em (5.1), fica garantida a possibilidade de prolongamento local;
por outro lado, $\mathcal J_0/\mathcal J_0^2$ é localmente livre sobre $Y_0$ e
pode-se escrever
\[
\mathcal A=\mathcal N_{Y_0/X_0}\otimes_{\mathcal O_{S_0}}\mathcal J,
\]
onde o primeiro fator do segundo membro é o feixe normal a $Y_0$ em $X_0$.
Utilizando o critério fundamental de lisura ([3], III, 3.1), obtém-se, por
exemplo:

\result{COROLÁRIO}{5.4} Nas condições de (5.3), suponhamos que
$\mathcal F=\mathcal O_X$, que $X$ seja plano sobre $S$ e que o
subpré-esquema fechado $Y_0$ de $X_0$ que corresponde a $\mathcal G_0$ seja
localmente uma interseção completa. Então o espaço tangente de Zariski a $Q_s$
em $x$ é canonicamente isomorfo a
$H^0(Y_0,\mathcal N_{Y_0/X_0})$. Se
$H^1(Y_0,\mathcal N_{Y_0/X_0})=0$, então o pré-esquema de Hilbert
$\operatorname{Hilb}_{X/S}$ é liso sobre $S$ no ponto $x$.
($\mathcal N_{Y_0/X_0}$ é o feixe normal a $Y_0$ em $X_0$.)

\result{OBSERVAÇÃO}{5.5} Este enunciado aplica-se, em particular, quando $Y_0$
é uma interseção completa em $X_0$ definida por uma equação, isto é, é um
``divisor de Cartier'' positivo. Então $\mathcal N_{Y_0/X_0}$ é isomorfo ao
feixe sobre $Y_0$ induzido pelo feixe invertível $\mathcal J_0^{-1}$ sobre $X_0$
definido pelo divisor $Y_0$. É essa a situação encontrada, em particular, no
estudo das famílias de divisores positivos sobre uma variedade projetiva não
singular $X_0$. O isomorfismo entre o espaço tangente de Zariski no ponto $x$
de $Q$ (ou, se se preferir, do aberto $D$ de $Q$ que corresponde aos divisores)
e $H^0(Y_0,\mathcal N_{Y_0/X_0})$ era conhecido em geometria algébrica
clássica sob o nome de ``homomorfismo característico'' do primeiro no segundo.
Ele só era definido quando $x$ era um ponto liso da variedade de parâmetros
$T$ de uma ``família contínua completa'' de divisores, isto é, de nosso ponto
de vista, de uma componente irredutível do esquema $D$, munida da estrutura
reduzida induzida. O espaço tangente a $T$ em $x$ é então um subespaço do
espaço tangente a $D$ em $x$; portanto, o homomorfismo característico clássico
é de fato injetivo, mas só é sobrejetivo sob condições adicionais, por exemplo,
se $D$ é integral em $x$. De fato, Zappa [8] construiu um exemplo
\src{272}
(com $X$ uma superfície projetiva não singular sobre o corpo dos complexos) no
qual, mesmo no ponto genérico de $T$, o homomorfismo característico não é
sobrejetivo. Isso significa, portanto, que $D$ não é integral nem sequer no
ponto genérico da componente irredutível considerada. Isso mostra de maneira
particularmente impressionante como as variedades com elementos nilpotentes
são necessárias para compreender fenômenos da mais clássica teoria das
superfícies.\footnote{A errata de 1962 acrescenta que Mumford acabava de
construir uma componente irredutível do esquema de Hilbert de
$\mathbf P^3_{\mathbf C}$, cujos pontos gerais representam curvas não
singulares de grau 14 e gênero 24, não reduzida em seu ponto genérico; a
explosão de uma dessas curvas fornece também um esquema projetivo regular de
dimensão 3 cujo esquema formal de módulos é genericamente não reduzido.}

\subsection*{5.6}

Em (5.4), deu-se um critério de lisura, aplicável em particular aos esquemas de
divisores. Kodaira deu em [6] um critério diferente, a saber, o anulamento de
$H^1(X_0,\mathcal L)$, onde $\mathcal L=\mathcal J_0^{-1}$ é o feixe
invertível sobre $X_0$ definido pelo divisor $Y_0$; critério válido quando $S$
é o espectro de um corpo de característica 0 e demonstrado em [6] por via
transcendente no caso em que o corpo de base é $\mathbf C$. Observemos aqui
que, em geral, sendo $S$ novamente arbitrário, a condição de Kodaira é
suficiente para que o morfismo canônico
\[
D\longrightarrow\operatorname{Pic}_{X/S}
\]
do pré-esquema dos divisores para o pré-esquema de Picard de $X/S$ seja liso
no ponto $x$ considerado (como se verifica facilmente pelo critério usual de
lisura, uma vez estabelecida a existência de $\operatorname{Pic}_{X/S}$). Se,
portanto, além disso, $\operatorname{Pic}_{X/S}$ é liso sobre $S$ no ponto
imagem de $x$ (por exemplo, se $\operatorname{Pic}_{X/S}$ é liso sobre $S$),
então $D$ é liso sobre $S$ em $x$. Por outro lado, Cartier demonstrou que todo
pré-esquema em grupos localmente de tipo finito sobre um corpo $k$ de
característica 0 é liso sobre $k$. Combinando esses dois resultados, recupera-se
o resultado de Kodaira. Observe-se que decorre dessas observações que, sobre um
corpo $K$ de característica $p>0$, quando $\operatorname{Pic}_{X/S}$ não é
liso sobre $K$ (o que ocorre quando $X$ é a superfície de Igusa), a condição
$H^1(X_0,\mathcal L)=0$ implica, ao contrário, que $D$ não é liso em $x$, e
nem sequer é reduzido em $x$ se $K$ é algebricamente fechado.

Para terminar, assinalemos ainda o resultado seguinte, que desempenha um papel
importante no estudo diferencial dos espaços fibrados:

\result{PROPOSIÇÃO}{5.7} Sejam $X$ um pré-esquema finito e plano sobre $S$
localmente noetheriano, e $Z$ um pré-esquema sobre $X$, tal que
$\Pi_{X/S}(Z/X)$ exista (o que ocorre se $Z$ é quase projetivo sobre $X$). Se,
então, $Z$ é liso sobre $X$, $\Pi_{X/S}(Z/X)$ é liso sobre $S$.

É uma consequência imediata da definição e do critério usual de lisura ([3],
III, 3.1). Observemos que, quando $X$ é finito e plano sobre $S$, a questão da
existência de $\Pi_{X/S}(Z/X)$\footnote{O texto impresso traz aqui $Z/K$;
nenhum $K$ está definido, e o funtor de seções considerado desde (5.7) é
$\Pi_{X/S}(Z/X)$.} pode ser tratada de maneira muito elementar, sem
\src{273}
utilizar a teoria dos esquemas de Hilbert. Verifica-se, por exemplo, que, se
$X$ é radicial sobre $S$, então $\Pi_{X/S}(Z/X)$ existe sem nenhuma restrição
sobre $Z$. Por exemplo, seja $T$ um $S$-pré-esquema, e seja $T_n$ a
``vizinhança infinitesimal de ordem $n$'' da diagonal de $T\times_ST$ em
$T\times_ST$, munida dos morfismos $p_1,p_2:T_n\longrightarrow T$ induzidos
pelas duas projeções. Considera-se $T_n$ como um pré-esquema finito sobre $T$
por meio de $p_1$; suponhamo-lo, além disso, plano sobre $T$ (o que ocorre se
$T$ é liso sobre $S$). Para todo pré-esquema $X$ sobre $T$, ponhamos
\[
(X/T/S)^{(n)}=
\Pi_{T_n/T}\bigl(p_2^*(X/T)/T_n\bigr),
\]
\footnote{O texto impresso põe $T_n/S$ sob o símbolo $\Pi$. A proposição 5.7
aplica-se aqui ao morfismo finito e plano $p_1:T_n\to T$, e o resultado é
explicitamente um pré-esquema sobre $T$; portanto, o denominador é $T_n/T$.}
este é um pré-esquema sobre $T$ chamado \emph{fibrado dos germes de seções de
ordem $n$ de $X$ sobre $T$} (relativamente a $S$). Ele depende funtorialmente
de $X$ e é liso sobre $T$ se $X$ o é.

\fsection{6}{Relações com a noção de norma e os produtos simétricos}

Sejam $S$ um pré-esquema, $X$ e $Y$ $S$-pré-esquemas,
\[
u:(X/S)^n\longrightarrow Y
\]
um $S$-morfismo simétrico da $n$-ésima potência cartesiana de $X/S$ em $Y$.
Suporemos, para simplificar, $S$ localmente noetheriano, e $X$ e $Y$ de tipo
finito sobre $S$. Pode-se então, a todo Módulo coerente $\mathcal F$ sobre $X$,
com suporte finito sobre $S$, plano sobre $S$ e de posto relativo sobre $S$
igual a $n$ (isto é, tal que $f_*(\mathcal F)$ seja um Módulo localmente livre
de posto $n$ sobre $S$), fazer corresponder de maneira natural uma seção de $Y$
sobre $S$:
\[
\pi^u_{X/S}(\mathcal F)\in\Gamma(Y/S).
\]
Não daremos aqui a definição formal, limitando-nos a assinalar que o formalismo
ao qual se chega é uma generalização natural do formalismo usual das normas e
dos traços. Para todo $X/S$ separado, a lei norma tem como alvo canônico o espaço de potências divididas $\Gamma^n_S(X)$. Se $X$ for plano sobre $S$ e a $n$-ésima potência simétrica existir (por exemplo, quando as órbitas de $\mathfrak S_n$ em $(X/S)^n$ estiverem contidas em abertos afins), o morfismo canônico $\operatorname{Symm}^n_S(X)\to\Gamma^n_S(X)$ será um isomorfismo, e obtém-se um elemento canônico
\[
\pi_{X/S}(\mathcal F)\in\Gamma(\operatorname{Symm}^n_S(X)/S),
\]
\src{274}
que permite recuperar os $\pi^u_{X/S}(\mathcal F)$. Outro caso importante é
aquele em que $X$ é um monoide comutativo sobre $S$, e $X=Y$, provindo o
morfismo $u$ da lei de composição de $X$. Escrever-se-á então simplesmente
$\pi(\mathcal F)$ para a seção de $X$ sobre $S$ associada ao Módulo
$\mathcal F$ sobre $X$.

Suponhamos agora que se tenha um Módulo coerente $\mathcal F$ sobre $X$, tal
que $\operatorname{Quot}_{\mathcal F/X/S}$ exista, ou pelo menos tal que o
funtor $\underline{\operatorname{Quot}}^n_{\mathcal F/X/S}$, que associa a
todo $S'$ sobre $S$ o conjunto dos feixes quocientes coerentes $\mathcal M$ de
$\mathcal F'=\mathcal F\otimes_{\mathcal O_S}\mathcal O_{S'}$ que são planos
sobre $S'$ e de posto relativo $n$, seja representável por um $S$-pré-esquema
$\operatorname{Quot}^n_{\mathcal F/X/S}$. (Quando $X$ é quase projetivo sobre
$S$, então $\operatorname{Quot}^n_{\mathcal F/X/S}$ existe de fato e não é
outro, com as notações do número 3, senão
$\operatorname{Quot}^P_{\mathcal F/X/S}$, onde $P$ é o polinômio reduzido ao
termo constante $n$.) Como a formação dos $\pi^u_{X/S}(\mathcal M)$ é
compatível com a mudança de base, obtém-se então um morfismo canônico
\[
\pi^u_{X/S}:\operatorname{Quot}^n_{\mathcal F/X/S}\longrightarrow Y,
\]
em geral:
\[
\pi_{X/S}:\operatorname{Quot}^n_{\mathcal F/X/S}
\longrightarrow\Gamma^n_S(X).
\]
Se $X$ for plano sobre $S$, identifica-se este alvo com $\operatorname{Symm}^n_S(X)$.
Independentemente da identificação precedente, quando $X$ é separado sobre $S$, a construção pela família finita e plana para $\mathcal F=\mathcal O_X$ fornece um morfismo:
\[
\pi_{X/S}:\operatorname{Hilb}^n_{X/S}
\longrightarrow\operatorname{Symm}^n_S(X).
\]
Este é evidentemente um isomorfismo para $n=0$ ou $n=1$. Mas, para
$n>1$\footnote{O texto impresso repete $n\geq1$, em contradição com o
isomorfismo que acaba de ser afirmado para $n=1$.}, mesmo se $S$ é o espectro
de um corpo $k$, e quando $X$ é liso sobre $S$, ele não é, em geral, um
isomorfismo, nem sequer um morfismo injetivo, pois um subesquema de dimensão 0
de $X$ (correspondente, por exemplo, a um ideal $\mathfrak i$ primário para o
ideal maximal em um anel local $\mathcal O_{X,x}$, para um ponto fechado $x$ de
$X$) não é conhecido quando se conhece o ciclo que ele define (no caso, quando
se conhece a codimensão sobre $k$ de $\mathfrak i$ em $\mathcal O_{X,x}$).
Pode-se dizer apenas o seguinte (sendo $S$ novamente arbitrário):

\medskip
\noindent a. Se $X$ é liso sobre $S$, então o morfismo-norma define um
isomorfismo do aberto de $\operatorname{Hilb}^n_{X/S}$ que corresponde à
classificação dos recobrimentos étales de posto $n$ contidos em $X$ (cf.
n\textsuperscript{o} 4, (b)) com o aberto de
$\operatorname{Symm}^n_S(X)$ que corresponde aos $n$-ciclos sem componentes
múltiplas.
\src{275}

\medskip
\noindent b. Se, além disso, $X$ tem dimensão relativa 1 sobre $S$, então o
morfismo norma define até mesmo um isomorfismo de
$\operatorname{Hilb}^n_{X/S}$ com $\operatorname{Symm}^n_{X/S}$.

Este último fato deve-se a que um subesquema de dimensão 0 de uma curva
algébrica não singular é conhecido quando se conhece o ciclo correspondente.
A mesma observação aplica-se, aliás, mais geralmente aos divisores de Cartier
positivos sobre um esquema algébrico não singular (e não se exclui que, nesse
caso muito particular, a variedade de Chow forneça a mesma coisa que o esquema
de Hilbert).

\fsection{7}{Complementos e questões}

Como observou J.-P. Serre, resulta de um exemplo bem conhecido de Nagata que se
pode encontrar um esquema $S$, espectro de um corpo $k$, um $S$-esquema $S'$,
espectro de uma extensão quadrática $k'$ de $k$, enfim, um $S'$-esquema $X$
próprio e liso (mas não projetivo) de dimensão 3, tal que
$\Pi_{S'/S}(X/S')$ não existe. Isso implica, a fortiori, que o esquema de
Hilbert $\operatorname{Hilb}^2_{X/S}$ não existe (nem sequer o $k$-esquema que
representaria os recobrimentos étales de posto 2 de $S$ contidos em $X$, nem,
a fortiori, o quadrado simétrico de $X$, cf. número precedente). Isso impõe,
portanto, limitações sérias às possibilidades de construções não projetivas em
geometria algébrica. (É, contudo, plausível que tais limitações não se
apresentem em geometria analítica, assim como não se apresentam em geometria
formal (cf. [2], II).) Em contrapartida, se $X$ é um esquema próprio sobre o
espectro $S$ de um corpo $k$, e se $Z$ é quase projetivo sobre $X$, então
$\Pi_{X/S}(Z/X)$ existe e é um esquema, soma disjunta de uma sequência de
esquemas quase projetivos sobre $S$ (como no caso projetivo (3.1)). Para vê-lo,
reduz-se, com efeito, ao caso em que o próprio $X$ é projetivo, dominando $X$
por um $S$-esquema projetivo $X'$; não daremos aqui os detalhes da demonstração,
que utiliza também o resultado de fatoração de um morfismo finito assinalado em
[2], I, A, n\textsuperscript{o} 2 (b). O êxito do método deve-se ao fato de que,
sendo $S$ o espectro de um corpo, o $X'$ que intervém no lema de Chow será
automaticamente plano sobre $S$.

Ignoro se o resultado continua válido sem hipótese sobre $S$, supondo apenas
$X$ próprio e plano sobre $S$, e $Z$ quase projetivo sobre $X$. Um caso
importante nas aplicações é aquele em que $Z$ é um subesquema fechado de $X$;
se, então, $\Pi_{X/S}(Z/X)$ existe, ele é necessariamente um subesquema fechado
de $S$. Pode-se
\src{276}
construí-lo diretamente de maneira bastante simples quando $X$ é projetivo
sobre $S$, sem utilizar a teoria dos esquemas de Hilbert, e o método empregado
mostra, mais geralmente, que, se $Z$ é afim sobre $X$, então
$\Pi_{X/S}(Z/X)$ existe e é afim sobre $S$. Ele mostra igualmente que, se $X$
é próprio e plano sobre $S$ (não necessariamente projetivo sobre $S$), então,
para todo fibrado vetorial $Z$ localmente trivial sobre $X$,
$\Pi_{X/S}(Z/X)$ existe e é um fibrado vetorial sobre $S$. Seria desejável que
esses resultados fossem retomados e unificados.

\section*{BIBLIOGRAFIA}
\addcontentsline{toc}{section}{Bibliografia}

\begin{description}
\item[{[1]}] BOREL (A.) e SERRE (J.-P.). -- Le théorème de Riemann--Roch,
\emph{Bull. Soc. math. France}, v. 86, 1958, p. 97--136.

\item[{[2]}] GROTHENDIECK (Alexander). -- Technique de construction et
théorèmes d'existence en géométrie algébrique, I, II, III, \emph{Séminaire
Bourbaki}, v. 13, 1959/60, n\textsuperscript{o} 190, 29 p.,
n\textsuperscript{o} 195, 22 p.; e v. 14, 1960/61,
n\textsuperscript{o} 212, 20 p.

\item[{[3]}] GROTHENDIECK (Alexander). -- Séminaire de Géométrie algébrique,
I, II, III, IV. -- Paris, Institut des hautes Études scientifiques, 1960/61
(mimeografado).

\item[{[4]}] GROTHENDIECK (Alexander). -- Technique de construction en
géométrie analytique, IV, V, \emph{Séminaire Cartan}, v. 13, 1960/61,
n\textsuperscript{o} 11 e 12.

\item[{[5]}] GROTHENDIECK (A.) e DIEUDONNÉ (J.). -- \emph{Éléments de
géométrie algébrique}, I: Le langage des schémas. -- Paris, Presses
universitaires de France, 1960 (Institut des hautes Études scientifiques,
Publications mathématiques, 4); II, III, IV (a aparecer).

\item[{[6]}] KODAIRA (K.). -- Characteristic linear systems of complete
continuous systems, \emph{Amer. J. of Math.}, v. 78, 1956, p. 716--744.

\item[{[7]}] SERRE (Jean-Pierre). -- Exemples de variétés projectives en
caractéristique $p$ non relevables en caractéristique 0, \emph{Proc. Nat. Acad.
Sc. U.S.A.}, v. 47, 1961, p. 108--109.

\item[{[8]}] ZAPPA (G.). -- Sull'esistenza, sopra le superficie algebriche, di
sistemi continui completi infiniti, la cui curva è a serie caratteristica
incompleta, \emph{Pont. Acad. Sc. Acta}, v. 9, 1945, p. 91--93.
\end{description}

\section*{ADENDO}
\addcontentsline{toc}{section}{Adendo}
\begin{center}
{\small [acrescentado durante a correção das provas]}
\end{center}

Parece agora que as conjecturas de [2], III, n\textsuperscript{o} 8, são falsas,
mesmo para variedades não singulares sobre um corpo de característica 0, tanto
no que concerne à existência quanto à quase projetividade do quociente, e mesmo
quando $G$ opera com um grafo fechado.

\par

\def\FgaUnitKey{e221}\clearpage
\chapter*{Errata da exposição n\textsuperscript{o} 221}
\addcontentsline{toc}{chapter}{Errata da exposição n\textsuperscript{o} 221}
\setcounter{footnote}{3}
\markboth{Errata da exposição 221}{Errata da exposição 221}
\sid{e221}
\src{302}

\noindent\textbf{6. TDTE IV: Os esquemas de Hilbert.}

\begin{center}
[Seminário Bourbaki, t. 13, 1960/61, n\textsuperscript{o} 221, 28 p.]
\end{center}

\noindent\textbf{Página 6, teorema 2.1, 3\textsuperscript{a} linha.} --- Em vez de
«é necessário», leia-se «é necessário e suficiente».

\noindent\textbf{Página 15, linha 13.} --- Em vez de «$P\neq Q$», leia-se
«$P<Q$».

\noindent\textbf{Linha 15.} --- Em vez de «Se, ao contrário, $P=Q$», leia-se
«No caso contrário, teremos $P(n)\geq Q(n)$ para $n$ grande».

\noindent\textbf{Linha 18, fórmula (*).} --- Em vez de «$<$» e «$>$»,
leia-se «$\leq$» e «$\geq$».

\noindent\textbf{5.ª linha a partir de baixo.} --- Em vez de «$r$», leia-se
«$Z$».

\noindent\textbf{Página 18, observação 3.9.} --- Assinalemos que o estudo das
componentes conexas dos esquemas de Hilbert sobre um corpo algebricamente
fechado foi realizado por HARTSHORNE. Esse autor demonstra que os
$\operatorname{Hilb}^{P}_{r}$ são conexos e determina os pares $(r,P)$ para
os quais $\operatorname{Hilb}^{P}_{r}\neq\varnothing$\footnote{HARTSHORNE
(R.). --- Connectedness of the Hilbert scheme (Tese, Harvard, 1962).}.

\noindent\textbf{Página 23, 5.ª linha a partir de baixo.} --- Em vez de «uma
componente», leia-se «de uma componente».

\noindent\textbf{Páginas 23 e 24, observação 5.5.} --- A propósito do exemplo
de ZAPPA, assinalemos que MUMFORD acaba até mesmo de construir uma componente
irredutível do esquema de Hilbert para $\mathbf P^3_{\mathbf C}$ (cujos pontos
gerais representam curvas não singulares de grau 14 e gênero 24), que não é
reduzida em seu ponto genérico. Realizando o blow-up de uma das curvas obtidas,
ele obtém também um esquema projetivo regular de dimensão 3 sobre $\mathbf C$,
cujo esquema formal de moduli não é reduzido em seu ponto genérico ou, o que
equivale ao mesmo, cuja variedade local de moduli, no sentido da geometria
analítica, sobre $\mathbf C$ (cf. Seminário Cartan, t. 13, 1960/61), não é
reduzida em nenhum de seus pontos.

\par

\def\FgaUnitKey{232}\unit{Exposição 232 -- Os esquemas de Picard: teoremas de existência}
\sid{232}
\markboth{Exposição 232}{Exposição 232}
\src{143}
\seminares{14.º ano, 1961/62, n.º 232\\Fevereiro de 1962}
\paperhead{TÉCNICA DE DESCIDA E TEOREMAS DE EXISTÊNCIA EM GEOMETRIA ALGÉBRICA\\
V. OS ESQUEMAS DE PICARD: TEOREMAS DE EXISTÊNCIA}%
  {por Alexander GROTHENDIECK}

\fsection{1}{Grupos e funtores de Picard relativos}

Para todo pré-esquema (mais geralmente, todo espaço anelado) $X$, chamamos
\emph{grupo de Picard} (absoluto) de $X$, e denotamos por
$\operatorname{Pic}(X)$, o grupo das classes, a menos de isomorfismo, de
Módulos invertíveis (isto é, localmente isomorfos a $\mathcal O_X$) sobre $X$.
Tem-se, portanto, um isomorfismo canônico
\begin{equation}
\operatorname{Pic}(X)\xrightarrow{\sim}H^1(X,\mathcal O_X^*),
\tag{1.1}
\end{equation}
onde $\mathcal O_X^*$ denota o feixe das unidades de $\mathcal O_X$ (que se
identifica, com efeito, com o feixe dos automorfismos do Módulo invertível
padrão $\mathcal O_X$). Observemos que
$X\longmapsto\operatorname{Pic}(X)$ é evidentemente um funtor contravariante
em $X$, e que o isomorfismo (1.1) é funtorial.

Se $X$ é um pré-esquema sobre outro pré-esquema $S$, então, quando $S'$ varia
na categoria $(\mathrm{Sch})/S$ dos pré-esquemas sobre $S$, obtém-se, pelo que
precede, um funtor contravariante
$S'\longmapsto\operatorname{Pic}(X\times_S S')$. Esse funtor não tem qualquer
possibilidade de ser «representável» ([4], II, A), pois, devido à existência de
automorfismos não triviais dos Módulos invertíveis que se pretende classificar,
esse funtor não é de «natureza local» ([5], IV, 5.4). Convém, portanto,
«torná-lo local», introduzindo, de modo geral, para todo pré-esquema relativo
$X/S$, um grupo de natureza relativa,
\begin{equation}
\operatorname{Pic}'(X/S)=H^0\bigl(S,R^1f_*(\mathcal O_X^*)\bigr),
\tag{1.2}
\end{equation}
(onde $f:X\longrightarrow S$ é o morfismo estrutural) (compare-se [4], II,
C 3). No lugar citado, esse grupo é chamado grupo de Picard relativo; será
preferível aqui chamá-lo \emph{grupo de Picard relativo restrito} de $X/S$,
por razões que aparecerão. Quando $S'$ varia em $(\mathrm{Sch})/S$,
$S'\longmapsto\operatorname{Pic}'(X\times_S S'/S')$ é um funtor contravariante
em $S'$, denotado também por $\operatorname{Pic}'_{X/S}$, portanto definido
essencialmente pela fórmula
\begin{equation}
\operatorname{Pic}'_{X/S}(S')=
\operatorname{Pic}'(X\times_S S'/S').
\tag{1.3}
\end{equation}
\src{144}
Este funtor é agora «de natureza local», visto que se fez o necessário para
isso. Intuitivamente, o segundo membro de (1.3) interpreta-se como o conjunto
das «famílias algébricas» de classes de feixes invertíveis sobre (as fibras de)
$X/S$, indexadas pelo pré-esquema de parâmetros $S'/S$. Quando o funtor
$\operatorname{Pic}'$ é representável, o pré-esquema sobre $S$ que o representa
é denotado por $\operatorname{Pic}'_{X/S}$ e chamado o
\emph{pré-esquema de Picard} de $X$ sobre $S$; portanto, teremos
\begin{equation}
\operatorname{Hom}_S(S',\operatorname{Pic}'_{X/S})
\simeq\operatorname{Pic}'_{X/S}(S')
=\operatorname{Pic}'(X\times_S S'/S').
\tag{1.4}
\end{equation}

Há, contudo, casos importantes em que $\operatorname{Pic}'_{X/S}$ não é
representável (exemplo: variedade de «Brauer--Severi» sobre um corpo $k$, sem
ponto racional sobre $k$), e nos quais, contudo, existe uma definição natural
de um pré-esquema de Picard relativo. Isso decorre do fato de que, na definição
do funtor $\operatorname{Pic}'$ a partir dos grupos de Picard absolutos
$\operatorname{Pic}(X\times_S S'/S')$, ainda não se localizou o suficiente;
mais precisamente, $\operatorname{Pic}'$ não é, em geral, «compatível com a
descida fielmente plana». Explicitemos.

Seja $\mathfrak M$ o conjunto dos morfismos de pré-esquemas que são
\emph{fielmente planos e quase-compactos}; esse conjunto é estável por mudança
de base e por composição. Seja $P$ um funtor contravariante de
$(\mathrm{Sch})/S$ na categoria dos conjuntos e, para todo $S$-morfismo
$u:T'\longrightarrow T$, $u\in\mathfrak M$, consideremos o diagrama
\begin{equation}
P(T)\longrightarrow P(T')\rightrightarrows P(T'\times_TT')
\tag{1.5}
\end{equation}
obtido pela aplicação de $P$ ao diagrama
\[
T\longleftarrow T'\mathrel{\substack{\xleftarrow{\,\mathrm{pr}_1\,}\\[-0.35ex]
\xleftarrow[\,\mathrm{pr}_2\,]{}}}T'\times_TT'.
\]
Se $P$ é representável, resulta da teoria da descida ([4], I, B, teor. 2) que
o diagrama (1.5) é exato para todo $u\in\mathfrak M$. Exprime-se esse fato
dizendo que $P$ é compatível com $\mathfrak M$, neste caso que $P$ é
«compatível com a descida fielmente plana», ou ainda que o «pré-feixe» $P$
sobre $(\mathrm{Sch})/S$ é um «feixe» para a noção de localização fornecida
pelo conjunto $\mathfrak M$. Quando $P$ é arbitrário, um procedimento padrão,
bem conhecido no caso da localização topológica usual, permite associar-lhe
um «feixe» $\mathcal P$ e um homomorfismo de funtores
$P\longrightarrow\mathcal P$, que seja universal num sentido evidente. O
cálculo de $\mathcal P$ pode ser explicitado da seguinte maneira: para definir
$\mathcal P(T)$, designa-se, para todo $T'$ sobre $T$ tal que o morfismo
$u:T'\longrightarrow T$ pertença a $\mathfrak M$, por
\src{145}
$\overline H^0(T'/T,P)$ o subconjunto de $P(T')$ formado pelos elementos $\xi$
cujas duas imagens $\xi_1,\xi_2$ em $P(T'\times_TT')$ são tais que existe um
morfismo $v:T''\longrightarrow T'\times_TT'$, $v\in\mathfrak M$, para o qual
$\xi_1$ e $\xi_2$ têm a mesma imagem em $P(T'')$. [N.B. -- O conjunto
$\overline H^0$ assim definido é, portanto, maior que o conjunto
$H^0(T'/T,P)$ introduzido em [4], I, A, 4(a).] Quando $T'$ varia sobre $T$
fixo (sempre com $u\in\mathfrak M$), os $\overline H^0(T'/T,P)$ formam um
sistema indutivo (quando se mune o conjunto dos $T'$ da pré-ordem definida pela
dominação), e põe-se
\begin{equation}
\mathcal P(T)=\varinjlim_{T'}\overline H^0(T'/T,P).
\tag{1.6}
\end{equation}
A lei funtorial em $T$ dessa expressão explicita-se de modo evidente.

Quando
\[
P(T)=\operatorname{Pic}(X\times_ST),
\]
o funtor contravariante sobre $(\mathrm{Sch})/S$ definido por (1.6) denomina-se
\emph{funtor de Picard relativo} de $X$ sobre $S$ e denota-se por
$\operatorname{Pic}_{X/S}$; denomina-se \emph{grupo de Picard relativo} de
$X$ sobre $S$, e denota-se por $\operatorname{Pic}(X/S)$, o grupo
$\operatorname{Pic}_{X/S}(S)$. Obtém-se então uma bijeção evidente
\begin{equation}
\operatorname{Pic}_{X/S}(T)\simeq\operatorname{Pic}(X\times_ST/T).
\tag{1.7}
\end{equation}

Um elemento de $\operatorname{Pic}(X/S)$ é, portanto, definido com o auxílio de
um elemento $\xi'$ de um grupo $\operatorname{Pic}(X\times_SS')$ (onde
$S'\longrightarrow S$ é fielmente plano e quase-compacto), tal que se possa
encontrar um morfismo fielmente plano e quase-compacto
$S''\longrightarrow S'\times_SS'$ para o qual as duas imagens inversas de
$\xi'$ em $\operatorname{Pic}(X\times_SS'')$ sejam iguais. Um elemento $\xi'$
de $\operatorname{Pic}(X\times_SS')$ e um elemento $\xi_1$ de
$\operatorname{Pic}(X\times_SS_1)$ (satisfazendo as condições que acabamos de
explicitar) definem o mesmo elemento de $\operatorname{Pic}(X/S)$ se e somente
se existe um morfismo fielmente plano e quase-compacto
$S'_1\longrightarrow S'\times_SS_1$ tal que as imagens dos dois elementos em
questão em $\operatorname{Pic}(X\times_SS'_1)$ sejam iguais. É frequentemente
conveniente continuar trabalhando com o funtor
$P'=\operatorname{Pic}'_{X/S}$ introduzido acima; constata-se imediatamente que
o morfismo canônico $P\longrightarrow P'$ define um isomorfismo
\begin{equation}
\mathcal P\xrightarrow{\sim}\mathcal P',
\tag{1.8}
\end{equation}
o que fornece uma descrição de $\operatorname{Pic}_{X/S}$ em termos de
$\operatorname{Pic}'_{X/S}=P'$, frequentemente mais cômoda. Em virtude de 2.3
abaixo, quando, na descrição de $\operatorname{Pic}(X/S)$ que acabamos de
explicitar, se substitui $P$ por $P'$, pode-se, com efeito,
\src{146}
tomar $S''=S'\times_SS'$, $S'_1=S'\times_SS_1$, ao menos sob as condições
explicitadas em loco citato.

Quando o funtor $\operatorname{Pic}_{X/S}$ é representável, diz-se que $X/S$
admite um pré-esquema de Picard, e o pré-esquema sobre $S$ que representa o
funtor denomina-se \emph{pré-esquema de Picard} de $X$ sobre $S$ e continua a
ser denotado por $\operatorname{Pic}_{X/S}$. Para isso, basta evidentemente que
$P'=\operatorname{Pic}'_{X/S}$ seja representável, pois então $P'$ já é um
``feixe'', e a fórmula (1.8) demonstra que o morfismo
$P'\longrightarrow\mathcal P'$ se identifica com o morfismo canônico
\begin{equation}
\operatorname{Pic}'_{X/S}\longrightarrow\operatorname{Pic}_{X/S},
\tag{1.9}
\end{equation}
que é então um isomorfismo. Portanto, nossa terminologia é compatível com a
introduzida acima em (1.4). Em geral, quando $\operatorname{Pic}_{X/S}$ existe,
ele é definido pelo isomorfismo funtorial
\begin{equation}
\operatorname{Hom}_S(S',\operatorname{Pic}_{X/S})\xrightarrow{\sim}
\operatorname{Pic}(X\times_SS'/S').
\tag{1.10}
\end{equation}

\fsection{2}{Relações entre os diversos grupos de Picard relativos e absolutos}

\result{PROPOSIÇÃO}{2.1} Seja $f:X\longrightarrow S$ um morfismo tal que
$\mathcal O_S\xrightarrow{\sim}f_*(\mathcal O_X)$. Tem-se então uma sequência
exata
\[
0\longrightarrow\operatorname{Pic}(S)\longrightarrow\operatorname{Pic}(X)
\longrightarrow\operatorname{Pic}'(X/S).
\]
Quando $X$ admite uma seção sobre $S$, o último morfismo é sobrejetivo, isto é,
tem-se um isomorfismo
\[
\operatorname{Pic}'(X/S)\xrightarrow{\sim}
\operatorname{Pic}(X)/\operatorname{Pic}(S).
\]

A sequência exata pode ser considerada como a sequência exata em graus baixos
correspondente à sequência espectral de Leray para $f$ e $\mathcal O_X^*$. A
segunda afirmação também é formal.

\result{PROPOSIÇÃO}{2.2} Seja $f:X\longrightarrow S$ um morfismo quase-compacto
e separado tal que $\mathcal O_S\xrightarrow{\sim}f_*(\mathcal O_X)$, e seja
$S'\longrightarrow S$ um morfismo fielmente plano e quase-compacto. Então:

\noindent (i) $\operatorname{Pic}'(X/S)\longrightarrow
\operatorname{Pic}'(X\times_SS'/S')$ é injetivo;

\noindent (ii) se $X$ admite localmente uma seção sobre $S$ (isto é, todo
$s\in S$ possui uma vizinhança aberta $U$ tal que $X|U$ tenha uma seção sobre
$U$), então o diagrama
\[
\operatorname{Pic}'(X/S)\longrightarrow
\operatorname{Pic}'(X\times_SS'/S')\rightrightarrows
\operatorname{Pic}'(X\times_SS''/S'')
\]
\src{147}
(onde $S''=S'\times_SS'$) é exato.

A primeira afirmação resulta, graças às propriedades elementares da descida
fielmente plana, da observação geral seguinte. Se $f:X\longrightarrow S$ é um
morfismo tal que $\mathcal O_S\xrightarrow{\sim}f_*(\mathcal O_X)$, então o
funtor $\mathcal F\longmapsto f^*(\mathcal F)$, da categoria dos Módulos
localmente livres de tipo finito sobre $S$ para a categoria dos Módulos
localmente livres de tipo finito sobre $X$, é plenamente fiel, e sua imagem
essencial é formada pelos Módulos $\mathcal G$ sobre $X$ tais que
$f_*(\mathcal G)$ seja localmente livre e que o homomorfismo canônico
\[
f^*\bigl(f_*(\mathcal G)\bigr)\longrightarrow\mathcal G
\]
seja um isomorfismo. A segunda afirmação foi demonstrada pela teoria da descida
em [4], I, B, 4.

Os resultados de 2.2 também se enunciam assim:

\result{COROLÁRIO}{2.3} Sob as condições de 2.2, o homomorfismo canônico
(1.9) $\operatorname{Pic}'_{X/S}\longrightarrow\operatorname{Pic}_{X/S}$ é
injetivo, e até bijetivo se $X$ admite uma seção localmente sobre $S$. (Logo,
neste último caso, há identidade entre o grupo de Picard relativo
$\operatorname{Pic}(X/S)$ e o grupo de Picard relativo restrito
$\operatorname{Pic}'(X/S)$.)

Combinando com 2.1, obtém-se, portanto:

\result{COROLÁRIO}{2.4} Sob as condições de 2.2, tem-se uma sequência exata
\[
0\longrightarrow\operatorname{Pic}(S)\longrightarrow\operatorname{Pic}(X)
\longrightarrow\operatorname{Pic}(X/S).
\]
Quando $X$ admite uma seção sobre $S$, o último homomorfismo é sobrejetivo,
isto é, tem-se então um isomorfismo
\[
\operatorname{Pic}(X/S)\xrightarrow{\sim}
\operatorname{Pic}(X)/\operatorname{Pic}(S).
\]

\result{OBSERVAÇÃO}{2.5} Seja $f:X\longrightarrow S$ um morfismo tal que
$\mathcal O_S\xrightarrow{\sim}f_*(\mathcal O_X)$, e seja $g$ uma seção de
$X$ sobre $S$. Seja $\mathcal L$ um Módulo invertível sobre $X$; chamemos
$g$-rigidificação de $\mathcal L$ um isomorfismo
$\mathcal O_S\xrightarrow{\sim}g^*(\mathcal L)$, e chamemos \emph{Módulo
invertível $g$-rigidificado} um Módulo invertível $\mathcal L$ sobre $X$ munido
de uma $g$-rigidificação. Todo automorfismo de tal estrutura é trivial, e
$\operatorname{Pic}'(X/S)$ identifica-se com o conjunto das classes, a menos de
isomorfismo, de Módulos invertíveis $g$-rigidificados sobre $X$\footnote{O
texto impresso traz «sobre $S$»; a rigidificação por uma seção $g:S\to X$
concerne aqui ao Módulo $\mathcal L$ sobre $X$.}. (É esse fato que permitiu
utilizar a teoria da descida para demonstrar 2.2(ii).) Isso fornece uma nova
interpretação de $\operatorname{Pic}(X/S)$, pelo menos quando $f$ é, além disso,
quase-compacto e separado, portanto
$\operatorname{Pic}(X/S)\simeq\operatorname{Pic}'(X/S)$ por 2.3.
\src{148}
\result{OBSERVAÇÃO}{2.6} Seja $f:X\longrightarrow S$ um morfismo como em
2.2, e seja $S'\longrightarrow S$ um morfismo fielmente plano e
quase-compacto tal que exista um $S$-morfismo $S'\longrightarrow X$, isto é,
tal que exista uma seção de $X'=X\times_SS'$ sobre $S'$. Sejam
$S''=S'\times_SS'$ e $X''=X\times_SS''$, e consideremos a sequência exata
\[
\operatorname{Pic}(X/S)\longrightarrow\operatorname{Pic}(X'/S')
\rightrightarrows\operatorname{Pic}(X''/S'').
\]
Aplicando 2.4, obtém-se uma sequência exata
\[
\operatorname{Pic}(X/S)\longrightarrow
\operatorname{Pic}(X')/\operatorname{Pic}(S')
\rightrightarrows
\operatorname{Pic}(X'')/\operatorname{Pic}(S''),
\]
em particular, todo elemento do grupo de Picard relativo já «provém» de um
elemento de $\operatorname{Pic}(X')$. Isso fornece, portanto, uma simplificação
substancial da descrição do grupo de Picard relativo dada no número anterior,
e o mesmo vale para o funtor de Picard de $X$ sobre $S$, pois, para todo $T$
sobre $S$, pode-se aplicar o que precede a $X\times_ST/T$ e ao morfismo
$T'=S'\times_ST\longrightarrow T$. Se, por exemplo, o próprio $f$ é fielmente
plano, pode-se tomar $S'=X$, o que permite, quando $f$ é, além disso, de tipo
finito (resp. suave, etc.), limitar-se, na descrição do funtor de Picard relativo
$T\longmapsto\operatorname{Pic}(X\times_ST/T)$, às mudanças de base
$T'\longrightarrow T$ que são de tipo finito (resp. suaves, etc.). Quando $f$
é projetivo e plano e $S$ é localmente noetheriano, demonstra-se que se pode
tomar acima um $S'\longrightarrow S$ tal que $S'$ seja soma disjunta de
recobrimentos planos $S'_i$ de abertos $S_i$ de $S$ que cobrem $S$; se $f$ é
até mesmo separável, pode-se tomar $S'_i$ étale sobre $S_i$.

\fsection{3}{O teorema principal de existência: enunciado}

Não se dispõe, nem mesmo a título conjectural, de um enunciado de existência
de pré-esquemas de Picard que englobe todos os casos conhecidos. Uma condição
«praticamente necessária», se assim se pode dizer, é que
$f:X\longrightarrow S$ seja próprio (assegurando propriedades essenciais de
finitude) e plano. Essas condições não são suficientes, mesmo se $S$ é o
espectro da álgebra dos números duais $k[t]/(t^2)$ sobre um corpo $k$ (digamos,
o corpo $\mathbf C$ dos números complexos), e $X$ tem dimensão 1. No momento de
escrever esta exposição, os teoremas de existência mais importantes para o
pré-esquema de Picard deduzem-se do teorema seguinte:

\result{TEOREMA}{3.1} Seja $f:X\longrightarrow S$ um morfismo de pré-esquemas
localmente noetherianos. Supõe-se:

\noindent (i) $f$ é projetivo;
\src{149}
\noindent (ii) $f$ é plano;

\noindent (iii) as fibras geométricas de $f$ são íntegras.

Sob essas condições, $\operatorname{Pic}_{X/S}$ existe.

A demonstração, que será esboçada nas duas seções seguintes, mostrará ao mesmo
tempo o seguinte: seja $\xi$ a seção de $\operatorname{Pic}_{X/S}$
correspondente a um feixe muito amplo $\mathcal O_X(1)$ sobre $X/S$ (isto é,
induzido por uma imersão projetiva
$X\longrightarrow\mathbf P(\mathcal E)$); existe então uma parte aberta $U$ de
$\operatorname{Pic}_{X/S}$, união disjunta de partes abertas quase-projetivas
sobre $S$, tal que $U$ seja estável pela translação por $\xi$, e que
$\operatorname{Pic}_{X/S}$ seja a união crescente dos abertos $U-n\xi$ (todos
isomorfos a $U$). Disso resulta, em particular, que, sob as condições de 3.1,
$\operatorname{Pic}_{X/S}$ é separado sobre $S$.

\result{OBSERVAÇÃO}{3.2} Vê-se por exemplos (com $S$ o espectro de um anel de
valorização discreta e $X$ de dimensão relativa 1 sobre $S$, por exemplo) que,
se em 3.1 se omite a hipótese (iii), substituindo-a pela hipótese mais fraca de
que, para todo $s\in S$, o homomorfismo
\[
k(s)\longrightarrow H^0(X_s,\mathcal O_{X_s})
\]
é um isomorfismo, então $\operatorname{Pic}_{X/S}$ não é necessariamente
separado sobre $S$; tanto no caso em que as fibras geométricas de $f$ são
reduzidas, mas uma fibra geométrica genérica íntegra ``explode'' por
especialização em duas componentes irredutíveis, como no caso em que as fibras
geométricas de $f$ são irredutíveis, mas uma fibra geométrica genérica íntegra
se especializa em uma ``fibra múltipla''. O primeiro caso ocorre, por exemplo,
com uma cônica que degenera em duas retas concorrentes; um exemplo do segundo me
foi fornecido por D. Mumford, com uma curva elíptica que degenera em uma curva
elíptica dupla. Esses exemplos são válidos em toda característica.

\result{OBSERVAÇÃO}{3.3} Sob as condições de 3.1, ignoro se
$\operatorname{Pic}_{X/S}$ é uma união disjunta de abertos que são de tipo
finito, portanto quase-projetivos, sobre $S$.
\footnote{A errata de 1962 assinala que D. Mumford resolveu afirmativamente essa
questão por meio da teoria dos esquemas de Picard desenvolvida aqui.} Observe-se
que a consideração dos polinômios de Hilbert $Q\in\mathbf Q[t]$ permite, como no
caso dos esquemas de Hilbert ([4], IV), decompor
$\operatorname{Pic}_{X/S}$ em uma soma disjunta de abertos
$\operatorname{Pic}^{Q}_{X/S}$, e parece plausível que estes sejam
quase-projetivos sobre $S$; isso é o que se verá, ao menos, na próxima exposição,
quando $f$ for um morfismo suave. Atente-se para o fato de que, se se substituir
a hipótese (i) pela hipótese: $X$ é localmente projetivo sobre $S$ (o que basta
para a validade de 3.1, pois a questão da existência de
$\operatorname{Pic}_{X/S}$ é manifestamente local sobre $S$), é fácil, por
outro lado,
\src{150}
dar exemplos em que $\operatorname{Pic}_{X/S}$ contém componentes conexas que
não são de tipo finito sobre $S$. Seja, por exemplo, $X_0$ uma variedade
algébrica projetiva não singular sobre um corpo $k$ algebricamente fechado,
munida de um automorfismo $u$ e de um elemento $\xi$ do grupo de
Néron--Severi de $X_0$, tais que os $u^n(\xi)$ sejam dois a dois distintos.
Pode-se, por exemplo, tomar para $X_0$ o produto de uma curva elíptica $E$ por
si mesma e, para $u$, o automorfismo
\[
(x,y)\longmapsto(x,y+x)
\]
de $E\times E$. Seja $S$ a união de duas curvas irredutíveis não singulares que
se cortam em dois pontos $a$ e $b$. Existe sobre $S$ um recobrimento principal
conexo $P$ de grupo $\mathbf Z$; utilizando as ações de $\mathbf Z$ sobre $X_0$
definidas por $u$, obtém-se um fibrado associado sobre $S$, de fibra $X_0$
(trivial sobre $S-a$ e $S-b$), na realidade um esquema abeliano sobre $S$ no
caso particular considerado. Vê-se facilmente que
$\operatorname{Pic}_{X/S}$, que é também o fibrado associado a $P$ e às ações
de $\mathbf Z$ sobre $\operatorname{Pic}_{X_0/k}$ por meio de $u$, contém uma
componente conexa isomorfa a
\[
P\times\operatorname{Pic}^{0}_{X_0/k},
\]
(onde $\operatorname{Pic}^{0}$ denota a componente conexa do elemento neutro
em $\operatorname{Pic}$), a qual não é de tipo finito sobre $S$. [Fenômenos
análogos também ocorrem em diversos casos de pré-esquemas de Picard não
separados sobre $S$, como os considerados em 3.2.]

\fsection{4}{Divisores de Cartier relativos e fibrados projetivos}

Teremos de utilizar apenas divisores positivos e omitiremos esse qualificativo
adicional no restante deste item.

Seja $X$ um pré-esquema. Um \emph{divisor de Cartier} sobre $X$, ou simplesmente
divisor sobre $X$, para simplificar, é um subpré-esquema fechado $D$ de $X$,
definido por um Ideal $\mathcal J$ que é um Módulo invertível, isto é, gerado
localmente por uma seção que não é divisor de zero de $\mathcal O_X$. A $D$
associamos um Módulo invertível, a saber,
\[
\mathcal L(D)=\mathcal J^{-1},
\]
e a injeção canônica $\mathcal J\longrightarrow\mathcal O_X$ fornece um
homomorfismo canônico
\[
s_D:\mathcal O_X\longrightarrow\mathcal J^{-1}=\mathcal L(D),
\qquad\text{isto é,}\quad s_D\in\Gamma(X,\mathcal L(D)).
\]
Além disso, o dado de um divisor é essencialmente equivalente ao dado de um
Módulo invertível $\mathcal L$ sobre $X$, munido de uma seção $s$ que não seja
divisor de zero em ponto algum, associando reciprocamente a tal par
$(\mathcal L,s)$ o «divisor» de $s$, denotado por $\operatorname{div}(s)$. Para
um $\mathcal L$ invertível dado sobre $X$, o conjunto dos
\src{151}
divisores $D$ que definem $\mathcal L$ está em correspondência biunívoca com o
conjunto quociente
\[
\Gamma(X,\mathcal L)^*/\Gamma(X,\mathcal O_X^*),
\]
onde $\Gamma(X,\mathcal L)^*$ designa a parte de $\Gamma(X,\mathcal L)$ formada
pelas seções que não são divisores de zero em nenhum ponto.

Suponhamos agora que se tenha um morfismo localmente de tipo finito
$f:X\longrightarrow S$, e suponhamos, para simplificar, $S$ localmente
Noetheriano. Sejam $\mathcal J$ um Ideal coerente sobre $X$, $D$ o subesquema de
$X$ que ele define, $x\in X$, e $s=f(x)$. Demonstra-se que as seguintes
condições são equivalentes:

\noindent (i) $\mathcal J$ é invertível em $x$ (isto é, $\mathcal J_x$ é gerado
por um elemento regular de $\mathcal O_{X,x}$) e $D$ é plano sobre $S$ em $x$;

\noindent (ii) $X$ e $D$ são planos sobre $S$ em $x$, e $D_s$ é um divisor de
Cartier sobre a fibra $X_s$ no ponto $x$;

\noindent (iii) $X$ é plano sobre $S$ em $x$, e $\mathcal J_x$ é gerado por um
elemento $f_x$ que induz sobre $X_s$ um germe não divisor de zero.

Diz-se então que $D$ é um \emph{divisor de Cartier relativo} sobre $X/S$ (ou
simplesmente divisor relativo sobre $X/S$) no ponto considerado. Observe-se, por
(i), que $D$ é também um divisor relativo nos pontos vizinhos de $x$; portanto,
se $X$ e $D$ são planos sobre $S$, e $D$ é próprio sobre $S$, então o conjunto
dos $s\in S$ tais que $D_s$ seja um divisor de Cartier em $X_s$ (isto é, tais
que $D$ seja um divisor de Cartier relativo nos pontos de $X_s$) é uma parte
aberta de $S$. Por outro lado, na definição precedente fez-se o necessário para
que a noção de divisor de Cartier relativo seja estável por uma mudança de base
arbitrária $S'\longrightarrow S$. Consideremos então o conjunto
$\operatorname{Div}(X/S)$ dos divisores relativos sobre $X/S$, e em seguida o
funtor contravariante em $S'$ variável sobre $S$ definido por
\[
\operatorname{Div}_{X/S}(S')=
\operatorname{Div}(X\times_SS'/S').
\]

Suponhamos $X$ plano e próprio sobre $S$. Então, pela caracterização (ii) dos
divisores de Cartier relativos, $\operatorname{Div}_{X/S}$ pode ser considerado
como um subfuntor do funtor $\operatorname{Hilb}_{X/S}$ definido em [4], IV, e o
morfismo de inclusão
\[
\operatorname{Div}_{X/S}\longrightarrow\operatorname{Hilb}_{X/S}
\]
é ``representável por imersões abertas'' (cf. [5], IV, 3.13), em virtude das
observações precedentes. Utilizando o teorema principal de existência de [4],
IV, obtém-se:
\src{152}
\result{PROPOSIÇÃO}{4.1} Suponhamos que $f:X\longrightarrow S$ seja projetivo
e plano. Então o funtor $\operatorname{Div}_{X/S}$ é representável e, mais
precisamente, é representado por um aberto de $\operatorname{Hilb}_{X/S}$.

Utilizando (para um feixe $\mathcal O_X(1)$ muito amplo dado sobre $X/S$) a
decomposição canônica de $\operatorname{Hilb}_{X/S}$ em soma dos abertos
$\operatorname{Hilb}^{Q}_{X/S}$ correspondentes aos polinômios de Hilbert
$Q\in\mathbf Q[t]$, deduz-se, aliás, uma decomposição análoga
\[
\operatorname{Div}_{X/S}=
\coprod_{Q\in\mathbf Q[t]}\operatorname{Div}^{Q}_{X/S}
\]
em soma de abertos disjuntos quase-projetivos sobre $S$.

Utilizando a aplicação $D\longmapsto\mathcal L(D)$, obtém-se, por outro lado,
um homomorfismo funtorial
\begin{equation}
\operatorname{Div}_{X/S}\longrightarrow\operatorname{Pic}_{X/S}
\tag{+}
\end{equation}
que nos propomos estudar; ver-se-á que ele é relativamente representável ([5],
IV, 3) sob condições bastante gerais. Partamos, portanto, de um elemento $\xi$
de $\operatorname{Pic}_{X/S}(S')$, supondo, para simplificar a notação, que
$S'=S$, e mostremos que o subfuntor correspondente de
$\operatorname{Div}_{X/S}$ é representável. Tomemos primeiro o caso em que
$\xi$ é definido por um Módulo invertível $\mathcal L$ sobre $X$. Suporemos
que $X$ seja próprio e plano sobre $S$ e que as fibras geométricas de $X$
sobre $S$ sejam íntegras, o que implica também ([2], III, par. 7) que se tem
$\mathcal O_S\xrightarrow{\sim}f_*(\mathcal O_X)$ e que essa relação permanece
válida após toda mudança de base $S'\longrightarrow S$. Então os divisores de
Cartier relativos $D$ sobre $X/S$ tais que $\mathcal L(D)$ e $\mathcal L$
definem o mesmo elemento de
$\operatorname{Pic}(X/S)=\operatorname{Pic}_{X/S}(S)$, isto é, em virtude de
2.4, tais que $\mathcal L(D)$ e $\mathcal L$ sejam isomorfos localmente sobre
$S$, estão em correspondência biunívoca com as seções do feixe quociente
$f_*(\mathcal L)^*/\mathcal O_S^*$. Essa correspondência é compatível com as
mudanças de base. Por outro lado, considerações gerais do tipo «Künneth» de
loco citato (cf. também [6]) mostram que o fato de $X/S$ ser próprio e de
$\mathcal L$ ser plano sobre $S$ implica a existência de um Módulo coerente
$\mathcal Q$ sobre $S$, definido a menos de isomorfismo único, e de um
isomorfismo de feixes
\[
f_*(\mathcal L)\xrightarrow{\sim}
\operatorname{Hom}_{\mathcal O_S}(\mathcal Q,\mathcal O_S),
\]
\footnote{O texto impresso coloca $\mathcal O_X$ como subíndice de
$\operatorname{Hom}$; como $\mathcal Q$ é um Módulo sobre $S$, seu dual é
tomado sobre $\mathcal O_S$.} sendo, além disso, a formação de $\mathcal Q$
compatível com as mudanças de base. Aqui $f_*(\mathcal L)^*$ designa o
subfeixe de conjuntos de $f_*(\mathcal L)$ cujas seções sobre $U$ são as
seções de $\mathcal L$ sobre $f^{-1}(U)$ que definem divisores de Cartier
relativos sobre
\src{153}
Cartier relativos sobre $f^{-1}(U)/U$, isto é, que induzem seções que não são
divisores de zero sobre os $X_s$ ($s\in U$). Utilizando a hipótese de que as
fibras $X_s$ são íntegras, isso significa simplesmente que as seções induzidas
sobre as fibras $X_s$ não são identicamente nulas, ou ainda, em termos de
homomorfismos locais $\mathcal Q\longrightarrow\mathcal O_S$, que tais
homomorfismos sejam sobrejetivos (Nakayama). Isso mostra que o conjunto das
seções de $f_*(\mathcal L)^*/\mathcal O_S^*$ está em correspondência biunívoca
com o conjunto dos Módulos quocientes invertíveis de $\mathcal Q$, ou ainda,
pela definição do fibrado projetivo $\mathbf P(\mathcal Q)$ associado ao Módulo
coerente $\mathcal Q$ (cf. [5], V, 2), com o conjunto das seções de
$\mathbf P(\mathcal Q)$ sobre $S$. Essa descrição é compatível com a formação
das imagens inversas, obtendo-se, por conseguinte, o teorema seguinte.

\result{TEOREMA}{4.3} Sejam $f:X\longrightarrow S$ um morfismo próprio e plano
com fibras geométricas íntegras, sendo $S$ localmente noetheriano, e
$\mathcal L$ um Módulo invertível sobre $X$. Para todo $S'$ sobre $S$, seja
$T(S')$ o conjunto dos divisores relativos $D$ sobre $X\times_SS'/S'$ tais que
$\mathcal L(D)$ seja isomorfo a
$\mathcal L\otimes_{\mathcal O_S}\mathcal O_{S'}$ localmente sobre $S'$ (isto
é, tais que $\mathcal L(D)$ e
$\mathcal L\otimes_{\mathcal O_S}\mathcal O_{S'}$ definam o mesmo elemento de
$\operatorname{Pic}(X\times_SS'/S')$). Então existe um Módulo coerente
$\mathcal Q$ sobre $S$, determinado a menos de isomorfismo único, tal que o
funtor $T$ seja representável pelo fibrado projetivo $\mathbf P(\mathcal Q)$.

\result{COROLÁRIO}{4.4} Se supusermos, além disso, que $f$ é projetivo, então o
homomorfismo funtorial
$\operatorname{Div}_{X/S}\longrightarrow\operatorname{Pic}_{X/S}$ é
representável por morfismos projetivos.

Com efeito, quando $X$ admite uma seção (resp. localmente uma seção) sobre $S$,
o homomorfismo precedente é representável por fibrados projetivos (resp. por
fibrados projetivos locais), graças a 4.3 e 2.1. No caso em que $f$ é
quase-projetivo, reduz-se facilmente ao caso precedente por um método de
descida, utilizando quase-seções finitas e planas de $X$ localmente sobre $S$.

\result{OBSERVAÇÃO}{4.5} Sob as condições de 4.3, o Módulo $\mathcal Q$ sobre
$S$ não é, em geral, localmente livre, o que se exprime pelo fato de que a
dimensão das fibras reduzidas de $\mathcal Q$, isto é, a dos
$H^0(X_s,\mathcal L_s)$ para $s\in S$ variável, pode dar saltos. Dado um Módulo
coerente $\mathcal Q$ sobre o pré-esquema localmente noetheriano $S$, verifica-se
facilmente que, para um $s\in S$ dado, $\mathcal Q$ é livre em $s$ se e somente
se $\mathbf P(\mathcal Q)$ é plano sobre $S$ nos pontos acima de $s$ (caso em
que será até liso sobre $S$ nos pontos acima de $s$); no caso em questão,
estando $\mathcal Q$ definido em termos de $\mathcal L$ como acima, isso
significa também que a formação da imagem direta $f_*(\mathcal L)$ «comuta com
a mudança de base na
\src{154}
vizinhança de $s$», ou ainda que
$f_*(\mathcal L)_s\longrightarrow H^0(X_s,\mathcal L_s)$ é sobrejetivo. Isso
ocorre, por exemplo, se $H^1(X_s,\mathcal L_s)=0$. Sob reserva da existência
dos pré-esquemas em questão, esses critérios aplicam-se, em particular, à
situação universal
$\operatorname{Div}_{X/S}\longrightarrow\operatorname{Pic}_{X/S}$, e fornecem
uma condição necessária e suficiente, resp. suficiente, para que esse morfismo
seja liso em um ponto dado de $\operatorname{Div}_{X/S}$.

\fsection{5}{Demonstração do teorema principal de existência}

Sob as condições de 3.1, escolhamos um Módulo $\mathcal O_X(1)$ muito amplo
sobre $X/S$, e seja $\xi$ o elemento correspondente de
$\operatorname{Pic}(X/S)$. Para abreviar, ponhamos
\[
P(S')=\operatorname{Pic}(X\times_SS'/S'),
\]
e suponhamos primeiro, para simplificar, que $X/S$ admita uma seção. Seja
$P^+(S')$ a parte de $P(S')$ formada pelas classes dos $\mathcal L$ invertíveis
sobre $X\times_SS'$ tais que
\[
R^if'_*(\mathcal L(n))=0\quad\text{para $i>0$ e todo $n\geq0$},
\qquad
f'_*(\mathcal L(n))\neq0\quad\text{para todo $n\geq0$}.
\]
Essas condições são estáveis por mudança de base, definindo, portanto, um
subfuntor $P^+$ de $P$, evidentemente estável pela translação por $\xi$.
Utilizando os «teoremas A e B» de Serre ([2], III, 2) e as generalidades
([5], IV, 5), vê-se facilmente que $P$ é representável se e somente se $P^+$ o
é, e, nesse caso, $P^+$ será representado por um aberto $U$ do pré-esquema
$\operatorname{Pic}_{X/S}$ que representa $P$, e este último será a união
crescente dos abertos $U-n\xi$.

Para abreviar, seja $D=\operatorname{Div}_{X/S}$, e seja $D^+$ a imagem inversa
de $P^+$ pelo morfismo canônico $D\longrightarrow P$. Tem-se, portanto, um
morfismo
\begin{equation}
D^+\longrightarrow P^+,
\tag{+}
\end{equation}
e já se sabe que $D^+$ é representável por um aberto $D^+$ do pré-esquema
$D=\operatorname{Div}_{X/S}$ (cuja existência é garantida por 4.1). Resulta
facilmente de 4.3 que o morfismo (+) é representável por morfismos lisos,
projetivos e sobrejetivos (e, mais precisamente, por fibrados projetivos
associados a Módulos localmente livres não nulos em toda parte): isso decorre
do fato de que, quando $\mathcal L$ sobre $X\times_SS'$ é como no início deste
item, então $f'_*(\mathcal L)$ é um Módulo localmente livre $\neq0$, cuja
formação comuta com a mudança de base; com as notações de 4.3, $\mathcal Q$
será então isomorfo ao dual de $f'_*(\mathcal L)$. Utilizando o critério geral
([5], IV, 4.7), concluiremos daí a representabilidade de $P^+$. No lugar citado,
tomaremos para $\mathfrak C$ o conjunto dos morfismos fielmente
\src{155}
planos e quase-compactos de pré-esquemas (que são de fato epimorfismos
efetivos em virtude de [4], I B). A condição (a) de loco citato, a saber, que
(+) é representável por morfismos pertencentes a $\mathfrak C$, está verificada,
como acabamos de ver; o mesmo ocorre com a condição (b), que significa que o
funtor $P^+$ é compatível com a descida fielmente plana quase-compacta, o que
é imediato. Resta provar a condição (c) de loco citato, a saber, que a
equivalência $R$ no pré-esquema $D^+$ deduzida do morfismo
$\mathfrak C$-representável (+) é $\mathfrak C$-efetiva, isto é, é efetiva e
tal que $D^+\longrightarrow D^+/R$ pertença a $\mathfrak C$. Para isso,
observamos primeiro que os abertos $D^{+Q}$ de $D^+$ correspondentes aos
diversos polinômios de Hilbert virtuais $Q\in\mathbf Q[t]$ são estáveis por $R$
(pois as fibras de $R$ são conexas), o que nos reduz a mostrar que, para todo
$Q$, a relação de equivalência $R^Q$ induzida em $D^{+Q}$ é
$\mathfrak C$-efetiva. Ora, $D^{+Q}$ é agora quase-projetivo e, por outro lado,
a relação de equivalência $R^Q$ é projetiva e plana. Estamos, portanto, nas
condições de aplicação de [4], III, 6.1, o que implica o resultado desejado.

No caso geral em que $X/S$ não admite necessariamente uma seção, reduzimo-nos
facilmente ao caso anterior pela técnica de descida, ou retomamos a demonstração
anterior com a modificação necessária na definição de $P^+$.

\result{OBSERVAÇÃO}{5.1} O método seguido é essencialmente o de Matsusaka para
a construção projetiva das variedades de Picard. O resultado de [4], III
invocado para a possibilidade da passagem ao quociente efetivo também pode ser
facilmente deduzido do teorema de existência dos esquemas de Hilbert (cf., por
exemplo, [6]) (classicamente, esses quocientes eram construídos mediante o uso
das coordenadas de Chow). Observe-se que a formação do aberto
$\operatorname{Pic}^{+}_{X/S}$ de $\operatorname{Pic}_{X/S}$ e sua decomposição
em abertos $\operatorname{Pic}^{+Q}_{X/S}$ quase-projetivos sobre $S$, segundo
os polinômios de Hilbert dos divisores que definem os Módulos invertíveis
considerados, é compatível com a mudança de base (o que permite a aplicação da
técnica de descida).

\result{OBSERVAÇÃO}{5.2} Não se exclui que $\operatorname{Pic}_{X/S}$ exista
sempre que $f:X\longrightarrow S$ seja próprio e plano e tal que os
homomorfismos
\[
k(s)\longrightarrow H^0(X_s,\mathcal O_{X_s})\qquad(s\in S)
\]
sejam isomorfismos (significando então esta última condição também que
$\mathcal O_S\xrightarrow{\sim}f_*(\mathcal O_X)$ e que essa relação permaneça
válida após toda mudança de base $S'\longrightarrow S$). É ao menos o que se
prova no âmbito dos espaços analíticos quando $f$ é, além disso, projetivo, por
um método diferente (o de Chow, salvo engano) explicado em [5], IX,
3.1.\footnote{A errata de 1962 assinala que a conjectura de existência formulada
no início desta observação é falsa, mas que os métodos de Mumford fornecem um
teorema ligeiramente mais fraco.} Nesse método, a passagem ao
\src{156}
quociente por uma relação de equivalência própria e plana em um aberto do
esquema dos divisores é substituída por uma passagem ao quociente pelo grupo
projetivo no esquema das imersões de $X$ em $\mathbf P^r_S$. Esse método poderá
provavelmente adaptar-se ao caso dos esquemas utilizando os resultados de
Mumford sobre a passagem ao quociente pelo grupo projetivo [8]; por ora, só
existe demonstração escrita quando $X$ tem ``muitas seções locais'' sobre $S$,
por exemplo, quando $X$ é separável sobre um anel local completo com corpo
residual algebricamente fechado.
\footnote{O texto impresso omite ``com corpo residual algebricamente fechado'';
a errata esclarece que o contraexemplo de Mumford torna essa restrição
indispensável.} Em princípio, o método em questão teria alcance mais geral,
pois também fornece a existência de pré-esquemas de Picard nos casos em que
estes não são separados e nos quais, portanto, o primeiro método necessariamente
falha. (Tecnicamente, a dificuldade provém do fato de que, quando as fibras
geométricas de $f$ deixam de ser íntegras, o funtor considerado em 4.3 deixa de
ser representável pelo próprio fibrado projetivo $\mathbf P(\mathcal Q)$ e
passa a sê-lo por um aberto deste, o que conduz à delicada questão da passagem
ao quociente por uma relação de equivalência plana, mas não própria.)

\result{OBSERVAÇÃO}{5.3} Observe-se que a demonstração apresentada aqui não
recorre nem à construção prévia das jacobianas das curvas ou famílias de curvas,
nem à teoria das variedades abelianas ou dos esquemas abelianos, e por isso se
distingue essencialmente de exposições tradicionais como a do livro de Lang [7]
ou a exposição de Chevalley [1], que seguem a via esboçada por A. Weil. Mesmo no
caso das jacobianas das curvas não singulares sobre um corpo algebricamente
fechado (digamos, o corpo dos números complexos), a construção da jacobiana aqui
apresentada é a única conhecida que fornece as propriedades muito fortes que
adotamos como definição no número 1 (essencialmente as de Chevalley, mas levando
em conta ``variedades de parâmetros'' com elementos nilpotentes). Que a
construção dos esquemas de Picard deva preceder, e não seguir, a teoria das
variedades abelianas é claro a priori pelo fato de que, em geral, os esquemas de
Picard não são esquemas abelianos nem se reduzem a eles, como já se vê no caso
das curvas singulares sobre um corpo algebricamente fechado, em que se encontram
as ``jacobianas generalizadas'' de Rosenlicht, que não são variedades abelianas.
Além disso, a teoria das variedades abelianas, e mais geralmente a dos esquemas
abelianos, ganha muito em simplicidade quando se dispõe de uma teoria dos
esquemas de Picard em geral. Em particular, a teoria de dualidade para os
esquemas abelianos, sobretudo os resultados do tipo Cartier, torna-se quase
formal a partir daí (cf., por exemplo, [10]).
\src{157}
\result{OBSERVAÇÕES}{5.4} O «princípio de compatibilidade» de Igusa para a
jacobiana de uma curva que degenera em uma curva singular só pode ser bem
compreendido como um teorema de existência do esquema de Picard de um esquema
relativo em curvas $X/S$, não necessariamente suave sobre $S$. Trata-se,
portanto, de um caso particular do teorema principal de existência 3.1 quando a
curva especializada é íntegra (isto é, em termos clássicos, irredutível de
multiplicidade 1). Observe-se que, por ora, o caso de uma curva especial
redutível (mesmo quando as componentes têm multiplicidade 1, isto é, quando a
curva especial é separável sobre o corpo residual) escapa aos teoremas de
existência conhecidos, salvo no caso em que se está sobre um anel de valorização
discreta completo com corpo residual algebricamente fechado; compare-se 5.2.
Essa questão de existência certamente surgirá por ocasião de uma construção
geométrica, na teoria dos esquemas, das «compactificações» de Baily--Satake dos
esquemas modulares das curvas de gênero $g$. (Essa compactificação é conhecida,
por ora, apenas para $g=1$, graças aos trabalhos de Igusa.)

\fsection{6}{Teoremas de existência relativos}

Esboçaremos aqui alguns casos úteis em que a existência de certos esquemas de
Picard implica a existência de alguns outros, o que permite deduzir do teorema
principal 3.1 diversos outros teoremas de existência.

\result{PROPOSIÇÃO}{6.1} Seja $f:X\longrightarrow S$ um morfismo plano e
projetivo tal que, na fatorização de Stein $f=f''f'$, o morfismo
$f':X\longrightarrow S'$ seja plano e tenha fibras geométricas íntegras
(portanto, satisfaça às hipóteses de 3.1), e o morfismo finito
$f'':S'\longrightarrow S$ seja plano. Então $\operatorname{Pic}_{X/S}$ existe
e (com as notações introduzidas em [4], II, C 2) tem-se um isomorfismo canônico
\[
\operatorname{Pic}_{X/S}\xrightarrow{\sim}
\Pi_{S'/S}\operatorname{Pic}_{X/S'}.
\]

Para demonstrá-lo, estabelece-se primeiro um isomorfismo de funtores
\[
\operatorname{Pic}_{X/S}\xrightarrow{\sim}
\Pi_{S'/S}\operatorname{Pic}_{X/S'},
\]
utiliza-se 3.1, que implica que $\operatorname{Pic}_{X/S'}$ é representável, e
utiliza-se a estrutura assinalada no n\textsuperscript{o} 3 de
$\operatorname{Pic}_{X/S'}$ (que implica que toda parte finita de uma fibra de
$\operatorname{Pic}_{X/S'}$ sobre $S$ está contida em um aberto afim) para a
existência de $\Pi_{S'/S}\operatorname{Pic}_{X/S'}$.
\src{158}
Por exemplo, quando $X$ é um esquema soma de esquemas $X_i$ sobre $S$ que
satisfazem as condições de 3.1, o enunciado 6.1 reduz-se ao enunciado trivial
\[
\operatorname{Pic}_{X/S}\xrightarrow{\sim}
\prod_i\operatorname{Pic}_{X_i/S}.
\]

\result{COROLÁRIO}{6.2} Seja $f:X\longrightarrow S$ um morfismo projetivo e
plano com fibras geométricas localmente íntegras (por exemplo, um morfismo
projetivo e normal); então $\operatorname{Pic}_{X/S}$ existe.

Nesse caso, aliás, $S'$ é um recobrimento étale de $S$ (o que é verdadeiro
sempre que $f$ é separável, isto é, plano com fibras geométricas reduzidas), e
verifica-se que o teorema de estrutura enunciado no n\textsuperscript{o} 3 para
$\operatorname{Pic}_{X/S}$ continua válido graças à estrutura análoga de
$\operatorname{Pic}_{X/S'}$.

Por outro lado, um procedimento de descida fornece um teorema de existência
relativo, cujo alcance depende, aliás, da solução das questões de descida não
plana levantadas em [4], I, A (c), e do qual nos limitamos a explicitar aqui um
caso particular:

\result{PROPOSIÇÃO}{6.3} Seja $f:X\longrightarrow S$ um morfismo próprio, sejam
$X_1$ e $X_2$ dois subpré-esquemas de $X$ planos sobre $S$, definidos por dois
Ideais coerentes $\mathcal I_1$ e $\mathcal I_2$ tais que
$\mathcal I_1\cap\mathcal I_2=(0)$ e que
$\mathcal O_X/(\mathcal I_1+\mathcal I_2)$ seja plano sobre $S$ (isto é, o
subpré-esquema de $X$ que é o supremo de $X_1$ e $X_2$ é $X$, enquanto seu ínfimo $Z$
é plano sobre $S$).
\footnote{A última letra está impressa como ``$X$'', mas a condição
imediatamente anterior diz respeito à planura de
$\mathcal O_X/(\mathcal I_1+\mathcal I_2)$ sobre $S$ e, portanto, define $Z$
plano sobre $S$.} Suponhamos, além disso, que, para todo $s\in S$, os
homomorfismos
\[
k(s)\longrightarrow H^0(X_{i,s},\mathcal O_{X_{i,s}}),\qquad i=1,2,
\]
sejam bijetivos. Então o homomorfismo natural de funtores
\[
\operatorname{Pic}_{X/S}\longrightarrow
\operatorname{Pic}_{X_1/S}\times\operatorname{Pic}_{X_2/S}
\]
é representável por morfismos afins; portanto, se
$\operatorname{Pic}_{X_1/S}$ e $\operatorname{Pic}_{X_2/S}$ existem, o mesmo
ocorre com $\operatorname{Pic}_{X/S}$, e o morfismo canônico
\[
\operatorname{Pic}_{X/S}\longrightarrow
\operatorname{Pic}_{X_1/S}\times\operatorname{Pic}_{X_2/S}
\]
é afim.

Por descida fielmente plana, reduzimo-nos ao caso em que $Z$ admite uma seção
sobre $S$, definindo, portanto, seções de $X$, $X_1$, $X_2$ sobre $S$, o que
permite eliminar os automorfismos nas estruturas consideradas, como explicado
em 2.5. A demonstração consiste então em observar que o dado de um Módulo
invertível
\src{159}
«rigidificado» $\mathcal L$ sobre $X$ equivale ao dado de uma tripla
$(\mathcal L_1,\mathcal L_2,u)$, em que $\mathcal L_i$ é um Módulo invertível
«rigidificado» sobre $X_i$, e $u$ é um isomorfismo de
$\mathcal L_1|Z$ em $\mathcal L_2|Z$, compatível com as rigidificações.
Resta apenas verificar que, para $\mathcal L_1$ e $\mathcal L_2$ fixados, o
dado de $u$ se exprime por uma seção de um esquema conveniente sobre $S$, afim
sobre $S$, o que é fácil. Deduz-se facilmente de 6.3:

\result{COROLÁRIO}{6.4} Seja $X$ um esquema próprio e separável sobre um corpo
$k$, e sejam $X_i$ as componentes irredutíveis de $X$. Se existem os
$\operatorname{Pic}_{X_i/k}$, também existe $\operatorname{Pic}_{X/k}$, e o
morfismo canônico
\[
\operatorname{Pic}_{X/k}\longrightarrow\prod_i\operatorname{Pic}_{X_i/k}
\]
é afim.

Em combinação com 6.2, isso mostra, por exemplo, a existência de
$\operatorname{Pic}_{X/k}$ sempre que $X$ é um esquema projetivo e separável
sobre um corpo $k$. Quando $X$ deixa de ser separável sobre $k$, tem-se também
um resultado de redução, utilizando um raciocínio devido a Oort [11]. O método
aplica-se igualmente sobre um esquema de base arbitrário (caso útil, por
exemplo, para demonstrar, na exposição seguinte, o resultado de finitude
anunciado em 3.3). Para evitar um enunciado excessivamente técnico,
limitemo-nos ao caso de um corpo de base:

\result{PROPOSIÇÃO}{6.5} Seja $X$ um esquema próprio sobre um corpo $k$, e seja
$X_0$ um subesquema com o mesmo conjunto subjacente (portanto, definido por um
Ideal nilpotente sobre $X$). Então o morfismo funtorial
\[
\operatorname{Pic}_{X/k}\longrightarrow\operatorname{Pic}_{X_0/k}
\]
é representável por morfismos afins. Em particular, se
$\operatorname{Pic}_{X_0/k}$ existe, o mesmo ocorre com
$\operatorname{Pic}_{X/k}$, e o morfismo
$\operatorname{Pic}_{X/k}\longrightarrow\operatorname{Pic}_{X_0/k}$ é afim.

Combinando isso com 6.4, conclui-se facilmente:

\result{COROLÁRIO}{6.6} Seja $X$ um pré-esquema projetivo sobre um corpo $k$;
então $\operatorname{Pic}_{X/k}$ existe.

\result{OBSERVAÇÃO}{6.6} É extremamente plausível que, para todo esquema $X$
próprio sobre um corpo $k$, $\operatorname{Pic}_{X/k}$ exista.\footnote{A
errata de 1962 assinala que Murre acabava de resolver afirmativamente essa
questão.} Os resultados precedentes reduzem essa questão ao caso em que $k$ é
algebricamente fechado e $X$ é íntegro. Sabe-se, então, que existe um esquema
íntegro $X'$ projetivo sobre $k$ e um morfismo dominante
$g:X'\longrightarrow X$ (lema de Chow). Bastaria, portanto, mostrar que o
\src{160}
morfismo funtorial correspondente
\[
\operatorname{Pic}_{X/k}\longrightarrow\operatorname{Pic}_{X'/k}
\]
é representável (e, sem dúvida, representável por morfismos afins), pois se
sabe, por outro lado, que $\operatorname{Pic}_{X'/k}$ é representável. Isso
levanta questões de descida não plana ainda não resolvidas. Observe-se que,
quando se restringe a consideração do funtor $\operatorname{Pic}_{X/k}$ aos
pré-esquemas reduzidos (sendo $X$ próprio e íntegro sobre $k$ algebricamente
fechado), obtém-se de fato um funtor representável, como mostrou Chevalley [1]
no caso em que $X$ é normal, e Seshadri [12], por um método de descida, no caso
geral. Mas, com nossas notações, o esquema construído por esses autores não é
$\operatorname{Pic}_{X/k}$, mas
$(\operatorname{Pic}_{X/k})_{\mathrm{réd}}$, o esquema reduzido correspondente
a $\operatorname{Pic}_{X/k}$.

\result{OBSERVAÇÃO}{6.7} Seja, mais geralmente, $f:X'\longrightarrow X$ um
morfismo sobrejetivo de pré-esquemas próprios sobre $k$. Então, considerações de
descida não plana levam a conjecturar que
\[
\operatorname{Pic}_{X/k}\longrightarrow\operatorname{Pic}_{X'/k}
\]
é um morfismo afim,\footnote{A errata de 1962 assinala que essa questão se
resolve afirmativamente e remete ao último parágrafo dos Comentários à
exposição seguinte.} o que implicaria, em particular (dividindo pelas
componentes conexas dos elementos neutros), que o homomorfismo correspondente
sobre os grupos de Néron--Severi é injetivo módulo torsão. É o que se pode
verificar pela teoria das interseções quando $X$ e $X'$ são não singulares. A
resposta não parece conhecida em nenhum outro caso.

\result{OBSERVAÇÃO}{6.8} Ao contrário do que se poderia pensar, a consideração
dos esquemas de Picard de esquemas algébricos com elementos nilpotentes é útil,
e até indispensável, em diversas questões. Assim, quando $X$ é um esquema
projetivo, digamos, suave sobre $k$,\footnote{O texto impresso traz «regular»;
a errata prescreve «suave sobre $k$».} e $Y$ uma seção hiperplana, convém
considerar as «vizinhanças infinitesimais» de todas as ordens $X_n$ de $Y$ e os
esquemas de Picard $\operatorname{Pic}_{X_n/k}$; quando $X$ é irredutível de
dimensão $\geq4$ (resp. $\geq3$), o morfismo canônico
\[
\operatorname{Pic}_{X/k}\longrightarrow\operatorname{Pic}_{X_n/k}
\qquad(n\text{ grande})
\]
é um isomorfismo (resp. induz um isomorfismo para as imagens inversas dos
subgrupos de torsão de Néron--Severi), resultado que nos será útil no estudo
qualitativo dos esquemas de Picard na exposição seguinte. Do mesmo modo, a
consideração dos esquemas de Picard de certas curvas com elementos nilpotentes
e os teoremas fundamentais de geometria formal [3] permitem verificar, no caso
de características iguais, uma conjectura de Mumford, a saber, que, para todo
anel local noetheriano normal completo $A$ de dimensão 2, o grupo das classes
de divisores de $A$ pode ser considerado como o conjunto dos pontos racionais
sobre $k$ de um
\src{161}
grupo algébrico $G$ sobre o corpo residual $k$ ($G$ sendo, além disso,
determinado canonicamente uma vez escolhido um corpo de coeficientes em $A$).
No caso em que $A$ tem dimensão arbitrária, é plausível que exista um grupo
proalgébrico sobre $k$ que desempenhe o mesmo papel que $G$ acima e que se
construa, quando se pode ``dessingularizar'' $\operatorname{Spec}(A)$, como um
limite projetivo de esquemas de Picard de esquemas projetivos convenientes (com
elementos nilpotentes) sobre $k$.

\section*{BIBLIOGRAFIA}
\addcontentsline{toc}{section}{Bibliografia}

\begin{description}
\item[{[1]}] CHEVALLEY (Claude). -- Sur la théorie de Picard,
\emph{Amer. J. of Math.}, v. 82, 1960, p. 435--490.

\item[{[2]}] GROTHENDIECK (A.) e DIEUDONNÉ (J.). -- Éléments de géométrie
algébrique, Chapitre I et suivants. -- Paris, Presses universitaires de France,
1960--1961 (Institut des Hautes Études Scientifiques, Publications
mathématiques 4, 8, \ldots).

\item[{[3]}] GROTHENDIECK (Alexander). -- Géométrie formelle et géométrie
algébrique, \emph{Séminaire Bourbaki}, v. 11, 1958/59,
n\textsuperscript{o} 182, 28 p.

\item[{[4]}] GROTHENDIECK (Alexander). -- Techniques de descente et théorèmes
d'existence en géométrie algébrique, I--IV, \emph{Séminaire Bourbaki}, v. 12,
1959/60, n\textsuperscript{o} 190, 29 p. e n\textsuperscript{o} 195, 22 p.;
v. 13, 1960/61, n\textsuperscript{o} 212, 20 p. e
n\textsuperscript{o} 221, 28 p.

\item[{[5]}] GROTHENDIECK (Alexander). -- Techniques de construction en
géométrie analytique, I--X, \emph{Séminaire Cartan}, v. 13, 1960/61:
Déformation des structures analytiques complexes, n\textsuperscript{o} 7--17.

\item[{[6]}] GROTHENDIECK (Alexander). -- Séminaire de géométrie algébrique,
redigido por Lichtenbaum, Harvard University, 1961 (a aparecer).

\item[{[7]}] LANG (Serge). -- \emph{Abelian varieties}. -- Nova York,
Interscience Publishers, 1959 (Interscience Tracts in pure and applied
Mathematics, 7).

\item[{[8]}] MUMFORD (D.). -- An elementary theorem in geometric invariant
theory, \emph{Bull. Amer. math. Soc.}, v. 67, 1961, p. 483--486.

\item[{[9]}] MUMFORD (D.). -- The topology of normal singularities of an
algebraic surface and a criterion for simplicity. -- Paris, Presses
universitaires de France, 1961 (Institut des Hautes Études Scientifiques,
Publications mathématiques, 9).

\item[{[10]}] MUMFORD (D.) e TATE (J.). -- Séminaire de géométrie algébrique,
Harvard University, Spring Term 1962 (a aparecer).

\item[{[11]}] OORT (F.). -- Sur le schéma de Picard,
\emph{Bull. Soc. math. France}, v. 90, 1962 (a aparecer).

\item[{[12]}] SESHADRI (C. S.). -- Tese a aparecer em
\emph{Annali di Matematica}.
\end{description}

\par

\def\FgaUnitKey{e232}\clearpage
\chapter*{Errata da exposição n.\textsuperscript{o} 232}
\addcontentsline{toc}{chapter}{Errata da exposição n.\textsuperscript{o} 232}
\markboth{Errata da exposição 232}{Errata da exposição 232}
\sid{e232}
\src{303}

\noindent\textbf{7. TDTE V: Esquemas de Picard. Teoremas de existência.}

\begin{center}
[Seminário Bourbaki, vol. 14, 1961/62, n.\textsuperscript{o} 232, 19 p.]
\end{center}

\noindent\textbf{Página 7, observação 3.3.} --- A questão levantada aqui foi
resolvida afirmativamente por MUMFORD (ver os comentários à exposição seguinte).

\noindent\textbf{Página 13, observação 5.2.} --- Como assinalo no início da
exposição seguinte, a conjectura de existência formulada aqui é falsa; contudo,
MUMFORD consegue demonstrar, por seus métodos, um teorema ligeiramente mais
fraco.

\noindent\textbf{Página 14, linha 7.} --- Depois de «completo», acrescentar «com
corpo residual algebricamente fechado». O contraexemplo de MUMFORD que acabamos
de mencionar mostra, aliás, que essa restrição é indispensável.

\noindent\textbf{Página 17, observação 6.6.} --- Como assinalamos no início da
exposição seguinte, a questão levantada aqui acaba de ser resolvida
afirmativamente por MURRE.

\noindent\textbf{Página 18, observação 6.7.} --- A questão levantada aqui tem
resposta afirmativa (cf. o último parágrafo dos comentários à exposição
seguinte).

\noindent\textbf{Página 18, observação 6.8, 4\textsuperscript{a} linha.} --- Em vez de
«regular», leia-se «suave sobre $k$».

\par

\def\FgaUnitKey{236}\unit{Exposição 236 -- Os esquemas de Picard: propriedades gerais}
\sid{236}
\markboth{Exposição 236}{Exposição 236}
\src{221}
\seminares{14.º ano, 1961/62, n.º 236\\Maio de 1962}
\paperhead{TÉCNICA DE DESCIDA E TEOREMAS DE EXISTÊNCIA EM GEOMETRIA ALGÉBRICA\\
VI. OS ESQUEMAS DE PICARD: PROPRIEDADES GERAIS}%
  {por Alexander GROTHENDIECK}

\fsection{0}{Complementos à exposição anterior ([3], V)\textsuperscript{(1)}}

Houve alguns avanços nas questões de existência dos pré-esquemas de Picard
levantadas em [3], V:

\medskip
\noindent a. (Mumford). Não é verdade, em geral, que, quando
$f:X\longrightarrow S$ é um morfismo projetivo e separável (= plano com fibras
separáveis), exista o pré-esquema de Picard $\operatorname{Pic}_{X/S}$, mesmo
que as fibras de $f$ tenham dimensão 1 e que $S$ seja o espectro de um anel de
valorização discreta completo. Um exemplo é obtido tomando
$S=\operatorname{Spec}\mathbf R[[t]]$ e, para $X$, o subesquema de
$\mathbf P^2_S$ definido pela equação (onde $x,y,z$ são variáveis homogêneas)
\[
x^2+y^2=tz^2,
\]
que representa, portanto, uma cônica que degenera geometricamente em duas
retas concorrentes, embora a fibra especial sobre o corpo $\mathbf R$ seja,
não obstante, irredutível (ela é dada pela equação $x^2+y^2=0$ sobre
$\mathbf R$). Vê-se facilmente que, após a extensão étale
$S'\longrightarrow S$, com $S'=\operatorname{Spec}\mathbf C[[t]]$, existe o
pré-esquema de Picard de $X'/S'$, do qual se obtém uma descrição explícita como
união disjunta de cópias de $\widetilde S'$, onde $\widetilde S'$ é obtido de
$S'$ duplicando infinitas vezes a origem. Verifica-se facilmente que o dado de
descida sobre $\operatorname{Pic}_{X'/S'}$ para $S'\longrightarrow S$ (dado
aqui pela ação do grupo de Galois $\mathbf Z/2\mathbf Z$ de $S'$ sobre $S$)
não é efetivo: o grupo permuta entre si certos pontos duplicados (de modo que
há órbitas que não estão contidas em um aberto afim). Mumford pode demonstrar,
contudo, que, se $f:X\longrightarrow S$ é um morfismo projetivo separável tal
que, para todo $s\in S$, as componentes irredutíveis de $X_s$ são
geometricamente irredutíveis sobre $k(s)$, então
$\operatorname{Pic}_{X/S}$ existe; sua demonstração apoia-se em um refinamento
de seu teorema de passagem ao quociente, cf. [7]. Observe-se, por outro lado,
que continua sendo possível que, sem hipóteses sobre as componentes
irredutíveis das fibras $X_s$, o esquema $\operatorname{Pic}^{\tau}_{X/S}$
(que será introduzido mais adiante) exista mesmo assim.
\src{222}
\noindent b. (Murre). Seja $X$ um esquema próprio sobre um corpo. Então
$\operatorname{Pic}_{X/k}$ existe. A demonstração retoma em parte a
demonstração de Chevalley [1] e apoia-se essencialmente na estrutura de grupo
do funtor de Picard.

Para alguns complementos relativos à teoria dos esquemas de Picard,
especialmente em relação com os esquemas abelianos, será útil consultar [7].
Por fim, uma lacuna notável da presente exposição é a ausência de «critérios
de equivalência» que permitam comparar o esquema de Picard de um esquema
projetivo com os de suas seções hiperplanas; os teoremas-chave para desenvolver
tais critérios encontram-se em [5], aos quais devem ser acrescentados os
teoremas de existência dos esquemas de Picard.

\fsection{1}{Propriedades topológicas dos pré-esquemas em grupos comutativos}

Sejam primeiro $k$ um corpo e $G$ um pré-esquema em grupos sobre $k$. Como o
elemento neutro $e$, por ser racional sobre $k$, é necessariamente fechado,
segue-se imediatamente que a diagonal de $G\times_kG$ é fechada e, portanto,
que $G$ é separado: todo pré-esquema em grupos sobre um corpo é separado.
Denotaremos por $G^\circ$ a componente conexa do elemento neutro $e$. Como $e$
é racional sobre $k$, $G^\circ$ é, de fato, geometricamente conexa, isto é, a
formação de $G^\circ$ é compatível com a extensão do corpo de base. Segue-se
também que $G^\circ$ é estável por multiplicação como conjunto e, quando $G$ é
localmente noetheriano, de modo que $G^\circ$ é aberto, pode-se considerar
$G^\circ$ como um subgrupo aberto de $G$. Doravante, suporemos que $G$ é
localmente de tipo finito sobre $k$; então $G^\circ$ é geometricamente
irredutível e de tipo finito sobre $k$. Com efeito, pode-se supor que $k$ é
algebricamente fechado e, além disso, que $G$ é reduzido (pois
$G_{\mathrm{red}}$ será então um subgrupo de $G$, tendo em vista que
$G_{\mathrm{red}}\times_kG_{\mathrm{red}}$ será ainda reduzido), logo liso
sobre $k$ em um aberto não vazio e, portanto, em toda parte, como se vê
transladando esse aberto. Mas então $G$ é localmente irredutível, de modo que
suas componentes irredutíveis coincidem com suas componentes conexas, e
$G^\circ$ é irredutível. Seja então $U$ uma vizinhança afim de $e$ em
$G^\circ$; usando a irredutibilidade de $G^\circ$, vê-se imediatamente que
$U\cdot U=G^\circ$, o que prova que $G^\circ$ é quase-compacto e, portanto, de
tipo finito sobre $k$.

Para simplificar, suponhamos que $G$ seja comutativo. Para todo inteiro $n>0$,
seja $G^{(n)}$ a imagem inversa de $G^\circ$ pelo homomorfismo $\varphi_n$ de
elevação à $n$-ésima potência em $G$; assim, $G^{(n)}$ é um subgrupo aberto de
$G$.
Ponhamos:
\[
G^\tau=\bigcup_nG^{(n)}.
\]
\src{223}
\[
G^\sigma=\bigcup_{(n,p)=1}G^{(n)},
\qquad
G^\rho=\bigcup_hG^{(p^h)},
\]
onde $p$ é o expoente característico do corpo $k$. Obtêm-se, portanto,
subgrupos abertos de $G$ que satisfazem
\[
G^\sigma\cap G^\rho=G^\circ,
\qquad
G^\sigma G^\rho=G^\tau.
\]

\result{OBSERVAÇÃO}{} Pode-se construir o esquema em grupos quociente
$G/G^\circ=N$ (cf. [3], IV) e definir então $G^\tau$, $G^\sigma$ e
$G^\rho$ como as imagens inversas em $G$, respectivamente, do subgrupo de
torção de $N$, da soma de suas componentes $q$-primárias para os primos
$q\neq p$ e de sua componente $p$-primária.\footnote{O texto impresso troca
as duas últimas correspondências: atribui a $G^\sigma$ a componente
$p$-primária e a $G^\rho$ seu complemento primo com $p$. As definições
precedentes de $G^\sigma$ e $G^\rho$ impõem a ordem corrigida adotada aqui.}
Observe-se, a esse respeito, que $N$ é um esquema em grupos discreto separável
sobre $k$ e, portanto, uma vez escolhida uma clausura algébrica $\overline k$
de $k$, que dá lugar a um grupo de Galois $\pi$, identifica-se com um grupo
discreto ordinário sobre o qual $\pi$ age por automorfismos. É dessa maneira
que se pode interpretar de modo evidente a construção do subgrupo de torção e
a decomposição deste em suas componentes $q$-primárias. Quando $G$ é o esquema
de Picard de um esquema $X$ próprio sobre $k$, $N$ pode ser chamado
\emph{esquema de Néron--Severi} (reduzido) de $X$ sobre $k$. Quando
$G^\circ_{\mathrm{red}}$ é um subesquema em grupos de $G^\circ$, o que
ocorre, em particular, sempre que $k$ é perfeito ou que $G^\circ$ é próprio
sobre $k$ (por exemplo, quando $X$ é geometricamente normal), convém introduzir
também o quociente $N'=G/G^\circ_{\mathrm{red}}$, que tende a comportar-se
melhor do que $N$ quanto à especialização, isto é, quando $X$ varia em uma
família de esquemas algébricos.

Sejam agora $S$ um pré-esquema localmente noetheriano e $G$ um pré-esquema em
grupos sobre $S$, localmente de tipo finito sobre $S$. Não supomos que $G$ seja
de tipo finito sobre $S$, nem separado sobre $S$. Poremos então
\[
G^\circ=\bigcup_{s\in S}(G_s)^\circ,
\]
e, quando $G$ for comutativo,
\[
G^\tau=\bigcup_{s\in S}(G_s)^\tau,
\qquad
G^\sigma=\bigcup_{s\in S}(G_s)^\sigma,
\qquad
G^\rho=\bigcup_{s\in S}(G_s)^\rho.
\]
\src{224}
Trata-se de subconjuntos de $G$, estáveis pela multiplicação em $G$, o que
evidentemente não implica que possam ser definidos mediante subpré-esquemas em
grupos de $G$. Em particular, parece não existir, em geral, um subpré-esquema
em grupos de $G$ cujo conjunto subjacente seja $G^\circ$. Naturalmente, se um
desses conjuntos for aberto, então, munido da estrutura induzida, será um
subpré-esquema aberto em grupos de $G$. Veremos que isso sempre ocorre com
$G^\tau$; assim, do ponto de vista dos funtores representáveis, e em particular
do ponto de vista das «especializações», a equivalência numérica comporta-se
de maneira mais satisfatória do que a equivalência algébrica. Eis as principais
propriedades gerais dos conjuntos que acabamos de definir:

\result{TEOREMA}{1.1} $G^\circ$, $G^\tau$, $G^\sigma$ e $G^\rho$ são
localmente construtíveis. Além disso:

\noindent (i) $G^\circ$ é quase-compacto sobre $S$. Quando os $G_s^\circ$ são
próprios e $G$ é separado sobre $S$, então $G^\circ$ é próprio sobre $S$ e,
portanto, fechado em $G$.

\noindent (ii) $G^\tau$ é aberto. Se $G^\circ$ é fechado, o mesmo vale para
$G^\tau$.

\noindent (iii) Se $G^\circ$ é fechado, o mesmo vale para $G^\sigma$, desde
que se esteja em característica igual, isto é, que todos os corpos residuais de
$S$ tenham a mesma característica. Se $G^\circ$ é fechado e
$G\longrightarrow S$ é universalmente aberto nos pontos de $G^\sigma$ (cf.
Corolário 1.5 abaixo), então $G^\sigma\longrightarrow S$ é universalmente
aberto.

\noindent (iv) Se $G^\circ$ é fechado, o mesmo vale para $G^\rho$. Suponhamos
que se esteja em característica igual e que, para todo inteiro $n>0$ tal que
$(n,p)=1$, o homomorfismo de elevação à $n$-ésima potência em $G$ seja aberto;
então $G^\rho$ é aberto.

Demos algumas indicações sobre a demonstração. O fato de $G^\circ$ ser
localmente construtível e quase-compacto sobre $S$ está contido no lema
seguinte:

\result{LEMA}{1.2} Suponhamos que $S\neq\varnothing$. Então existem um aberto
não vazio $U$ em $S$, um esquema em grupos $H$ de tipo finito sobre $U$ com
fibras conexas e, por fim, um homomorfismo de esquemas em grupos
$H\longrightarrow G|U$ que é uma imersão aberta cuja imagem é
$G^\circ|U$.

Para demonstrar o lema, pode-se supor que $S$ seja irredutível; seja $\eta$ seu
ponto genérico, e façamos a mudança de base
$S'=\operatorname{Spec}(\mathcal O_{S,\eta})\longrightarrow S$. Obtém-se um
esquema em grupos $G'$ sobre um anel artiniano local
$\mathcal O_{S,\eta}=A$, no qual há um subesquema aberto em grupos
$G'{}^\circ$ de tipo finito sobre $A$, como se disse acima. Este provém,
portanto, de um esquema em grupos $H$ de tipo finito sobre uma vizinhança
aberta $U$ de $\eta$, e a imersão canônica
$G'{}^\circ\longrightarrow G'$ provém de uma imersão
\src{225}
aberta $H\longrightarrow G|U$, que será um homomorfismo de esquemas em grupos
para $U$ suficientemente pequeno. Como as fibras de $H$ são conexas quando se
toma $U$ suficientemente pequeno, e têm todas a mesma dimensão, igual à das
fibras de $G$ para $U$ suficientemente pequeno, segue-se que, para todo
$s\in U$, a imagem de $H_s$ em $G_s$ é exatamente $G_s^\circ$ (após reduzir
novamente $U$), o que demonstra 1.2.

A segunda afirmação de (i) está contida no lema seguinte (que se aplica a uma
vizinhança aberta quase-compacta de $G^\circ$ em $G$):

\result{LEMA}{1.3} Sejam $X$ um pré-esquema separado de tipo finito sobre um
$S$ localmente noetheriano, $g$ uma seção de $X$ sobre $S$ e $X^\circ$ a união
das componentes conexas dos $g(s)$ nos $X_s$. Seja $s\in S$ tal que
$X_s^\circ$ seja próprio sobre $k(s)$. Então existe uma vizinhança aberta $U$
de $s$ tal que $X^\circ|U$ seja próprio sobre $U$ e, a fortiori, fechado em
$X|U$.

Por descida fielmente plana da base, reduzimo-nos ao caso em que $S$ é o
espectro de um anel local completo e $s$ é seu ponto fechado. Aplicando [2],
III, 5.5.1, vê-se que $X$ se decompõe na união disjunta de dois abertos $X'$ e
$X''$, o primeiro próprio sobre $S$ e tal que $X'_s=X_s^\circ$. Isso nos reduz
ao caso em que $X=X'$, isto é, em que $X$ é próprio sobre $S$. Nesse caso,
conclui-se por uma demonstração padrão que utiliza o critério valuativo para
que uma parte seja própria (omitido em [2], cap. II).

Demonstremos que $G^\tau$ é aberto ou, o que equivale ao mesmo, tendo em vista
que a formação de $G^\tau$ (assim como a de $G^\circ$, $G^\sigma$ e $G^\rho$)
comuta com a extensão da base: para toda seção $g$ de $G$ sobre $S$,
$g^{-1}(G^\tau)$ é aberto. Isso significa duas coisas:

\medskip
\noindent a. Seja $y\in S$ tal que $g(y)\in G_y^\tau$. Então, para todo
$y'\in\overline{\{y\}}$ suficientemente próximo de $y$, tem-se
$g(y')\in G_{y'}^\tau$.

\medskip
\noindent b. Seja $y'\in S$ tal que $g(y')\in G_{y'}^\tau$. Então, para toda
generalização $y$ de $y'$, tem-se $g(y)\in G_y^\tau$.

Para (a), observe-se que existe um inteiro $n>0$ tal que
$g^n(y)\in G_y^\circ$; como $G^\circ$ é construtível,
$(g^n)^{-1}(G^\circ)$ é construtível e, portanto, tem-se
$g^n(y')\in G_{y'}^\circ$ para $y'\in\overline{\{y\}}$ suficientemente
próximo de $y$. Para (b), observe-se que as potências
$g(y')^n=g^n(y')$ permanecem em um aberto quase-compacto de $G_{y'}$ (pois
estão contidas em um número finito de classes módulo $G_{y'}^\circ$); logo,
existe um aberto quase-compacto $U$ em $G$ que contém os $g^n(y')$ e, portanto,
também suas generalizações $g^n(y)$. Assim, as potências de $g(y)$ permanecem
em um aberto quase-compacto de $G_y$, o que implica facilmente que
$g(y)\in G_y^\tau$.
\src{226}
Suponhamos $G^\circ$ fechado; demonstremos que o mesmo vale para $G^\tau$.
Como já se sabe que $G^\tau$ é aberto e, portanto, localmente construtível,
resta demonstrar que é estável por especialização, o que decorre do fato de ser
uma união de fechados, a saber, as imagens inversas, pelos homomorfismos
$\varphi_n$ de elevação à $n$-ésima potência, do fechado $G^\circ$.

O mesmo argumento demonstrará que $G^\sigma$ e $G^\rho$ são fechados se
$G^\circ$ o for (sob a condição, no primeiro caso, de se estar em
característica igual), uma vez demonstrado que $G^\sigma$ e $G^\rho$ são
localmente construtíveis. Ora, para $x\in G^\tau$, seja $\nu(x)$ o menor
inteiro $n>0$ tal que o homomorfismo $\varphi_n$ de elevação à $n$-ésima
potência envie $x$ a $G^\circ$. Assim, $G^\sigma$ (respectivamente,
$G^\rho$) é formado pelos $x\in G$\footnote{O texto impresso traz aqui
«$X$», que não está definido nesta passagem; os quatro conjuntos considerados
são subconjuntos de $G$.} tais que $\nu(x)$ seja primo com $p$
(respectivamente, uma potência de $p$), e nossa afirmação de construtibilidade
decorre então da seguinte afirmação mais precisa:

\result{LEMA}{1.4} A função $\nu$ sobre $G^\tau$ é localmente construtível.

Isso significa, com efeito, que, para todo inteiro $n>0$, o conjunto dos
$x\in G^\tau$ tais que $\nu(x)=n$ é localmente construtível; ora, ele é a
diferença entre $\varphi_n^{-1}(G^\circ)$ e a união dos
$\varphi_d^{-1}(G^\circ)$, onde $d$ percorre os divisores próprios de $n$.
Como $G^\circ$ é localmente construtível, o mesmo vale para os
$\varphi_d^{-1}(G^\circ)$ e, portanto, também para a diferença anterior.

Suponhamos $G^\circ$ fechado e $G\longrightarrow S$ universalmente aberto nos
pontos de $G^\sigma$; demonstremos que $G^\sigma\longrightarrow S$ é aberto,
isto é, que transforma uma vizinhança em $G^\sigma$ de um $x\in G^\sigma$
em uma vizinhança de $y=f(x)$. Como $G^\sigma$ é localmente construtível,
reduzimo-nos a demonstrar que, para toda generalização $y'$ de $y$, existe uma
generalização $x'$ de $x$ em $G^\sigma$ sobre $y'$. Isso nos reduz, por
mudança de base, ao caso em que $S$ é o espectro de um anel de valorização
discreta e em que $y,y'$ são, respectivamente, o ponto fechado e o ponto
genérico. Usando o fato de que $G\longrightarrow S$ é aberto em $x$ (logo,
existe uma generalização $x_1$ de $x$ em $G$ sobre $y'$), pode-se supor ainda
que existe uma seção $g$ de $G$ sobre $S$ tal que $x=g(y)$ (após fazer outra
mudança de base). Se $k(y')$ tem característica 0, basta tomar qualquer
generalização $x'$ de $x$ em $G$ sobre $y'$; ela pertence a $G^\tau$ porque
$G^\tau$ é aberto e, portanto, a $G^\sigma$ porque
$G_{y'}^\sigma=G_{y'}^\tau$. Se a característica de $k(y')$ é $p>0$,
ponhamos
\[
\nu(g(y'))=p^hm\qquad\text{com }(m,p)=1 ;
\]
\src{227}
sejam $a$ e $b$ dois inteiros tais que $ap^h+bm=1$; ponhamos
\[
g_1=g^{ap^h},\qquad g_2=g^{bm},\qquad\text{logo}\qquad g=g_1g_2 ;
\]
portanto, por construção, tem-se $g_1(y')\in G^\sigma$ e
$g_2(y')\in G^\rho$. Como $G^\circ$ é fechado, segue-se que
$g_1(y)\in G^\sigma$ e $g_2(y)\in G^\rho$; daí, em virtude de
\[
g(y)=g_1(y)g_2(y)\in G^\sigma,
\]
também $g_2(y)\in G^\sigma$ e, portanto,
\[
g_2(y)\in G_y^\sigma\cap G_y^\rho=G_y^\circ.
\]
Ora, da hipótese e do fato de que $S$ é o espectro de um anel de valorização
discreta, resulta que $G-(G_y-G_y^\circ)$ é um aberto de $G$ sobre o qual
$G\longrightarrow S$ induz um morfismo aberto; logo, por todo ponto de
$G_y^\circ$ passa uma «quase-seção». Por conseguinte, após fazer ainda uma
extensão da base $S$, pode-se supor que existe uma seção $g'_2$ de $G^\circ$
sobre $S$ tal que $g'_2(y)=g_2(y)$. Ponhamos $g'=g_1g'_2$; então, por
construção, $g'(y)=g(y)=x$, e
\[
g'(y')=g_1(y')g'_2(y')\in G^\sigma,
\]
logo $g'(y')=x'$\footnote{O texto impresso traz aqui «$g(y')=x'$»; a seção
construída na frase precedente e cuja fibra genérica está sendo calculada é
$g'$.} é uma generalização de $x$ em $G^\sigma$ sobre $y'$, o que demonstra
(iii).

Demonstremos, por fim, a última afirmação (iv). Reduzimo-nos a provar que, se
$x\in G^\rho$, então toda generalização $x'$ de $x$ pertence a $G^\rho$.
Pode-se supor (após tomar as imagens por $\varphi_{p^h}$ para um $h$
conveniente)\footnote{O texto impresso traz $\varphi_{n^h}$; para
$x\in G^\rho$, a definição dada na p.~223 fornece uma potência do expoente
característico $p$ que envia $x$ a $G^\circ$.} que se tem até mesmo
$x\in G^\circ$. Então, para todo inteiro $n$ primo com a característica, $x$
pertence à imagem de $\varphi_n$ (pois a elevação à $n$-ésima potência em um
grupo conexo de tipo finito sobre um corpo de característica prima com $n$ é
sobrejetiva). Como $\varphi_n$ é aberta, segue-se que $x'$ também pertence à
imagem de $\varphi_n$. Mais precisamente, seja $U$ um aberto quase-compacto de
$G^\tau$ que contém $G_y^\circ$; então se tem $x'\in\varphi_n(U)$ para todo
$n$ primo com $p$. Tomando para $n$ um múltiplo comum dos fatores primos com
$p$ dos $\nu(z)$, para $z\in U$, obtém-se $x'\in G^\rho$.

A afirmação (1.1), (iii), completa-se assim:

\result{COROLÁRIO}{1.5} Seja $n>1$ um inteiro tal que o homomorfismo de
elevação à $n$-ésima potência $\varphi_n:G\longrightarrow G$ seja
universalmente aberto (por exemplo, étale); ponhamos
\[
G^{(n)}=\varphi_n^{-1}(G^\circ),
\]
e suponhamos que, em cada fibra geométrica $G_{\bar s}$, o subconjunto
\src{228}
$\bigcup_{h\geq0}|{}_{n^h}G_{\bar s}|$ seja denso em
$\bigcup_{h\geq0}G_{\bar s}^{(n^h)}$. Então $G\longrightarrow S$ é universalmente aberto nos
pontos dos $G^{(n^h)}$. Em particular, se os $G_s^\circ$ não contêm componente
aditiva e se, para todo inteiro $n>1$, o homomorfismo $\varphi_n$ é
universalmente aberto nos pontos de característica residual prima com
$n$\footnote{O texto impresso traz «prima com $p$». Nesta proposição
quantificada por $n$, o critério de etalidade e o argumento pelos núcleos
${}_nG$ exigem característica residual prima com $n$.}, então
$G\longrightarrow S$ é universalmente aberto nos pontos de $G^\sigma$ e,
portanto, por (1.1), (iii), se $G^\circ$ é fechado,
$G^\sigma\longrightarrow S$ é universalmente aberto. Nesses casos,
$s\longmapsto\dim G_s$ é uma função localmente constante sobre $S$.

Com efeito, da hipótese resulta que o núcleo ${}_nG$ de $\varphi_n$ é
universalmente aberto sobre $S$; logo, $G\longrightarrow S$ é universalmente
aberto nos pontos de ${}_nG$ e, portanto (substituindo $n$ por $n^h$), nos
pontos dos ${}_{n^h}G$. A hipótese de densidade diz exatamente que, fibra a fibra, os pontos de ordem
uma potência de $n$ são densos na união dos $G^{(n^h)}$. Para a afirmação
particular, a ausência de componentes aditivas implica essa densidade para os
inteiros $n$ invertíveis no corpo residual; daí resulta facilmente que
$G\longrightarrow S$ é universalmente aberto em todos os pontos dessa união,
em particular ao longo da seção neutra, de onde decorre facilmente que
$s\longmapsto\dim G_s$ é uma função localmente constante.

\result{OBSERVAÇÃO}{1.6} Recordemos que um morfismo $X\longrightarrow S$ é
dito universalmente aberto se transforma todo aberto em um aberto e conserva
essa louvável propriedade após toda mudança de base $S'\longrightarrow S$.
Isso significa também (quando $S$ é localmente noetheriano e
$X\longrightarrow S$ é localmente de tipo finito) que toda componente
irredutível de $X$ domina $S$ e que essa propriedade é preservada após toda
mudança de base $S'\longrightarrow S$. Em ambas as afirmações, basta, aliás
(sob as hipóteses de finitude acima), testar mudanças de base
$S'\longrightarrow S$, em que $S'$ é o espectro de um anel de valorização
discreta (completo, com corpo residual algebricamente fechado, se assim se
quiser~...). A definição estende-se de maneira evidente ao caso de um
subconjunto $Z$ de $X$ (como o subconjunto $G^\sigma$ de $G$). É um fenômeno
perfeitamente normal, mesmo partindo de um morfismo projetivo liso
$X\longrightarrow S$ com fibras geométricas conexas (por exemplo, o produto
fibrado da família modular de curvas elípticas consigo mesma sobre
$S=\operatorname{Spec}\mathbf C[j]$), que $\operatorname{Pic}_{X/S}$ não
seja universalmente aberto sobre $S$, isto é, que possa haver componentes
irredutíveis de $\operatorname{Pic}_{X/S}$ situadas inteiramente sobre um
único ponto de $S$: isso está ligado ao fato de que o posto do grupo de
Néron--Severi das fibras de $X/S$ pode apresentar saltos para cima (fenômenos
de «multiplicação complexa»). Em contrapartida, (1.5) assegura-nos que, nos
casos favoráveis, $\operatorname{Pic}_{X/S}^\sigma$ (e, ao que parece, na
maioria das vezes até mesmo $\operatorname{Pic}_{X/S}^\tau$) é
universalmente aberto sobre $S$.
\src{229}
Por fim, eis um caso útil em que, excepcionalmente, $G^\circ$ consente em ser
aberto:

\result{COROLÁRIO}{1.7} Suponhamos que $G\longrightarrow S$ seja
universalmente aberto nos pontos de $G^\circ$ (cf. (1.5)) e que as fibras
$G_s$ sejam separáveis e, portanto, lisas sobre $k(s)$ (esta última condição é automaticamente satisfeita em
característica residual zero, em virtude de um resultado de Cartier). Então
$G^\circ$ é aberto em $G$. Se, além disso, $S$ é reduzido, $G^\circ$ é liso
(em particular, plano) sobre $S$.

A primeira afirmação pode, aliás, ser precisada observando-se que, se, para um
$s\in S$ dado, $G\longrightarrow S$ é universalmente aberto nos pontos de
$G_s^\circ$ e $G_s^\circ$ é separável sobre $k(s)$, então $G^\circ$ é uma
vizinhança de $G_s^\circ$ (e, aliás, a hipótese feita em $s$ continuará
verificada nos pontos vizinhos). Esse enunciado, independente de toda estrutura
de grupos, será demonstrado em [2], IV, §~7. A última afirmação de (1.7), também
independente de toda estrutura de grupos e de toda questão de componente
conexa, é um caso particular de um critério de planicidade dado em [2], IV, §~5,
que implica, mais geralmente:

\result{COROLÁRIO}{1.8} Seja $U$ um subconjunto aberto de uma fibra $G_s$ tal
que $G\longrightarrow S$ seja universalmente aberto nos pontos de $U$ (cf.
(1.5)). Se $G_s$ é separável sobre $k(s)$ e $S$ é reduzido em $s$, então $G$ é
plano sobre $S$ nos pontos de $U$ (e, portanto, neste caso, liso sobre $S$ nos
pontos de $U$).

\result{OBSERVAÇÕES}{1.9} Ignoro se, quando $G$ é separado sobre $S$,
$G^\circ$ ou $G^\tau$ é sempre fechado em $G$; isso parece pouco provável.
Encontram-se, em todo caso, contraexemplos evidentes se se abandona a hipótese
de separação, por exemplo tomando a reta afim cuja origem é duplicada uma
infinidade enumerável de vezes, obtendo-se assim um esquema em grupos $G$ sobre
a reta afim cujas fibras se reduzem todas ao grupo unidade, salvo uma cuja
fibra é $\mathbf Z$. Esse esquema em grupos é um subesquema aberto em grupos,
a saber, o fecho da seção unidade, no pré-esquema de Picard do $S$-esquema $X$
correspondente a uma família de cônicas que degenera em duas retas concorrentes.

Além disso, mesmo partindo de um esquema em grupos finito e plano $G$ sobre o
espectro $S$ de um anel de valorização discreta $V$ e, por conseguinte, se
$G^\circ$ se reduz à seção unidade, sendo, portanto, fechado, várias
conclusões tornam-se falsas quando se abandonam certas hipóteses. Suponhamos
que $V$ seja de característica igual $p>0$ e seja $G$ o núcleo do homomorfismo
$\mathbf G_a\longrightarrow\mathbf G_a$ definido pelo homomorfismo de
funtores
\[
f\longmapsto f^p-tf,
\]
onde $t$ é um uniformizante de $V$.
\src{230}
(Lembremos que, por definição, o «grupo aditivo» $\mathbf G_a$ sobre $S$
representa o funtor $S'\longmapsto\Gamma(S',\mathcal O_{S'})$.) Então, como
$S$-esquema,
\[
G=\operatorname{Spec}V[x]/(x^p-tx)
\]
é a união de $p$ «retas» concorrentes, representando um grupo
$\mathbf Z/p\mathbf Z$ que degenera em um grupo infinitesimal de tipo aditivo.
Neste exemplo, $G^\circ=G^\sigma$ (e é uma das $p$ retas), portanto não é
aberto em $G$. No caso de característica mista, de característica residual
$p>0$, partamos do esquema em grupos $H=\boldsymbol\mu_p$, núcleo da aplicação
de elevação à $p$-ésima potência em $\mathbf G_m$, portanto
\[
H=\operatorname{Spec}V[x]/(x^p-1) ;
\]
é ainda a união de $p$ retas concorrentes, representando um grupo
$\mathbf Z/p\mathbf Z$ (em característica zero) que degenera em um grupo
infinitesimal de tipo multiplicativo em característica $p$. Seja $H'$ o
«esquema em grupos constante» definido pelo grupo finito ordinário
$\mathbf Z/p\mathbf Z$, portanto a soma disjunta, indexada por
$\mathbf Z/p\mathbf Z$, de cópias de $S$, ou ainda
$\operatorname{Spec}V^{\mathbf Z/p\mathbf Z}$. Então $G=H\times_S H'$ descreve
um grupo $(\mathbf Z/p\mathbf Z)^2$ em característica zero, que degenera em um
grupo infinitesimal vezes $\mathbf Z/p\mathbf Z$ em característica $p$. Aqui
$G^\rho$ é a união da seção unidade e da fibra especial, portanto $G^\rho$ não
é aberto, ao passo que $G^\sigma$ é a união da seção unidade e da fibra
genérica, portanto não é fechado, contrariamente ao que se afirma nos casos de
característica igual em (1.1), (iii) e (iv). Naturalmente, esses fenômenos estão
ligados à característica $p>0$. Os resultados precedentes dão, com efeito:

\result{COROLÁRIO}{1.10} Suponhamos que as características residuais de $S$
sejam todas nulas, portanto $G^\tau=G^\sigma$ e $G^\circ=G^\rho$. Então
$G^\tau=G^\sigma$ é aberto, e até aberto e fechado se $G$ é separado sobre $S$
e se os $G_s^\circ$ são próprios; sob essa mesma hipótese, $G^\circ=G^\rho$ é
próprio sobre $S$, portanto fechado em $G$. Suponhamos, por fim, que, para todo
inteiro $n>1$, o homomorfismo de elevação à $n$-ésima potência em $G$ seja
universalmente aberto; então $G^\circ$ é aberto; e, se além disso os
$G_s^\circ$ não têm componente aditiva (por exemplo, se os $G_s^\circ$ são
próprios, como acima), então $G^\tau\longrightarrow S$ é universalmente aberto,
e até liso se $S$ é reduzido.

Assinalemos, por fim, o seguinte resultado fácil:

\result{PROPOSIÇÃO}{1.11} Existe uma parte aberta $U$ de $S$ tal que o conjunto
$G'$ dos pontos de $G$ nos quais $G$ é liso (resp. plano) sobre $S$ seja o
conjunto subjacente a um subesquema em grupos aberto induzido de $G|U$. Além
disso, toda seção de $G$ sobre $U$ é uma seção de $G'$ sobre $U$.
\src{231}
\result{COROLÁRIO}{1.12} Se $G$ é liso (resp. plano) sobre $S$ nos pontos da
seção unidade, também o é nos pontos de toda seção de $G$
sobre $S$ e em todos os pontos de $G^\circ$. Se, além disso, para todo inteiro
$n>0$, o homomorfismo de elevação à $n$-ésima potência
$\varphi_n:G\longrightarrow G$ é étale nos pontos de característica prima com
$n$, então $G$ é liso (resp. plano) sobre $S$ em todos os pontos de $G^\sigma$.

\fsection{2}{Aplicação às propriedades locais dos esquemas de Picard}

\result{TEOREMA}{2.1}

\noindent(i) Seja $f:X\longrightarrow S$ um morfismo próprio e liso, e
suponhamos que $\operatorname{Pic}_{X/S}$ exista (por exemplo, que $f$ seja um
morfismo projetivo). Então $\operatorname{Pic}_{X/S}$ é separado sobre $S$ e,
para toda parte fechada $Z$ de $\operatorname{Pic}_{X/S}$ que seja de tipo
finito sobre $S$, $Z$ é próprio sobre $S$.

\noindent(ii) Seja $X$ um pré-esquema sobre um corpo $k$, próprio e
geometricamente normal. Então $\operatorname{Pic}_{X/k}^\circ$\footnote{O texto
impresso traz $\operatorname{Pic}_{X/S}^\circ$, embora $S$ não intervenha neste
enunciado sobre o corpo $k$.} é próprio sobre $k$.

\begin{proof}
(i) Os critérios valuativos ([2], II, parágrafo 7) reduzem-nos a provar o
seguinte: se $S$ é o espectro de um anel de valorização discreta completo e $U$
o aberto reduzido ao ponto genérico de $S$, então toda seção racional de
$\operatorname{Pic}_{X/S}$ sobre $S$, isto é, toda seção sobre $U$, prolonga-se
de maneira única a uma seção sobre $S$. Tendo em conta a definição de
$\operatorname{Pic}_{X/S}$, isso equivale ao seguinte enunciado: para todo
módulo invertível $\mathcal L$ sobre $V=f^{-1}(U)$, existe um módulo invertível
sobre $X$ que prolonga $\mathcal L$, único a menos de isomorfismo. Ora, isso
decorre facilmente da descrição dos módulos invertíveis sobre $V$, resp. $X$,
em termos de classes de divisores de Cartier, tendo em conta que os anéis
locais de $X$ são regulares ($X$ sendo liso sobre $S$ regular), portanto
fatoriais pelo teorema de Auslander, o que implica que todo divisor sobre
$X$\footnote{O texto impresso traz aqui «sobre $S$»; a fatorialidade invocada é
a dos anéis locais de $X$.} é um divisor de Cartier. Com efeito, todo divisor
sobre $V$ pode ser prolongado a um divisor sobre $X$ tomando-se o seu fecho.
\end{proof}

\result{OBSERVAÇÃO}{2.2} A demonstração permanece válida supondo-se apenas que
$f$ é plano e que suas fibras $X_s$ são localmente interseções completas e
lisas sobre $k(s)$ em codimensão $\leq2$, tendo em conta o seguinte resultado
demonstrado em [5]: um anel local noetheriano de interseção completa que é
regular em codimensão $\leq3$ é fatorial («conjectura de Samuel»). Note-se que o
\src{232}
resultado torna-se falso se se substitui «codimensão $\leq2$» por «codimensão
$\leq1$», isto é, pela hipótese «normal», como se pode constatar no exemplo de
uma família de quádricas não singulares que degenera em um cone quadrático.

(ii) Utilizando o lema de Chow, reduzimo-nos ao caso em que $X$ é projetivo,
portanto imerso em $\mathbf P_k^n$; podemos supor $X$ conexo. Se $\dim X=1$,
então $X$ é liso sobre $k$, e aplica-se (i). Se
$\dim X\geq2$, pode-se utilizar uma variante dos «critérios de equivalência»
conhecidos, que implica existir um número finito de curvas $Y_i$ lisas em $X$
(obtidas como interseções de $X$ com subespaços lineares convenientes de
$\mathbf P_k^n$), tais que
\[
\operatorname{Pic}_{X/k}^\tau\longrightarrow
\prod_i\operatorname{Pic}_{Y_i/k}^\tau
\]
seja um monomorfismo e induza, a fortiori, um monomorfismo nas componentes
conexas. Como a componente conexa do segundo membro é própria sobre $k$ pelo
que precede, e como se trata de um homomorfismo de esquemas em grupos, que será
necessariamente uma imersão fechada, segue-se que
$\operatorname{Pic}_{X/k}^\circ$ também é próprio sobre $k$. Pode-se evitar o
recurso aos delicados critérios de equivalência utilizando-se a estrutura dos
grupos algébricos comutativos sobre um corpo algebricamente fechado (devida a
Chevalley--Borel); reduzimo-nos a provar que todo morfismo da reta afim privada
da origem\footnote{O texto impresso traz «da reta afim». A errata vinculada,
p.~303, prescreve «reta afim privada da origem».} em
$\operatorname{Pic}_{X/k}^\tau$ é constante, o que equivale a dizer que todo
módulo invertível sobre $X[t,t^{-1}]$\footnote{O texto impresso traz $X[t]$; a
errata vinculada, p.~303, prescreve $X[t,t^{-1}]$.} provém de um módulo
invertível sobre $X$, resultado bem conhecido e muito elementar, que nem sequer
utiliza o fato de $X$ ser próprio sobre $k$ (a hipótese de $X$ ser normal
permitindo reduzir-nos imediatamente ao caso em que $X$ é regular).

\result{COROLÁRIO}{2.3} Seja $f:X\longrightarrow S$ um morfismo próprio e
normal (isto é, plano e com fibras geométricas normais), e suponhamos que
$\operatorname{Pic}_{X/S}$ exista; então $\operatorname{Pic}_{X/S}^\circ$ é
próprio sobre $S$, portanto fechado em $\operatorname{Pic}_{X/S}$; além disso,
$\operatorname{Pic}_{X/S}^\tau$ e $\operatorname{Pic}_{X/S}^\rho$ também são
fechados, assim como $\operatorname{Pic}_{X/S}^\sigma$ em «característica
igual».

Aplicam-se (1.1) e (2.1), (ii).

\result{COROLÁRIO}{2.4} Seja $f:X\longrightarrow S$ um morfismo próprio e liso,
tal que $\operatorname{Pic}_{X/S}$ exista e seja uma soma disjunta de esquemas
$P^{(i)}$ de tipo finito sobre $S$ (cf. [3], V, 4.1). Então cada $P^{(i)}$ é
próprio sobre $S$.

Isso decorre de (2.1), (i). Como se assinalou em (2.2), o resultado pode ser
generalizado fazendo-se hipóteses menos restritivas sobre as fibras de $f$, mas
torna-se falso se nos limitamos a supor $f$ normal. Nesse caso, ignoro, aliás,
se $\operatorname{Pic}_{X/S}^\tau$ é, não obstante, próprio sobre $S$, mesmo
admitindo que seja de tipo finito sobre $S$.
\src{233}
\result{TEOREMA}{2.5} Seja $f:X\longrightarrow S$ um morfismo próprio e plano,
tal que $\operatorname{Pic}_{X/S}$ exista, e, para todo inteiro $n$, seja
$\varphi_n$ o homomorfismo de elevação à $n$-ésima potência nesse pré-esquema
em grupos. Então $\varphi_n$ é étale em todo ponto
$x\in\operatorname{Pic}_{X/S}$\footnote{O texto impresso traz «$x\in X$»;
$\varphi_n$ é o endomorfismo de potência do pré-esquema de Picard.} de
característica residual prima com $n$.

Pela caracterização infinitesimal dos morfismos étales, essa afirmação equivale
à seguinte:

\result{LEMA}{2.6} Suponhamos que $S$ seja o espectro de um anel artiniano local
$A$ cujo ideal maximal $\mathfrak m$ tenha potência $(\nu+1)$-ésima nula, e
sejam
\[
A_{\nu-1}=A/\mathfrak m^\nu,
\qquad X_{\nu-1}=X\otimes_A A_{\nu-1},
\]
$\mathcal L$ um módulo invertível sobre $X$ e $\mathcal L'_{\nu-1}$ um módulo
invertível sobre $X_{\nu-1}$ cuja potência tensorial $n$-ésima seja isomorfa a
\[
\mathcal L_{\nu-1}=\mathcal L\otimes_A A_{\nu-1}.
\]
Então existe um módulo invertível $\mathcal L'$ sobre $X$ cuja potência
tensorial $n$-ésima é isomorfa a $\mathcal L$ (se $n$ é primo com a
característica residual de $k=k(A)$).

Demonstremos o lema. Ponhamos
\[
V=\mathfrak m^\nu=\mathfrak m^\nu/\mathfrak m^{\nu+1} ;
\]
este é um espaço vetorial sobre $k=k(A)$. Começamos por prolongar
$\mathcal L'_{\nu-1}$ a um módulo invertível qualquer $\mathcal L'$ sobre $X$.
A obstrução a fazê-lo encontra-se em
\[
H^2(X_0,\mathcal O_{X_0})\otimes_k V,
\]
mas, em virtude da hipótese feita sobre $\mathcal L'_{\nu-1}$ e do fato de
$\mathcal L_{\nu-1}$ poder ser prolongado, vê-se que o produto dessa obstrução
por $n$ é nulo, portanto ela própria é nula, pois $n$ é primo com a
característica. Além disso, a arbitrariedade na extensão efetuada encontra-se
em
\[
H^1(X_0,\mathcal O_{X_0})\otimes_k V,
\]
e o desvio $\xi$ de $\mathcal L'^{\otimes n}$ em relação a $\mathcal L$
encontra-se no mesmo módulo; se se procura corrigir $\mathcal L'$ de modo a
tornar nulo esse desvio, reduzimo-nos a encontrar um $\eta$ no referido módulo
tal que $n\eta=\xi$. Isso é novamente possível porque $n$ é primo com a
característica.

\result{COROLÁRIO}{} Sob as condições de (2.5), suponhamos, além disso, que os
esquemas de Picard das fibras $X_s$ não contenham componente aditiva (por
exemplo, que os $X_s$ sejam geometricamente normais, cf. (2.1), (ii)). Então
$\operatorname{Pic}_{X/S}\longrightarrow S$ é universalmente aberto nos pontos
de $\operatorname{Pic}_{X/S}^\circ$. Se
$\operatorname{Pic}_{X/S}^\circ$ é fechado (por exemplo, se os $X_s$ são
geometricamente normais, cf. (2.3)), então
$\operatorname{Pic}_{X/S}^\sigma$ é ele próprio universalmente aberto sobre
$S$. Por fim, no caso de característica igual,
$\operatorname{Pic}_{X/S}^\rho\longrightarrow S$ é universalmente aberto.

Aplicam-se (1.5) e (1.1).
\src{234}
\result{COROLÁRIO}{2.7} Seja $f:X\longrightarrow S$\footnote{O texto impresso
traz «$f:X\longrightarrow Y$», embora tanto o pré-esquema de Picard quanto a
função seguinte tenham base $S$.} um morfismo próprio e plano, tal que
$\operatorname{Pic}_{X/S}$ exista. Então a função
\[
s\longmapsto\dim\operatorname{Pic}_{X_s/k(s)}
\]
sobre $S$ é semicontínua superiormente (isto é, pode dar saltos para cima, mas
não para baixo), e é até contínua (isto é, localmente constante) se os
$\operatorname{Pic}_{X_s/k(s)}$ não contêm componente aditiva.

A primeira afirmação é trivial, ou quase, para todo pré-esquema em grupos
localmente de tipo finito sobre uma base localmente noetheriana, pois basta
examinar ao longo da seção unidade. A segunda afirmação decorre de (2.5).

\result{OBSERVAÇÃO}{2.8} Sejam $s,s'\in S$ tais que $s$ seja uma especialização
de $s'$. Então (2.7) equivale a uma desigualdade (resp. igualdade) entre as
dimensões dos esquemas de Picard de $X_{s'}$ e de sua «especialização» $X_s$.
Aliás, Serre havia observado, ainda antes da construção dos esquemas de Picard,
que a invariância da dimensão dos Picard dos $X_s$, no caso de um morfismo liso
$f:X\longrightarrow S$, era uma consequência formal da teoria da
especialização do grupo fundamental ([4], X) e das relações clássicas de Kummer
entre os pontos de ordem finita na variedade de Picard e o grupo fundamental
abelianizado ([4], XI). Se designamos, de modo geral, por $\alpha,\mu,\lambda$
as dimensões das partes abeliana, multiplicativa e aditiva de
$\operatorname{Pic}_{X_s/k(s)}^\circ$, e definimos do mesmo modo
$\alpha',\mu',\lambda'$, as relações conhecidas exprimem-se pelas seguintes
desigualdades:
\[
\tag{*}\alpha+\mu+\lambda\geq\alpha'+\mu'+\lambda'
\]
(válida contanto que $\operatorname{Pic}_{X/S}$ exista, portanto provavelmente
em todos os casos), desigualdade que, para $\lambda=\lambda'=0$, se reduz a uma
igualdade, válida sob as mesmas hipóteses de existência:
\[
\alpha+\mu=\alpha'+\mu' ;
\]
por outro lado, tem-se
\[
\tag{**}2\alpha+\mu\leq2\alpha'+\mu'
\]
\src{235}
pelo argumento de Serre, se os $X_s$ são separáveis (sem sequer supor a
existência de $\operatorname{Pic}_{X/S}^\tau$), ou se os
${}_n\operatorname{Pic}_{X/S}$ (núcleos de $\varphi_n$ no funtor de Picard) são
separados sobre o aberto de $S$ onde $n$ é invertível (tendo em conta que são étales sobre esse aberto graças a (2.5)).
Tender-se-ia a conjecturar que (*) é uma igualdade em todos os casos, ou pelo
menos se os $X_s$ são separáveis, e que se têm as desigualdades
\[
\tag{***}\alpha\leq\alpha',\qquad\lambda\geq\lambda',
\]
que deveriam ser válidas sempre que se tenha um pré-esquema em grupos,
localmente de tipo finito sobre $S$ localmente noetheriano, no qual a dimensão
das fibras seja constante (ver (1.3) para um resultado positivo nessa direção).

\result{OBSERVAÇÃO}{2.9} Em todos os casos conhecidos,
$\operatorname{Pic}_{X/S}^\tau$ é universalmente aberto sobre $S$, mas sem
dúvida não se deve tirar daí ilusões excessivas, mesmo que
$f:X\longrightarrow S$ seja liso; em todo caso, Mumford
construiu um exemplo (é verdade que com $S$ não reduzido, sendo de fato $S$ o
espectro de um anel artiniano) no qual $\operatorname{Pic}_{X/S}^\tau$ não é
plano sobre $S$, fazendo variar infinitesimalmente a superfície de Igusa. O
ponto considerado por Mumford encontra-se, aliás, em
$\operatorname{Pic}_{X/S}^\rho$, e continua possível (para
$f:X\longrightarrow S$ liso) que $\operatorname{Pic}_{X/S}^\tau$ seja plano
sobre $S$ nos pontos de $\operatorname{Pic}_{X/S}^\sigma$; o conferencista
duvida, contudo, que isso seja sempre assim, mesmo restringindo-se aos pontos de
$\operatorname{Pic}_{X/S}^\circ$ e supondo que $S$ seja o espectro de um anel
de valorização discreta. A questão está, aliás, ligada ao estudo dos pontos
fixos de um esquema abeliano sob um grupo finito de automorfismos, situação
para a qual parecem faltar exemplos.

Parece que, mesmo limitando-nos a $f:X\longrightarrow S$ liso e projetivo, os
resultados de regularidade local sobre $\operatorname{Pic}_{X/S}$ enunciados no
presente número e as conjecturas assinaladas em (2.8) esgotam aproximadamente o
que se pode dizer a esse respeito sem hipóteses mais particulares sobre a
natureza das fibras de $f$. Lembremos, contudo, que, se as fibras geométricas
de $\operatorname{Pic}_{X/S}$ são reduzidas e sem componente aditiva, então
decorre de (1.8) e (2.5) que $\operatorname{Pic}_{X/S}$ é liso sobre $S$ nos
pontos de $\operatorname{Pic}_{X/S}^\sigma$ quando $S$ é reduzido; esse
resultado vale até mesmo, se $f:X\longrightarrow S$ é normal, sem hipótese
sobre $S$, como veremos em (3.5). Assinalemos, a esse respeito:

\result{PROPOSIÇÃO}{2.10}

\noindent(i) Se $\operatorname{Pic}_{X/S}$ é liso (resp. plano) sobre $S$ nos
pontos da seção unidade, também o é em todos os pontos de
$\operatorname{Pic}_{X/S}^\sigma$ e nos pontos de toda seção de
$\operatorname{Pic}_{X/S}$ sobre $S$.
\src{236}
\noindent(ii) Seja $s\in S$ tal que
\[
H^2(X_s,\mathcal O_{X_s})=0,
\]
então existe uma vizinhança aberta $U$ de $s$ tal que
$\operatorname{Pic}_{X/S}|U$ seja liso sobre $U$.

\noindent(iii) Seja $X$ um esquema próprio sobre um corpo; então tem-se
\[
\dim\operatorname{Pic}_{X/k}\leq\dim H^1(X,\mathcal O_X),
\]
com igualdade se e somente se $\operatorname{Pic}_{X/k}$ é liso sobre $k$;
isso ocorre sempre se $k$ tem característica zero.

(i) decorre de (2.5) e (1.12), (ii) do critério infinitesimal de lisura e de um
cálculo de obstruções bem conhecido, tendo em conta que (pelo «teorema de
semicontinuidade») a hipótese feita em $s$ continuará verificada nos pontos
vizinhos. Por fim, (iii) decorre do fato de que $H^1(X,\mathcal O_X)$ é
isomorfo ao espaço tangente de Zariski no elemento neutro de
$\operatorname{Pic}_{X/k}$; a última afirmação é um caso particular de um
teorema de Cartier segundo o qual um «grupo formal» em característica zero é
formalmente liso sobre $k$.

\fsection{3}{O subesquema abeliano canônico de
\texorpdfstring{$\operatorname{Pic}_{X/S}$}{Pic(X/S)} e o esquema de Albanese}

\result{PROPOSIÇÃO}{3.1} Sejam $k$ um corpo e $G$ um esquema em grupos de tipo
finito sobre $k$, comutativo e «sem componente aditiva». Então
$G_{\mathrm{red}}^\circ$ é separável sobre $k$ e, por conseguinte, é um
subesquema em grupos liso sobre $k$.

Sendo isso trivial se $k$ é perfeito, em particular para $G_{\bar k}$, onde
$\bar k$ é o fecho algébrico de $k$, basta ver que
$(G_{\bar k}^\circ)_{\mathrm{red}}$ provém de um subesquema de $G$. Ora, da
hipótese sobre $G_{\bar k}$ (ausência de componente aditiva)\footnote{O texto
impresso traz uma construção equivalente a «Ora, a hipótese ... resulta»; a
errata vinculada, p.~303, insere a preposição necessária para «Ora, da hipótese
... resulta».} resulta facilmente que existe um inteiro $m$ tal que
$(G_{\bar k}^\circ)_{\mathrm{red}}$ seja a imagem esquemática do homomorfismo de
elevação à $m$-ésima potência em $G_{\bar k}$. Como esse último homomorfismo
provém do homomorfismo análogo para $G$, a imagem esquemática deste último
responde à questão.

\result{COROLÁRIO}{3.2} Seja $X$ um esquema geometricamente normal e próprio sobre $k$ tal que
$\operatorname{Pic}_{X/k}$ exista; então existe um subesquema abeliano $A$ de
$\operatorname{Pic}_{X/k}$ cujo conjunto subjacente é
$\operatorname{Pic}_{X/k}^\circ$.

Com efeito, em virtude de (2.1), (ii), sendo
$\operatorname{Pic}_{X/k}^\circ$ próprio sobre $k$, ele satisfaz a condição de
(3.1).
\src{237}
O resultado precedente mostra, portanto, que, em certos casos, a clássica
«variedade de Picard» (que, na teoria atual, seria denotada por
$(\operatorname{Pic}_{X/k})_{\mathrm{red}}^\circ$) «é definida sobre $k$»,
sem supor que o corpo $k$ seja perfeito.

Seja agora $f:X\longrightarrow S$ um esquema relativo próprio e plano, com
\[
\mathcal O_S\xrightarrow{\sim}f_*(\mathcal O_X)
\]
para simplificar, tal que $\operatorname{Pic}_{X/S}$ exista e que
$\operatorname{Pic}_{X/S}^\circ$ seja próprio sobre $S$. Suporemos, além
disso, para (3.3), (ii), que existe um aberto de $\operatorname{Pic}_{X/S}$
contendo $\operatorname{Pic}_{X/S}^\circ$ que seja quase-projetivo sobre $S$;
essa condição é satisfeita, como vimos, se $f$ é projetivo e tem fibras
geométricas separáveis e irredutíveis. Recordemos que se chama esquema
abeliano sobre $S$ um esquema em grupos sobre $S$, próprio e liso sobre
$S$, com fibras geométricas conexas. Propomo-nos
examinar se existe um subesquema em grupos $A$ de
$\operatorname{Pic}_{X/S}$ que seja um esquema abeliano e cujo conjunto
subjacente seja $\operatorname{Pic}_{X/S}^\circ$. Acabamos de ver que sempre
existe um se $S$ é o espectro de um corpo. Eis o que se sabe dizer no caso
geral considerado aqui:

\result{TEOREMA}{3.3} Sob as condições precedentes:

\noindent(i) Se existe um subesquema abeliano de
$\operatorname{Pic}_{X/S}$ cujo conjunto subjacente é
$\operatorname{Pic}_{X/S}^\circ$, ele é único. Sua formação é, por
conseguinte, compatível com as mudanças de base.

\noindent(ii) Para que exista tal subesquema abeliano, basta que exista um
após toda mudança de base $S'\longrightarrow S$, em que $S'$ é artiniano
local; se $S$ é o espectro de um anel local, basta até mesmo testar com os
$S'=\operatorname{Spec}(A_n)$, em que $A_n=A/\mathfrak m^{n+1}$. Se $S$ é
reduzido, basta igualmente testar com as mudanças de base
$S'\longrightarrow S$, em que $S'$ é o espectro de um anel de valorização
discreta (completo com corpo residual algebricamente fechado, se assim se
desejar).

\noindent(iii) Suponhamos que $A$ exista e seja
$B=\operatorname{Alb}^0(X/S)$ o esquema abeliano dual (isto é,
$B=\operatorname{Pic}_{A/S}^\circ$ [7]). Pode-se então construir
canonicamente um espaço homogêneo principal
$P=\operatorname{Alb}^1(X/S)$ sob $B$ e um $S$-morfismo
$X\longrightarrow P$ que seja universal para os $S$-morfismos de $X$ em
esquemas para-abelianos (isto é, espaços homogêneos principais sob esquemas
abelianos). A formação de $\operatorname{Alb}^0(X/S)$,
$\operatorname{Alb}^1(X/S)$ e
$X\longrightarrow\operatorname{Alb}^1(X/S)$ comuta com a mudança de base.

Esbocemos a demonstração:

(i) é uma propriedade geral de rigidez dos subesquemas abelianos dos esquemas
em grupos comutativos: se dois desses subesquemas coincidem como conjuntos em
um ponto $s\in S$, eles coincidem sobre toda a componente
\src{238}
conexa de $s$ [7]. (Este resultado generaliza um teorema clássico de Chow.)

(ii) Utilizando os esquemas de Hilbert, vê-se que o funtor que, a todo $S'$
sobre $S$, associa o conjunto (reduzido a 1 ou 0 elementos) dos subesquemas
abelianos canônicos de $(\operatorname{Pic}_{X/S})\times_S S'$ é
representável por um esquema $T$ de tipo finito sobre $S$. Em virtude de (i),
$T\longrightarrow S$ é um monomorfismo e, em virtude de (3.2), é
sobrejetivo. Dizer que existe um subesquema abeliano canônico de
$\operatorname{Pic}_{X/S}$ significa que $T$ tem uma seção sobre $S$, ou ainda
que $T\longrightarrow S$ é um isomorfismo. Isso equivale também a dizer que
$T\longrightarrow S$ é étale ou, no caso em que $S$ é reduzido, que
$T\longrightarrow S$ é próprio. Daí resulta imediatamente (ii).

(iii) Utilizando simplesmente a definição de $\operatorname{Pic}_{X/S}$,
constata-se que, para todo esquema abeliano $C$ sobre $S$, o dado de um
$S$-morfismo de $X$ em um espaço homogêneo principal sob $C$ equivale ao dado
de um homomorfismo de grupos
$C'\longrightarrow\operatorname{Pic}_{X/S}$, em que $C'$ é o esquema abeliano
dual de $C$. Ora, se existe o subesquema abeliano canônico $A$ de
$\operatorname{Pic}_{X/S}$, esses homomorfismos fatoram-se necessariamente por
$A$ (como se pode ver ainda utilizando os pontos de ordem finita). Daí resulta
imediatamente (iii).

\result{OBSERVAÇÕES}{3.4} Denotaremos por
$\operatorname{Pic}_{X/S}^{\circ\circ}$ o subesquema abeliano canônico de
$\operatorname{Pic}_{X/S}$, caso exista. Infelizmente, nem sempre é esse o
caso, como se pode ver fazendo variar infinitesimalmente a superfície de Igusa
(pela deformação modular de primeira ordem). É possível, contudo, que
$\operatorname{Pic}_{X/S}^{\circ\circ}$ exista ao menos se $S$ é reduzido ou,
o que equivale ao mesmo segundo (ii), se $S$ é o espectro de um anel de
valorização discreta. Sejam então $X_0,X_1$ as fibras especial e
genérica\footnote{O francês mantido traz no singular «la fibre spéciale et
générique». Os testemunhos inglês “the special and generic fibers” e espanhol
«las fibras especial y genérica» exigem o plural; o ramo corrente usa «as
fibras» e o diplomático preserva «a fibra».} de $X$, e seja
\[
A_1=\operatorname{Pic}_{X_1/K}^\circ,
\]
em que $K$ é o corpo de frações do anel de valorização $V$. Em virtude de
Koizumi, existe um esquema abeliano $A$ sobre $S$, essencialmente único, cuja
fibra genérica é $A_1$, e vê-se facilmente, como em (2.1), (i) (supondo
doravante $X$ liso sobre $S$), que o morfismo identidade de $A_1$ se prolonga em um
morfismo
\[
A\longrightarrow\operatorname{Pic}_{X/S}.
\]
Daí um homomorfismo
\[
\tag{*}A_0\longrightarrow\operatorname{Pic}_{X_0/k}^{\circ\circ},
\]
que se demonstra facilmente ser um homomorfismo sobrejetivo cujo núcleo é um
$p$-grupo finito, em que $p$ é o expoente característico do corpo residual $k$
\src{239}
(ainda utilizando os pontos de ordem finita). Isto posto, as condições
seguintes sobre $X/S$ são equivalentes:

\medskip
\noindent a. O homomorfismo precedente (*) é um isomorfismo (o que Shimura
exprimiria dizendo que a formação da «variedade de Picard» é «compatível com a
especialização»).

\noindent b. $\operatorname{Pic}_{X/S}^{\circ\circ}$ existe (então não será
outro senão $A$).

\noindent c. (Para lembrar) Os
$\operatorname{Pic}_{X_n/S}^{\circ\circ}$ existem.

Segundo a observação feita sobre o núcleo de (*), a condição (a) é satisfeita
se a característica residual é nula, mas esse resultado será generalizado de
modo notável em (3.5).

Naturalmente, se $\operatorname{Pic}_{X/S}$ é liso sobre $S$ nos pontos de
$\operatorname{Pic}_{X/S}^\circ$, este último é
aberto em $\operatorname{Pic}_{X/S}$ (cf. (1.7)) e, portanto, munido da
estrutura induzida, é um subesquema abeliano de
$\operatorname{Pic}_{X/S}$, logo igual a
$\operatorname{Pic}_{X/S}^{\circ\circ}$, que existe nesse caso. Mas há muito
mais:

\result{TEOREMA}{3.5} Sob as condições de (3.3), seja $s\in S$ tal que
$\operatorname{Pic}_{X_s/k(s)}$ seja liso sobre $k(s)$ (ou, o que equivale ao
mesmo,
\[
\dim\operatorname{Pic}_{X_s/k(s)}=\dim H^1(X_s,\mathcal O_{X_s})).
\]
Então existe uma vizinhança aberta $U$ de $s$ tal que
$\operatorname{Pic}_{X/S}$ seja liso sobre $S$ nos pontos de
$\operatorname{Pic}_{X/S}^\circ|U$\footnote{O francês mantido traz a
construção impossível «en les points». O ramo corrente verte a leitura
corrigida «aux points» por «nos pontos»; o diplomático preserva literalmente
«em os pontos».}, que é, portanto, um subesquema abeliano aberto em
$\operatorname{Pic}_{X/S}|U$. A fortiori,
$\operatorname{Pic}_{X|U/U}^{\circ\circ}$ existe.

Indiquemos o princípio da demonstração. O que precede reduz-nos ao caso em que
$S$ é o espectro de um anel artiniano local $A$, e raciocina-se por indução
sobre a ordem infinitesimal de $A$. Pode-se, portanto, supor que
$\operatorname{Pic}_{X_n/A_n}^\circ$ seja liso sobre $A_n$ e limitar-se a
provar que $\operatorname{Pic}_{X_{n+1}/A_{n+1}}^\circ$ é liso sobre
$A_{n+1}$. Constata-se que, para isso, basta construir um esquema abeliano
$P_{n+1}$ sobre $A_{n+1}$ que prolongue
\[
P_n=\operatorname{Pic}_{X_n/A_n}^\circ,
\]
e, ao mesmo tempo, um Módulo invertível $\mathcal L_{n+1}$ sobre
$X_{n+1}\times_{A_{n+1}}P_{n+1}$ que prolongue o Módulo invertível
$\mathcal L_n$ sobre $X_n\times_{A_n}P_n$ que intervém na definição do esquema
de Picard $\operatorname{Pic}_{X_n/A_n}$ como solução de um problema universal.
(N.B. --- Podemos supor que $X$ seja munido de uma seção sobre $S$~...) Para
realizar a construção, deve-se utilizar o seguinte resultado-chave: todo
esquema abeliano definido sobre um quociente de anel artiniano local pode ser
prolongado (em outros termos, o «esquema formal dos Módulos» absoluto ([3], II)
para um esquema abeliano sobre um corpo algebricamente fechado é liso sobre o
anel dos
\src{240}
vetores de Witt sobre $k$); esse próprio resultado é obtido utilizando as
propriedades formais gerais da obstrução ao levantamento e as operações do
grupo. Admitido esse resultado, começa-se por prolongar $P_n$ de uma maneira
qualquer a $P_{n+1}$; encontra-se então uma obstrução ao levantamento de
$\mathcal L_n$, obstrução que pertence a
\[
H^2(X_0\times P_0,\mathcal O_{X_0\times P_0})\otimes_k V
\qquad(\text{onde }V=\mathfrak m^{n+1}/\mathfrak m^{n+2}),
\]
e até mesmo ao subespaço
\[
H^1(X_0,\mathcal O_{X_0})\otimes
H^1(P_0,\mathcal O_{P_0})\otimes_k V
\]
(tendo em conta que a restrição de $\mathcal L_n$ aos dois fatores $X_n$ e
$P_n$ é trivial). Ora, este último espaço não é outro senão
\[
H^1(P_0,\mathcal T_{P_0/k})\otimes V,
\]
em que $\mathcal T_{P_0/k}$ é o feixe tangente a $P_0/k$; portanto, é também o
espaço que exprime a indeterminação que havia no levantamento de $P_n$ a
$P_{n+1}$ ([3], II). Por conseguinte, pode-se corrigir esse levantamento (de
uma e uma só maneira, como convém) de modo a anular a obstrução ao levantamento
de $\mathcal L_n$.

\result{COROLÁRIO}{3.6} Sob as condições de (3.5),
$R^1f_*(\mathcal O_X)$ é um Módulo localmente livre sobre $S$ em uma
vizinhança de $s$, e sua formação comuta com a mudança de base.

Com efeito, esse Módulo não é outro senão o Módulo tangente a
$\operatorname{Pic}_{X/S}$ ao longo da seção unidade.

\result{OBSERVAÇÃO}{3.7} Utilizando o mesmo argumento que em (3.5), pode-se
mostrar que, se $S'$ é um subesquema de $S$ definido por um Ideal coerente
nilpotente e se se supõe apenas que $\operatorname{Pic}_{X'/S'}^\circ$ existe e
tem fibras lisas, então $\operatorname{Pic}_{X/S}^\circ$ existe
necessariamente e é um esquema abeliano sobre $S$. Isso permite construir
esquemas de Picard em certos casos apesar da ausência de uma hipótese
projetiva; por exemplo, o esquema abeliano dual de um esquema abeliano sobre um
anel artiniano sempre existe. Por outro lado, utilizando (3.6) no caso em que
$X$ é um esquema abeliano sobre $S$ e utilizando a estrutura conhecida de
$H^*(X_s,\mathcal O_{X_s})$ como álgebra exterior sobre
$H^1(X_s,\mathcal O_{X_s})$ (Rosenlicht--Serre), conclui-se que, para todo
$i$, $R^if_*(\mathcal O_X)$ é localmente livre e, mais precisamente, isomorfo
à $i$-ésima potência exterior de $R^1f_*(\mathcal O_X)$.

\result{OBSERVAÇÃO}{3.8} No caso de um morfismo projetivo liso
$f:X\longrightarrow S$, com $S$ reduzido e de características residuais
nulas, o resultado (3.6) era conhecido por métodos transcendentes, como
consequência da teoria de Hodge. De fato, todos os
$R^pf_*(\Omega^q_{X/S})$ são então localmente livres. Há, em contrapartida,
contraexemplos
\src{241}
no caso de características mistas para «$R^1f_*(\mathcal O_X)$ localmente
livre», com o auxílio das variedades de Serre ([8], p.~50). Parece não haver
nenhum contraexemplo em característica igual.

\result{COROLÁRIO}{3.8} Sob as condições de (3.5),
$\operatorname{Pic}_{X/S}|U$ é liso sobre $U$ em todos os pontos de
$\operatorname{Pic}_{X/S}^\sigma|U$.

Aplica-se (2.10), (i). Em particular, tendo em conta (2.10), (iii):

\result{COROLÁRIO}{3.9} Suponhamos que $S$ tenha características residuais
nulas. Então $\operatorname{Pic}_{X/S}^\tau$ é liso sobre $S$.

\result{OBSERVAÇÃO}{3.10} Deduz-se daí, por exemplo, se
$\operatorname{Pic}_{X/S}^\tau$ é também próprio sobre $S$, que o grupo de torção
dos grupos de Néron--Severi das fibras geométricas de $f$ permanece constante
sobre toda componente conexa de $S$ (o que, aliás, é evidente por métodos
transcendentes quando $f$ é liso e projetivo). Assinalemos que a utilização
direta de (2.5) permite mostrar, mais geralmente, que, se
$\operatorname{Pic}_{X/S}^\tau$ é próprio sobre $S$ (por exemplo,
$f:X\longrightarrow S$ liso e projetivo) e se $q$ é um número primo distinto
das características residuais de $S$, então a componente $q$-primária dos
grupos de torção de Néron--Severi das fibras geométricas de $X/S$ permanece
constante sobre toda componente conexa de $S$. Contudo, já não ocorre o mesmo
em característica $p>0$ para a componente $p$-primária da torção. Continua
possível, porém, que o posto sobre $k(s)$ de
\[
\operatorname{Pic}_{X_s/k(s)}^\tau/
\operatorname{Pic}_{X_s/k(s)}^{\circ\circ}=T_{X_s/k(s)}
\]
permaneça localmente constante; quando $S$ é reduzido, pode-se mostrar que isso
equivale a que $\operatorname{Pic}_{X/S}^{\circ\circ}$ exista e que
$\operatorname{Pic}_{X/S}^\tau$ seja plano sobre $S$, e basta testar a questão
para $S$ espectro de um anel de valorização discreta. Foi isso que verifiquei
nos poucos exemplos que examinei; mas, como o enunciado correspondente com $S$
artiniano é falso (cf. observações (2.9) e (3.4)), sem dúvida não convém
alimentar ilusões excessivas.

\fsection{4}{O teorema de finitude para o esquema de Picard}

Seja $f:X\longrightarrow S$ um morfismo projetivo e plano tal que
$\operatorname{Pic}_{X/S}$ exista; então a consideração dos «polinômios de
Hilbert» $Q\in\mathbf Q[t]$ permite decompor $\operatorname{Pic}_{X/S}$ em uma
soma de abertos $\operatorname{Pic}_{X/S}^Q$. Quando não se fazem hipóteses
adicionais, que assegurem, por exemplo, que $\operatorname{Pic}_{X/S}$ seja
separado sobre $S$, não será verdade, em geral, que esses abertos sejam de tipo
finito sobre $S$; obtém-se um contraexemplo quando $X$ é «uma cônica que
degenera em duas retas». É
\src{242}
possível, contudo, que assim seja se $f$ é separável e tem fibras
geometricamente irredutíveis. A questão está ligada à de saber se
$\operatorname{Pic}_{X/S}^\tau$ é de tipo finito sobre $S$, o que talvez seja
verdade sem hipótese alguma sobre as fibras de $X/S$. Quando
$f:X\longrightarrow S$ é liso, note-se que
$\operatorname{Pic}_{X/S}^\tau$ está contido em um dos
$\operatorname{Pic}_{X/S}^Q$ (graças ao fato de que, sobre uma variedade
projetiva não singular, a equivalência «por torção» é mais fina que (na
realidade, «idêntica à», segundo Matsusaka) a equivalência numérica para os
divisores), portanto é de tipo finito sobre $S$ se os
$\operatorname{Pic}_{X/S}^Q$ o são. Note-se que as questões de finitude que
acabamos de levantar conservam um sentido evidente mesmo sem supor a existência
de $\operatorname{Pic}_{X/S}$, pois se exprimem dizendo que certas famílias de
Módulos invertíveis são «limitadas» no sentido de [3], IV: consideremos, para
todo corpo algebricamente fechado $k$, os subesquemas íntegros de
$\mathbf P_k^r$, de dimensão $n$ e de grau $d$, e os Módulos invertíveis sobre
tais pré-esquemas, que tenham polinômio de Hilbert $Q$ (onde $r,n,d,Q$ são
dados); trata-se de mostrar que a família desses Módulos (considerados como
Módulos coerentes sobre os $\mathbf P_k^r$) é limitada, isto é, pode ser
parametrizada por um esquema de tipo finito sobre $\mathbf Z$.

Utilizando o método de Matsusaka [6], uma demonstração bastante técnica (que
recorre aos «critérios de equivalência» sob uma forma conveniente) permite
responder afirmativamente quando nos restringimos às subvariedades não
singulares de $\mathbf P_k^r$. Mais precisamente, obtém-se o resultado seguinte:

\result{TEOREMA}{4.1} Seja $f:X\longrightarrow S$\footnote{O texto impresso
traz «Sejam: $X\longrightarrow S$»; o morfismo utilizado a seguir é $f$.} um
morfismo liso e projetivo com fibras geometricamente conexas, sendo $S$
noetheriano, $\mathcal O_X(1)$ um Módulo muito amplo sobre $X$ relativamente a
$S$, $E$ um subconjunto de $\operatorname{Pic}_{X/S}$, correspondente a uma
família $(\mathcal L_i)$ de Módulos invertíveis sobre as fibras geométricas
$\overline X_{s_i}$ de $X/S$, $D_i$ um divisor (não necessariamente efetivo)
sobre $\overline X_{s_i}$ que define $\mathcal L_i$,
\[
a_n^{(i)}t^n+\cdots+a_0^{(i)}
\]
o polinômio de Hilbert de $\mathcal L_i$, $\xi=\xi_i$ um divisor que define
$\mathcal O_X(1)$, isto é, uma seção hiperplana. As condições seguintes são
equivalentes:

\noindent a. $E$\footnote{O texto impresso traz aqui «$Q$»; o subconjunto
introduzido no enunciado, cuja quase-compacidade exprime que a família é
limitada, é $E$.} é quase-compacto, isto é, está contido em um aberto de tipo
finito sobre $S$, isto é, a família dos $(\mathcal L_i)$ é limitada.

\noindent b. $E$ está contido na reunião de um número finito dos conjuntos
$\operatorname{Pic}_{X/S}^Q$, $Q\in\mathbf Q[t]$, isto é, o conjunto dos
polinômios de Hilbert dos $\mathcal L_i$ é finito.

\noindent b'. (Se as fibras de $X/S$ têm todas a mesma dimensão $n$.) Os dois
coeficientes $a_{n-1}^{(i)}$ e $a_{n-2}^{(i)}$ dos polinômios de Hilbert dos
$\mathcal L_i$ permanecem em um conjunto finito.
\src{243}
\result{COROLÁRIO}{4.2} Seja $f:X\longrightarrow S$ um morfismo liso e
projetivo com fibras geometricamente conexas; então os esquemas
$\operatorname{Pic}_{X/S}^Q$ e $\operatorname{Pic}_{X/S}^\tau$
($Q\in\mathbf Q[t]$) são projetivos sobre $S$.

Como são de tipo finito sobre $S$ segundo (4.1), pode-se aplicar (2.4).

\section*{BIBLIOGRAFIA}
\addcontentsline{toc}{section}{Bibliografia}

\begin{description}
\item[{[1]}] CHEVALLEY (Claude). -- Sur la théorie de Picard,
\emph{Amer. J. of Math.}, t. 82, 1960, p. 435--490.

\item[{[2]}] GROTHENDIECK (A.) e DIEUDONNÉ (J.). -- Éléments de géométrie
algébrique, Capítulo I e seguintes. -- Paris, Presses universitaires de France,
1960--1961 (Institut des Hautes Études Scientifiques, 4, 8, 11, \ldots).

\item[{[3]}] GROTHENDIECK (Alexander). -- Technique de descente et théorèmes
d'existence en géométrie algébrique, I--V, \emph{Séminaire Bourbaki}, t. 12,
1959/60, n\textsuperscript{o} 190, 29 p. e n\textsuperscript{o} 195, 22 p.;
t. 13, 1960/61, n\textsuperscript{o} 212, 20 p. e
n\textsuperscript{o} 221, 28 p.; t. 14, 1961/62,
n\textsuperscript{o} 232, 19 p.

\item[{[4]}] GROTHENDIECK (Alexander). -- Séminaire de géométrie algébrique,
1960/61. -- Paris, Institut des Hautes Études Scientifiques, 1961.

\item[{[5]}] GROTHENDIECK (Alexander). -- Séminaire de géométrie algébrique,
1961/62, redigido por um grupo de ouvintes. -- Paris, Institut des Hautes Études
Scientifiques (a publicar).

\item[{[6]}] MATSUSAKA (T.). -- The criteria for algebraic equivalence and the
torsion group, \emph{Amer. J. Math.}, t. 79, 1957, p. 53--66.

\item[{[7]}] MUMFORD (D.) e TATE (J.). -- Séminaire de géométrie algébrique,
Harvard University, semestre da primavera de 1962 (a publicar).

\item[{[8]}] SERRE (Jean-Pierre). -- Sur la topologie des variétés algébriques
en caractéristique $p$, \emph{Symposium internacional de topología algebraica}
[1956, México], p. 24--53. -- México, Universidad Nacional Autónoma, 1958.
\end{description}

\par

\def\FgaUnitKey{e236}\clearpage
\chapter*{Errata e complemento à exposição n\textsuperscript{o} 236}
\addcontentsline{toc}{chapter}{Errata e complemento à exposição
n\textsuperscript{o} 236}
\markboth{Errata da exposição 236}{Errata da exposição 236}
\sid{e236}
\src{303}

\noindent\textbf{8. TDTE VI: Os esquemas de Picard. Propriedades gerais.}

\begin{center}
[Seminário Bourbaki, t. 14, 1961/62, n\textsuperscript{o} 236, 23 p.]
\end{center}

\noindent\textbf{Página 12, linha 16.} --- Em vez de «reta afim», leia-se
«reta afim privada da origem».

\noindent\textbf{Linha 18.} --- Em vez de «$X[t]$», leia-se
«$X[t,t^{-1}]$».

\noindent\textbf{Página 16, proposição 3.1, 6\textsuperscript{a} linha.} ---
Em vez de «Ora, a hipótese ...», leia-se «Ora, da hipótese ...».

\noindent\textbf{Página 21, § 4: O teorema de finitude para o esquema de
Picard.} --- As questões de finitude do gênero das levantadas neste parágrafo
foram resolvidas quase totalmente desde a redação da presente exposição.
Indiquemos os principais fatos agora conhecidos nessa direção. Para simplificar
os enunciados, supomos implicitamente que todos os pré-esquemas de Picard que
neles intervêm existem, embora uma modificação evidente desses enunciados
permita eliminar toda hipótese de existência. Doravante, $S$ designa um esquema
noetheriano, e $X,Y$, esquemas próprios sobre $S$.

\medskip
\noindent(i) Seja $f:X\longrightarrow Y$ um morfismo sobrejetivo; então
\[
f^*: \operatorname{Pic}_{Y/S}\longrightarrow\operatorname{Pic}_{X/S}
\]
é um morfismo de tipo finito.
\src{304}
\noindent(ii) Suponhamos $Y$ projetivo sobre $S$, munido de um Módulo
invertível amplo relativamente a $S$, e seja $X$ o esquema dos zeros de uma
seção qualquer desse Módulo, $f:X\longrightarrow Y$ a imersão canônica.
Suponhamos, enfim, que as componentes irredutíveis das fibras de $Y/S$ tenham
dimensão $\geq3$. Então
\[
f^*: \operatorname{Pic}_{Y/S}\longrightarrow\operatorname{Pic}_{X/S}
\]
é de tipo finito.

\noindent(iii) Suponhamos que $X$ seja projetivo sobre $S$ e que todas as suas
fibras geométricas sejam íntegras e de dimensão $n$. Seja
$\mathcal O_X(1)$ um Módulo invertível amplo sobre $X$, que permita definir
polinômios de Hilbert. Seja $M$ um subconjunto de
$\operatorname{Pic}_{X/S}$; então $M$ é quase-compacto se, e somente se, nos
polinômios de Hilbert
\[
a_0x^n+a_1x^{n-1}+a_2x^{n-2}+\cdots
\]
dos elementos de $M$, os coeficientes $a_1$ e $a_2$ permanecem limitados, omitindo-se todo coeficiente cujo monômio teria grau negativo; assim, apenas $a_1$ intervém se $n=1$, e nenhum dos dois se $n=0$.

\noindent(iv) Para todo inteiro $n\neq0$, o homomorfismo de elevação à
$n$-ésima potência no pré-esquema em grupos $\operatorname{Pic}_{X/S}$ é um
morfismo de tipo finito.

Notemos que (i) e (ii) significam também que (sob as hipóteses consideradas
nesses enunciados (i) e (ii)) um subconjunto $M$ de
$\operatorname{Pic}_{Y/S}$ é quase-compacto se, e somente se, sua imagem em
$\operatorname{Pic}_{X/S}$ o é. Conclui-se daí que um Módulo invertível
$\mathcal L$ sobre $Y$ é $\tau$-equivalente a zero se, e somente se, sua imagem
inversa sobre $X$ o é; em outros termos, $\operatorname{Pic}_{Y/S}^\tau$ é a
imagem inversa de $\operatorname{Pic}_{X/S}^\tau$. Em particular, para mostrar
que o primeiro pré-esquema é de tipo finito sobre $S$, basta prová-lo para o
segundo, pois $\operatorname{Pic}_{Y/S}\longrightarrow
\operatorname{Pic}_{X/S}$ é de tipo finito. Utilizando então (i), o lema de
Chow e (iii), obtemos:

\medskip
\noindent(v) $\operatorname{Pic}_{X/S}^\tau$ é de tipo finito sobre $S$.

De modo geral, a conjunção de (i) para um morfismo finito e de (ii) permite
reduzir, para a maioria das questões de finitude, ao caso em que $X/S$ tem
fibras geométricas íntegras normais de dimensão $\leq2$; muitas vezes, aplicando
(i) a um morfismo sobrejetivo não necessariamente finito e a resolução das
singularidades das superfícies algébricas (devida, em característica qualquer,
a Abhyankar), reduzimos ao caso em que $X/S$ é até mesmo liso, portanto tem
fibras geométricas não singulares de dimensão 2. Isso permite, por exemplo,
tendo em conta (v) e a desigualdade de Picard--Igusa que majora o posto do grupo
de Néron--Severi de uma superfície projetiva não singular, provar a seguinte
generalização do teorema de finitude de Néron:

\medskip
\noindent(vi) Seja $X/S$ próprio sobre $S$ e, de resto, arbitrário; então os
grupos de Néron--Severi
\[
\operatorname{Pic}_{X_i/k_i}/\operatorname{Pic}_{X_i/k_i}^\circ
\]
das fibras geométricas $X_i/k_i$ de $X/S$
\src{305}
são de tipo finito, e seu posto e a ordem de seu subgrupo de torção permanecem
limitados.

O mesmo método de redução ao caso das superfícies não singulares e os teoremas
conhecidos nesse caso (a saber, o teorema de finitude de Néron e o fato de que a
forma de interseção sobre o grupo de Néron--Severi é não degenerada) implicam:

\medskip
\noindent(vii) Seja $X$ um esquema próprio sobre um corpo algebricamente
fechado. Existe um número finito de curvas fechadas íntegras $C_i$
($1\leq i\leq r$) em $X$, com a seguinte propriedade: para todo subconjunto
$M$ de $\operatorname{Pic}_{X/k}$, $M$ é quase-compacto se, e somente se, os
inteiros
\[
\deg_{C_i'}\mathcal L\qquad(\mathcal L\in M)
\]
permanecem limitados (onde $C_i'$ designa a normalização de $C_i$).

Pode-se tomar aqui, para $r$, o posto do grupo de Néron--Severi. Uma vez sabido
que este último é de tipo finito, (vii) reduz-se ao fato de que as formas
lineares sobre o grupo de Néron--Severi definidas por curvas $C$ em $X$ só se
anulam simultaneamente sobre os elementos de torção do grupo de Néron--Severi.
No caso em que $X$ é projetivo e não singular, esse resultado, assim como (v),
era devido a Matsusaka. Utilizando (vii), obtém-se facilmente a seguinte
caracterização dos Módulos invertíveis $\tau$-equivalentes a zero sobre $X$:

\medskip
\noindent(viii) Sejam $X/k$ um esquema próprio sobre um corpo e $\mathcal L$ um
Módulo invertível sobre $X$. Condições equivalentes:

\noindent a. $\mathcal L$ é $\tau$-equivalente a zero;

\noindent b. Para todo Módulo coerente $F$, tem-se
\[
\chi(F\otimes\mathcal L)=\chi(F),
\]
onde $\chi$ designa a característica de Euler--Poincaré;

\noindent b'. Como (b), mas com $F=\mathcal O_Y$, sendo $Y$ um subesquema
fechado íntegro de dimensão 1 em $X$;

\noindent c. Para todo $Y$ como acima, designando por $Y'$ a curva normalizada,
tem-se
\[
\deg\mathcal L_{Y'}=0.
\]

Se $X/k$ é projetivo e munido de um Módulo invertível amplo
$\mathcal O_X(1)$, as condições precedentes são também equivalentes às
seguintes:

\noindent d. Para todo inteiro $m$, $\mathcal L^{\otimes m}(1)$ é amplo.

\noindent e. (Se $X$ é íntegro.) Para todo par de inteiros $m,n$, tem-se
\[
\chi(\mathcal L^{\otimes m}(n))=\chi(\mathcal O_X(n))
\]
\src{306}
(isto é, (b) é verdadeiro tomando $F=\mathcal L^{\otimes m}(n)$).

Para a suficiência desta última condição, note-se que ela significa que os
polinômios de Hilbert dos $\mathcal L^{\otimes m}$ são todos iguais; portanto,
em virtude do critério de Mumford (iii), os $\mathcal L^{\otimes m}$ permanecem
em um subconjunto quase-compacto de $\operatorname{Pic}_{X/k}$, isto é, vale
(a). As condições (b), (b') e (c) devem ser consideradas variantes (sobre um
esquema próprio qualquer) da noção de equivalência numérica, definida
habitualmente sobre variedades projetivas não singulares. Para estas últimas,
a equivalência de (a) e (c) era evidentemente bem conhecida (Matsusaka).

O critério (e) implica também o resultado seguinte:

\medskip
\noindent(ix) Seja $f:X\longrightarrow S$ um morfismo projetivo e plano cujas
fibras geométricas são íntegras. Então $\operatorname{Pic}_{X/S}^\tau$ é aberto
e fechado em $\operatorname{Pic}_{X/S}$.

Limitar-nos-emos a alguns comentários sobre a demonstração dos resultados-chave
(i), (ii), (iii) (o resultado (iv) é um pouco apartado dos demais e se demonstra
utilizando somente (i) para os morfismos finitos sobrejetivos radiciais, mais
precisamente, para um morfismo de Frobenius). Para (i), utilizam-se de modo
essencial as ideias da descida não plana (ver, em particular, TDTE I, página 9).
Ocorre que (tendo em vista apenas resultados de finitude) a falta de critérios
de efetividade para dados de descida é inofensiva. Por outro lado, Mumford
demonstrou recentemente uma forma um pouco menos forte de (iii), a saber, o
critério que envolve todos os coeficientes do polinômio de Hilbert. Seu
argumento é extremamente simples e inspirado pela demonstração de um critério
de amplitude de Nakai (enunciado por este último para as superfícies não
singulares e generalizado por Mumford aos morfismos projetivos quaisquer).
Parece-me, aliás, que esse argumento só é válido mediante uma ligeira restrição
adicional sobre as fibras de $X/S$, a propriedade $(S_2)$ de Serre, verificada
se as fibras geométricas são normais. Utiliza-se então esse critério restrito na
demonstração de (ii): o critério (i) permite, com efeito, reduzir ao caso em que
$Y/S$ é plano com fibras geométricas íntegras e normais e, aplicando o critério
de Mumford, reduz-se facilmente ao caso em que $X/S$ satisfaz as mesmas
condições. Da hipótese de dimensão resulta então que as fibras geométricas de
$Y$ e de $X$ têm profundidade $\geq2$ em seus pontos fechados, o que permite
aplicar os «critérios de equivalência» sob a forma que lhes é dada em SGA 1962
(exposição XII) e concluir a demonstração de (ii). Uma vez que dispomos de (i) e
(ii), não é difícil eliminar do critério de Mumford toda hipótese de normalidade
sobre as fibras e demonstrá-lo também sob a forma mais forte dada em (iii).
\src{307}
Notemos, enfim, que a demonstração de (i) e (ii) mostra igualmente que, no
caso em que $S$ é o espectro de um corpo $k$, o morfismo
\[
\operatorname{Pic}_{Y/k}\longrightarrow\operatorname{Pic}_{X/k}
\]
é afim (e não apenas de tipo finito).

\par

\def\FgaUnitKey{com}\unit{Comentários de 1962}
\sid{com}
\markboth{Comentários de 1962}{Comentários de 1962}
\src{297}
\seminares{14.º ano, 1961/62, Suplemento\\Maio de 1962}
\paperhead{FUNDAMENTOS DA GEOMETRIA ALGÉBRICA.\\COMENTÁRIOS}%
  {por Alexander GROTHENDIECK}

\section*{Advertência}
\addcontentsline{toc}{section}{Advertência}

Apresentamos, no Seminário Bourbaki, entre 1957 e 1962, oito exposições
relativas aos Fundamentos da geometria algébrica. Com exceção da primeira,
essas exposições são redigidas no âmbito dos esquemas. Todos os resultados
enunciados encontrarão seu lugar nos «Éléments de géométrie algébrique» de Jean
DIEUDONNÉ e Alexander GROTHENDIECK. Contudo, o conteúdo de nenhuma dessas
exposições está atualmente coberto nem por um dos capítulos já publicados ou
em preparação, nem por qualquer outro livro ou artigo, e sem dúvida não o
estará ainda por alguns anos. Essa é a principal razão que nos levou a reunir
essas exposições, que colocarão à disposição dos usuários certo número de
noções e resultados-chave da teoria dos esquemas, enquanto se aguarda uma
redação sistemática. Além disso, a leitura desses textos permitirá ao leitor
familiarizar-se rapidamente com os referidos resultados e noções, sem que os
pormenores, necessariamente volumosos, de um tratado sistemático constituam
um obstáculo, e fornecerá as motivações indispensáveis para o estudo de tal
tratado.

Para comodidade do leitor, reunimos aqui certo número de comentários e
erratas, agrupados por exposições, que assinalarão em particular, de passagem,
os progressos realizados desde a redação do texto e indicarão algumas
referências adicionais.

Vários resultados apresentados nessas exposições foram tratados em detalhe no
Seminário de Geometria Algébrica do IHES, bem como em dois seminários
consecutivos realizados em Harvard em 1961/62, conduzidos o primeiro por mim e
o segundo por MUMFORD--TATE, para os quais estão atualmente em preparação
notas mimeografadas (redigidas pelo Sr. LICHTENBAUM).

\section*{Siglas utilizadas}
\addcontentsline{toc}{section}{Siglas utilizadas}

\begin{description}
\item[SGA :] Seminário de Geometria Algébrica, Institut des Hautes Études
Scientifiques, 1960--1962.

\src{298}
\item[SHG :] Seminário GROTHENDIECK, Harvard University, 1961/62.

\item[SHMT :] Seminário MUMFORD--TATE, Harvard University, 1961/62.

\item[EGA :] GROTHENDIECK (A.) e DIEUDONNÉ (J.). --- Éléments de géométrie
algébrique. --- Paris, Presses universitaires de France, 1960, 1961, …
(Institut des Hautes Études Scientifiques, Publications mathématiques, 4, 8,
11, …).

\item[TDTE :] GROTHENDIECK (Alexander). --- Technique de descente et théorèmes
d'existence en géométrie algébrique, I--VI, \emph{Séminaire Bourbaki}, t. 12,
1959/60, n.os 190 e 195; t. 13, 1960/61, n.os 212 e 221; t. 14, 1961/62,
n.os 232 e 236.
\end{description}

\par

\end{document}
